% !TeX program = pdflatex
% Final filename: 1929 von Neumann_On the Algebra of Functional Operations and the Theory of Normal Operators_GPT.tex
\documentclass[
  paper=b5,
  twoside=true,
  fontsize=13pt,
  headings=normal,
  bookmarkpackage=false
]{scrartcl}

\usepackage[T1]{fontenc}
\usepackage[utf8]{inputenc}
\usepackage{amsmath,mathtools}
\usepackage{newpxtext,newpxmath}
\usepackage[
  b5paper,
  left=0.7cm,
  right=0.7cm,
  top=0.91cm,
  bottom=0.91cm,
  % Count the page-number header in the vertical layout: 0.91 cm lies above it.
  includehead,
  headheight=0.55cm,
  headsep=0.12cm
]{geometry}
\usepackage{microtype}
\usepackage{xparse}
\usepackage{xcolor}
\usepackage[singlespacing=true]{scrlayer-scrpage}

\clearpairofpagestyles
\ohead*{\pagemark}
\setkomafont{pageheadfoot}{\normalfont}
\pagestyle{scrheadings}
\raggedbottom

\definecolor{InternalLink}{HTML}{1E5A91}
\definecolor{CitationLink}{HTML}{773F91}
\definecolor{ExternalLink}{HTML}{087F76}
\usepackage[
  unicode,
  colorlinks=true,
  linkcolor=InternalLink,
  citecolor=CitationLink,
  urlcolor=ExternalLink,
  pdfborder={0 0 0},
  pdfpagelayout=OneColumn
]{hyperref}
\hypersetup{
  pdftitle={Zur Algebra der Funktionaloperationen und Theorie der normalen Operatoren / On the Algebra of Functional Operations and the Theory of Normal Operators},
  pdfauthor={John von Neumann},
  pdfsubject={English translation of Mathematische Annalen 102 (1929), 370--427},
  pdfkeywords={von Neumann algebra, functional operations, normal operators, spectral theorem}
}

\setkomafont{disposition}{\rmfamily\bfseries}
\setkomafont{section}{\rmfamily\bfseries\large}
\setkomafont{subsection}{\rmfamily\bfseries\normalsize}
\renewcommand{\thesection}{\Roman{section}}
\RedeclareSectionCommand[tocnumwidth=2.8em]{section}
\setcounter{secnumdepth}{1}
\setcounter{tocdepth}{1}

\linespread{1.1}
\AtBeginDocument{\fontsize{13pt}{15.6pt}\selectfont}
\setlength{\parindent}{1.2em}
\setlength{\parskip}{0pt}
\setlength{\emergencystretch}{3em}
\interfootnotelinepenalty=0
\allowdisplaybreaks[2]
\urlstyle{same}

\newcounter{statement}
\NewDocumentEnvironment{theorem}{o}
  {\par\medskip\refstepcounter{statement}\noindent
   \textbf{\IfNoValueTF{#1}{Theorem}{#1}.}\ \ignorespaces}
  {\par\medskip}
\NewDocumentEnvironment{definition}{o}
  {\par\medskip\refstepcounter{statement}\noindent
   \textbf{\IfNoValueTF{#1}{Definition}{#1}.}\ \ignorespaces}
  {\par\medskip}
\NewDocumentEnvironment{proof}{o}
  {\par\smallskip\noindent
   \textit{\IfNoValueTF{#1}{Proof}{#1}.}\ \ignorespaces}
  {\hfill\(\square\)\par\medskip}

\newcommand{\NoteRef}[1]{{\hypersetup{linkcolor=CitationLink}%
  \hyperref[fn:#1]{note~#1}}}
\newcommand{\TheoremRef}[1]{\hyperref[thm:#1]{Theorem~#1}}
\newcommand{\DefinitionRef}[1]{\hyperref[def:#1]{Definition~#1}}
\newcommand{\SectionRef}[2]{\hyperref[#1]{#2}}
\newcommand{\Eref}{{\hypersetup{linkcolor=CitationLink}%
  \hyperref[fn:E-paper]{\emph{E.}}}}
\newcounter{translatornote}
\newcounter{savedfootnote}
\newcommand{\TranslatorNote}[1]{%
  \stepcounter{translatornote}%
  \setcounter{savedfootnote}{\value{footnote}}%
  \begingroup
    \renewcommand{\thefootnote}{T\arabic{translatornote}}%
    \footnote{\textit{Translator's note.} #1}%
  \endgroup
  \setcounter{footnote}{\value{savedfootnote}}}
\newcommand{\TranslatorNoteMark}{%
  \stepcounter{translatornote}%
  \setcounter{savedfootnote}{\value{footnote}}%
  \begingroup
    \renewcommand{\thefootnote}{T\arabic{translatornote}}%
    \footnotemark
  \endgroup
  \setcounter{footnote}{\value{savedfootnote}}}
\newcommand{\TranslatorNoteText}[1]{%
  \setcounter{savedfootnote}{\value{footnote}}%
  \begingroup
    \renewcommand{\thefootnote}{T\arabic{translatornote}}%
    \footnotetext{\textit{Translator's note.} #1}%
  \endgroup
  \setcounter{footnote}{\value{savedfootnote}}}

\begin{document}
\begin{titlepage}
  \setcounter{page}{1}
  \centering
  \vspace*{0.10\textheight}
  {\Large\bfseries
   Zur Algebra der Funktionaloperationen\\[0.35em]
   und Theorie der normalen Operatoren\par}
  \vspace{1.4em}
  {\large\bfseries
   On the Algebra of Functional Operations\\[0.35em]
   and the Theory of Normal Operators\par}
  \vspace{2.2em}
  {\large J. v. Neumann (Berlin)\par}
  \vfill
  {\normalsize
   English translation of\\
   \textit{Mathematische Annalen} \textbf{102} (1929), 370--427\\[0.5em]
   \href{https://doi.org/10.1007/BF01782352}{doi:10.1007/BF01782352}\par}
  \vspace*{0.10\textheight}
  \thispagestyle{empty}
\end{titlepage}
\setcounter{page}{2}

\pdfbookmark[1]{Contents}{contents}
\tableofcontents
\clearpage

% [Source p. 370 / PDF 1]
\section*{Introduction}
\phantomsection\label{sec:introduction}
\addcontentsline{toc}{section}{Introduction}

\phantomsection
\noindent\textbf{1.}\label{intro:1}
The present paper falls into two essentially independent parts. The first
(\S\S~I--III) is devoted to the investigation of the linear and bounded
operators (i.e.\ matrices) of Hilbert space $\mathfrak H$, by considering the
algebraic properties of the (noncommutative) ring $\mathfrak B$ formed by
them. The subject of the second part, on the other hand, is formed by those
operators, not necessarily everywhere defined (in $\mathfrak H$) and
bounded, which admit the so-called Hilbert spectral representation with
complex eigenvalues (cf.\ the more detailed explication of these notions in
\hyperref[intro:4]{paragraph 4 of the Introduction}). These are the operators
to be designated as ``normal,'' which hitherto have been considered only in
the bounded case,\footnote{\label{fn:1}Until recently, only in the completely continuous
case; cf.\ the encyclopedia article by Hellinger and Toeplitz,
\emph{Encyklop\"adie der mathematischen Wissenschaften}, II C 13, p.~1562.
See also note~\ref{fn:normal-spectral-form}.} and for which we shall give a
new, more general definition (cf.\ the passage just cited).

Before we discuss these matters more precisely, let us recall the definition
of (complex) Hilbert space $\mathfrak H$. It may be realized as the set of
all sequences of complex numbers
\[
  \{x_1,x_2,\ldots\}
  \qquad\text{with}\qquad
  \sum_{n=1}^{\infty}|x_n|^2<\infty;
\]
we denote its elements by $f,g,\ldots,\varphi,\psi,\ldots$.\footnote{\label{fn:2}By the
well-known theorem of Fischer and F.~Riesz it may equally well be realized
as the set of all functions $f(x)$ on an interval $a<x<b$ for which
\[
  \int_a^b |f(x)|^2\,dx<\infty
\]
($a<b$, and either endpoint may also be infinite); or as the set of all
functions $f(P)$, where $P$ ranges over the surface of the unit sphere, for
which
\[
  \iint |f(P)|^2\,d\sigma<\infty
\]
($\iint\cdots d\sigma$ denoting integration over the surface of the unit
sphere); and so forth. Cf., for example, the paper cited in \NoteRef{7},
especially Chapter~I, Appendix~I, and the Introduction; for the theorem of
Fischer and F.~Riesz, see, for instance, F.~Riesz, \emph{G\"ottinger
Nachrichten} (1907), pp.~210--273.}

% [Source p. 371 / PDF 2]
One may calculate with these elements as with vectors, so that the meaning
of constructions such as $f\pm g$ and $af$ is immediately
clear.\footnote{\label{fn:3}When $\mathfrak H$ is interpreted as the ``sequence space''
of the $\{x_1,x_2,\ldots\}$ for which $\sum_{n=1}^{\infty}|x_n|^2$ is
finite, one defines
\[
  \{x_1,x_2,\ldots\}\pm\{y_1,y_2,\ldots\}
  =\{x_1\pm y_1,x_2\pm y_2,\ldots\},
\]
\[
  a\{x_1,x_2,\ldots\}=\{ax_1,ax_2,\ldots\}.
\]
In the ``function space'' of the $f(x)$ for which
$\int_a^b|f(x)|^2\,dx$ is finite, the corresponding definitions are
$(f\pm g)(x)=f(x)\pm g(x)$ and $(af)(x)=af(x)$; analogously for the other
examples. Cf.\ the reference in \NoteRef{2}.} There is, moreover, an ``inner
product'' $(f,g)$ as for vectors,\footnote{\label{fn:4}One defines
\[
 \bigl(\{x_1,x_2,\ldots\},\{y_1,y_2,\ldots\}\bigr)
   =\sum_{n=1}^{\infty}x_n\overline{y_n}
\]
(the series converges absolutely), or, respectively,
\[
  (f(x),g(x))=\int_a^b f(x)\overline{g(x)}\,dx,
\]
and so forth. Cf.\ the preceding notes.} which permits the definition of
the ``absolute value'' in Hilbert space,
\[
  |f|=\sqrt{(f,f)},
\]
and of the distance $|f-g|$.\footnote{\label{fn:5}Thus
\[
  |\{x_1,x_2,\ldots\}|
     =\sqrt{\sum_{n=1}^{\infty}|x_n|^2},
  \qquad
  |f(x)|=\sqrt{\int_a^b |f(x)|^2\,dx}.
\]
Unfortunately the notation is ambiguous here: on the left of the second
formula stands the absolute value of the function $f(x)$, while on the
right stand the absolute values of its values. Such an ambiguity will never
arise in our considerations in $\mathfrak H$. The other realizations are
treated analogously; cf.\ the preceding notes.} In this way Hilbert space
$\mathfrak H$ becomes a topological (indeed, metric) linear
space,\footnote{\label{fn:6}Cf.\ Hausdorff, ``Umgebungsaxiome und
topologisch-lineare R\"aume,'' in \emph{Grundz\"uge der Mengenlehre},
1st ed., Leipzig and Berlin, 1914.} in which geometrical-topological
expressions such as linear manifold, accumulation point, continuity, and so
forth may be used without further explanation.

For a systematic characterization of $\mathfrak H$ on this basis, one may
compare an earlier paper by the author,\footnote{\label{fn:7}\label{fn:E-paper}
``Zur allgemeinen Eigenwerttheorie Hermitescher Funktionaloperatoren,''
\emph{Mathematische Annalen} \textbf{102}, no.~1 (1929), pp.~49--131.
This paper, which will be cited repeatedly below, will always be denoted by
E. Several abbreviations from E.\ that we shall use are important; they are
collected here: closed = cl., bounded = bd., extension = ext., Hermitian =
H., hypermaximal = hypermax., linear = lin., linear manifold = lin.~mfd.,
maximal = max., normalized = norm., operator = op., orthogonal = orth.,
projection = P., complete = compl., resolution of the identity = res.~id.
We shall further abbreviate the term ``commutable'' (cf.\ \NoteRef{16}) by
comm.} whose notation and concepts will also be used here. In this
connection compare especially the portions named in \NoteRef{2}; nevertheless,
we shall always cite precisely the sections of E.\ that are used.

The definitions from E.\ most important for our purposes are the following
(abbreviations as in \NoteRef{7}). An operator is a function whose range and
domain are subsets of $\mathfrak H$; we denote operators by
$A,B,\ldots,R,S,\ldots$, and the value of $A$ at $f$ by $Af$. (Thus $Af$
need by no means be defined everywhere in $\mathfrak H$.) It is linear if,
whenever $Af$ and $Ag$ are defined, $A(af)$ and $A(f+g)$ are also defined
(i.e.\ its domain

% [Source p. 372 / PDF 3]
is a linear manifold),\footnote{\label{fn:8}That is, together with $f,g$ it contains
$af,f+g$. A closed linear manifold must in addition be closed, i.e.\ contain
all its accumulation points.} and their values are respectively $aAf$ and
$Af+Ag$. It is called closed if it has the following property: whenever all
$Af_n$ ($n=1,2,\ldots$) are defined and
\[
  f_n\longrightarrow f,\qquad Af_n\longrightarrow f^*
  \qquad(n\longrightarrow\infty),
\]
then $Af$ is defined and equals $f^*$.\footnote{\label{fn:9}In compact spaces this
continuity-like property is equivalent to continuity, but in $\mathfrak H$
it is much weaker; cf.\ \Eref, note~29.}

An operator is continuous if it is so, in the usual sense, throughout its
domain: if $Af_0$ is defined, then for every $\epsilon>0$ there is a
$\delta>0$ such that
\[
  |f-f_0|<\delta \quad\Longrightarrow\quad |Af-Af_0|<\epsilon
\]
whenever $Af$ is defined. One sees at once that the domain of a closed,
continuous operator is necessarily a closed set; if the operator is also
linear, its domain is therefore a closed linear
manifold.\footnote{\label{fn:10}For linear operators continuity can be reduced to
several equivalent formulations; cf.\ \Eref, Theorem~12. One of these is:
there is a fixed $c$ such that $|f|\leq 1$ implies $|Af|\leq c$.}

In the first half of this paper, as already announced, we shall concern
ourselves with everywhere-defined, linear, continuous operators, which (in
accordance with Hilbert) will be called bounded. (Cf., also for what follows,
Appendix~I of \Eref) For these everywhere-defined operators the simplest
arithmetic operations are defined as follows. If $A,B$ are bounded, then so
are $aA,A\pm B,AB$, and
\[
  (aA)f=a\cdot Af,\qquad
  (A\pm B)f=Af\pm Bf,\qquad
  (AB)f=A(Bf).
\]
The rules of algebra hold, with two exceptions: there are zero divisors,
and multiplication is not commutative.\footnote{\label{fn:11}$0$ and $1$ are defined by
$0f=0$ and $1f=f$.} Thus the set $\mathfrak B$ of all bounded operators is
a ring---with zero divisors and noncommutative. Another important operation
is $*$, defined as follows: for every bounded $A$ there is exactly one
bounded $A^*$ such that
\[
  (Af,g)=(f,A^*g),\qquad (f,Ag)=(A^*f,g).
\]\footnote{\label{fn:12}Either relation follows from the other by interchanging $f,g$
and taking complex conjugates. The operation $*$ corresponds, for
matrices, to taking the complex-conjugate transpose. More precisely, if
$\mathfrak H$ is realized as the ``sequence space'' of the
$\{x_1,x_2,\ldots\}$ for which $\sum_{n=1}^{\infty}|x_n|^2$ is finite, and
$A$ has the matrix $\{a_{\mu\nu}\}$, so that
\[
  A\{x_1,x_2,\ldots\}=\{y_1,y_2,\ldots\},
  \qquad
  y_\nu=\sum_{\mu=1}^{\infty}a_{\mu\nu}x_\mu
\]
always holds---as in Appendix~II of \Eref{}---then $A^*$ has the matrix
$\{\overline{a_{\nu\mu}}\}$. Cf.\ \Eref, Appendix~II.}

% [Source p. 373 / PDF 4]
The following rules are easily verified:
\[
  0^*=0,\qquad 1^*=1,\qquad
  (aA)^*=\overline a\,A^*,\qquad
  (A\pm B)^*=A^*\pm B^*,\qquad
  (AB)^*=B^*A^*.
\]

For the unbounded (more precisely: not necessarily bounded) operators that
form the subject of the second half of this paper, we must proceed somewhat
differently. For such operators $R,S$, in defining $aR,R\pm S,RS$, one
must first observe that the domains of these operators are thereby
restricted, since $Rf,Sf$ need not be defined everywhere. In order that
$(R+S)f$ be defined, both $Rf$ and $Sf$ must be defined, and likewise for
$(R-S)f$; for $RSf$, both $Sf$ and $R(Sf)$ must be defined. In the case of
$R^*$ it is even uncertain whether such an operator exists. Yet the
operation $*$ is of particular importance for us. We shall therefore, when
necessary, begin directly with an already specified pair $R,R^*$, stipulating the
following.

\medskip
\noindent\textbf{Conjugate operator pair.}
A conjugate operator pair (for short, conj.\ op.\ pair) consists of two
operators $R,R^*$ having the same domain and satisfying, whenever both
sides are defined,
\[
  (Rf,g)=(f,R^*g),\qquad (f,Rg)=(R^*f,g)
\]
(cf.\ \NoteRef{12}). In addition, in analogy with the H.\ operators and for the
same reasons (cf.\ \Eref, Introduction~V, Definition~6, and note~30), we
require that the domain span the closed linear manifold $\mathfrak H$,
i.e.\ that it be dense in $\mathfrak H$. This definition will prove to be
the correct generalization of the operation $*$ in $\mathfrak B$; H.\
operators are included as the special case $R=R^*$.

\medskip
\phantomsection
\noindent\textbf{2.}\label{intro:2}
We first make a few remarks about the points of view governing the
investigation of $\mathfrak B$. We regard this ring as a hypercomplex number
system. Indeed, it is probably the simplest system of this kind with
infinitely many units---that is, the simplest one that goes beyond the
systems usually investigated and characterized essentially by the validity
of the so-called ``chain conditions.''\footnote{\label{fn:13}Cf.\ E.~Noether,
\emph{Mathematische Annalen} \textbf{96}, no.~1 (1926), pp.~26--61,
especially \S~2; also E.~Artin, \emph{Abhandlungen aus dem Mathematischen
Seminar der Universit\"at Hamburg} \textbf{5}, no.~3 (1927), pp.~251--260.}
Whereas in hypercomplex systems with finitely many units (or with the
``chain conditions'' that replace this assumption) the topology is so
trivial that it plays no noteworthy role, in our infinite system
$\mathfrak B$ topology plays a decisive role. Thus, for example, in the
definition of which subsets $\mathfrak M$ of $\mathfrak B$ are also rings
(in the terminology of the works cited in \NoteRef{13}: ``orders''), if we
required only that $\mathfrak M$ contain $aA,A\pm B,AB$ whenever it contains
$A,B$, we should lose ourselves in an inextricable thicket of pathological
constructions. One must also require that $\mathfrak M$ be closed in a
(still to be formulated) topology of $\mathfrak B$.

% [Source p. 374 / PDF 5]
We shall, moreover, impose one further simplifying requirement on the rings
$\mathfrak M$: together with $A$ they must also contain $A^*$. This is a
significant restriction, which makes it possible to avoid the complications
of elementary-divisor theory.\footnote{\label{fn:14}For matrices in finite-dimensional
spaces, E.~Fischer was the first to introduce this condition successfully.}
There is good reason to exclude these difficulties at first; compared with
the finite-dimensional case,\footnote{\label{fn:15}Every hypercomplex system with
finitely many units can, of course, be represented by matrices (and hence
operators) of a finite-dimensional space.} the theory of $\mathfrak B$
already furnishes enough new possibilities for complication.

It follows from what has been said that we must first concern ourselves
with the topology of $\mathfrak B$. A difficulty is that there are several
ways of topologizing $\mathfrak B$ (i.e.\ of defining its notion of
neighborhood) that are relevant for us. We shall discuss these precisely in
Section~\ref{sec:topologies}; here it is enough to say that rings
$\mathfrak M$ in $\mathfrak B$ will be required to be closed in the
so-called weak topology.

Since the intersection of arbitrarily many rings is again a ring, the
intersection of all rings containing $\mathfrak M$ exists in particular
($\mathfrak M$ being any subset of $\mathfrak B$). It is plainly
characterized uniquely as the smallest ring containing $\mathfrak M$; we
call it $R(\mathfrak M)$. Another important construction is the following.
Let $\mathfrak M$ be a subset of $\mathfrak B$. The set of all
$A\in\mathfrak B$ for which both $A$ and $A^*$ commute with every
$B\in\mathfrak M$\footnote{\label{fn:16}Two elements $A,B$ of $\mathfrak B$ are
commutable if $AB=BA$.} will be denoted by $\mathfrak M'$. The operation
$'$ can be iterated, thereby producing successively the sets
\[
  \mathfrak M',\ \mathfrak M'',\ \mathfrak M''',\ldots
\]
from $\mathfrak M$. It has a number of simple properties; among other things
we shall show that
\[
  \mathfrak M\subset\mathfrak M'',\qquad
  \mathfrak M\subset\mathfrak N\Longrightarrow
  \mathfrak M'\supset\mathfrak N',
\]
and, in general,
\[
  \mathfrak M'=\mathfrak M'''=\mathfrak M^{\mathrm V}=\cdots,\qquad
  \mathfrak M''=\mathfrak M^{\mathrm{IV}}=\mathfrak M^{\mathrm{VI}}=\cdots.
\]\footnote{\label{fn:17}$\mathfrak M\subset\mathfrak N$ or
$\mathfrak N\supset\mathfrak M$ means that $\mathfrak M$ is a subset of
$\mathfrak N$.}

Together, the two operations $R(\mathfrak M)$ and $\mathfrak M'$ permit a
simple characterization of the algebraic relations in $\mathfrak B$.

\medskip
\phantomsection
\noindent\textbf{3.}\label{intro:3}
An operator $A$ is called normal if $A$ and $A^*$ commute (cf.\ \Eref,
Appendix~I). A ring $\mathfrak M$ is called Abelian if all its elements
commute with one another. One readily shows that a ring $\mathfrak M$ is
Abelian if and only if all its elements are normal, and that the ring
$R(A)$ is Abelian if and only if $A$ is normal. We shall now prove the
converse: every Abelian ring $\mathfrak M$ can be generated by a normal
operator (indeed, by an H.\ operator); that is, there is such an $A$ with
\[
  \mathfrak M=R(A).
\]

% [Source p. 375 / PDF 6]
There is a further simplification in the case of Abelian rings. Let
$\mathfrak M$ be a subset of $\mathfrak B$, and let $R(\mathfrak M)$ be the
ring generated by $\mathfrak M$. One easily shows that $R(\mathfrak M)$ is
obtained by first forming, from the elements of $\mathfrak M$, all operators
arising through finitely many applications of the operations
\[
  aA,\qquad A\pm B,\qquad AB,\qquad A^*,
\]
and denoting their set by $r(\mathfrak M)$, and then adjoining to
$r(\mathfrak M)$ all its accumulation points. Here, as mentioned in
\hyperref[intro:2]{paragraph 2}, accumulation point is to be understood in
the sense of the so-called weak topology. If $R(\mathfrak M)$ is Abelian
(this is the case if and only if all elements of $\mathfrak M$ commute with
one another and with their adjoints), we shall show that it suffices to
adjoin to $r(\mathfrak M)$ a much smaller set of accumulation points---the
so-called limits of strongly doubly convergent sequences; cf.\ the
discussion of these notions in Section~\ref{sec:topologies}. This already
produces the whole of $R(\mathfrak M)$. The result is especially important
for the theory of functions of commuting normal (hence, in particular, H.\
and unitary) operators, although that theory lies beyond the scope of the
present paper.

For arbitrary (not necessarily Abelian) rings we shall establish a relation
between $R(\mathfrak M)$ and $\mathfrak M''$, where $\mathfrak M$ is any
subset of $\mathfrak B$. In particular, we shall show (our most general
result is slightly broader; cf.\
\hyperref[sec:ring-and-commutant]{\S\,II}, especially \TheoremRef{5}) that if
$1\in\mathfrak M$, then
\[
  R(\mathfrak M)=\mathfrak M''.
\]
To see the significance of this theorem for $\mathfrak B$, regarded as a
hypercomplex number system, consider what it says when Hilbert space
$\mathfrak H$ is replaced by a finite-dimensional, say $k$-dimensional,
Euclidean space. Let $\mathfrak M$ be, in particular, a finite or continuous
group of unitary matrices. It is then immediately clear that
$R(\mathfrak M)$ is the set of all linear combinations of matrices in
$\mathfrak M$ (among which there can be at most $k^2$ linearly independent
ones).

First suppose that $\mathfrak M$, as a set of matrices, is irreducible. As
is well known, this means that only the matrices $a1$ commute with all
elements of $\mathfrak M$;\footnote{\label{fn:18}Cf.\ I.~Schur, \emph{Berliner Berichte}
(1905), p.~406. The finiteness of $\mathfrak M$ always assumed there is, as
is well known, immaterial in this case.} that is, $\mathfrak M'$ consists
of these alone. Hence $\mathfrak M''$ comprises all matrices whatsoever,
and by our theorem the same must hold for $R(\mathfrak M)$. Thus every
matrix is a linear combination of matrices from $\mathfrak M$, or,
equivalently (since there are $k^2$ linearly independent matrices), there
are $k^2$ linearly independent matrices in $\mathfrak M$. Put another way:
no fixed linear relation among the $k^2$ matrix entries can hold for every
element of $\mathfrak M$. This, however, is Burnside's theorem on
irreducible matrix groups.\footnote{\label{fn:19}Cf.\ Burnside, \emph{Proceedings of the
London Mathematical Society}, ser.~2, \textbf{3} (1905), p.~430---though
there only for unitary matrices.} If we do not assume that $\mathfrak M$

% [Source p. 376 / PDF 7]
is irreducible, one sees just as readily that our theorem corresponds to
the Frobenius--I.~Schur strengthening of Burnside's
theorem.\footnote{\label{fn:20}Cf.\ Frobenius and I.~Schur, \emph{Berliner Berichte}
(1906), p.~209.}

We have therefore proved the analogue of the theorem that, in finitely many
dimensions, can serve as the foundation of the theory of unitary
representations of groups. We shall not investigate here to what extent an
equivalent theory in Hilbert space can be based upon our theorem.

These remarks exhaust the first half of the present paper.

\medskip
\phantomsection
\noindent\textbf{4.}\label{intro:4}
The subject of the second half (\S\S~IV--V) is the conjugate operator pairs
mentioned at the end of \hyperref[intro:1]{paragraph 1}, independently of
any boundedness assumption. The first task is to characterize the normal
operators among them. A difficulty now arises: the pair $R,R^*$ ought to be
called normal when $R$ and $R^*$ commute; but since $R,R^*$ need not be
defined everywhere, the definition of the products $RR^*,R^*R$, and hence
the reduction of commutability to them, is at first
uncertain.\footnote{\label{fn:21}The equality $RS=SR$ can by no means characterize the
commutability of $R,S$. Suppose, for example, that $Rf$ is not defined
everywhere. The same is then true of $0Rf$, whereas $R0f$ is defined
everywhere; hence $0R\neq R0$, i.e.\ $R$ would fail to commute with $0$,
which must not occur under any reasonable notion of commutability.

Quite generally, forming $AB$ or $A\pm B$ for operators $A,B$ that are not
everywhere defined is a questionable matter. There are, for example,
hypermaximal H.\ operators $A,B$ for which $Af$ and $Bf$ are never
simultaneously defined (except for $f=0$), so that $A+B$ cannot be formed at
all. Cf.\ the author's forthcoming paper in \emph{Journal f\"ur die reine
und angewandte Mathematik}, ``Zur Theorie der unbeschr\"ankten Matrizen,''
where very general classes of counterexamples are given.} In
\hyperref[app:commutativity-unbounded]{Appendix~III}
we shall see that the natural matrix definition of commutability likewise
fails.

We proceed as follows. If, of two operators $R,A$, at least $A$ is defined
everywhere, we define their commutability by requiring that $AR$ be an
extension of $RA$. That is, whenever $Rf$ is defined, $R(Af)$ must also be
defined and must equal $A(Rf)$ (where $Af$ and $A(Rf)$ are in any event
defined).\footnote{\label{fn:22}The example $R,0$ in \NoteRef{21} is precisely what suggests
this definition. Moreover, if $E$ is a projection operator, the
commutability of $R,E$ says that $R$ is reduced by the closed linear
manifold belonging to $E$ (cf.\ \Eref, Definition~13), exactly as one should
reasonably expect.} We may consequently form $\mathfrak M'$ even for a set
$\mathfrak M$ of entirely arbitrary operators (which need be neither linear
nor everywhere defined): it is the set of all $A\in\mathfrak B$ for which
both $A$ and $A^*$ commute with every $R\in\mathfrak M$. Thus
$\mathfrak M'$ is always a subset of $\mathfrak B$. In particular, for
every operator $R$ we have the sets
\[
  (R)',\ (R)'',\ldots\qquad
  \text{(all subsets of $\mathfrak B$).}
\]\footnote{\label{fn:23}It will surely cause no misunderstanding that, just as in
$R(A)$, we denote the set having the single element $R$ also by $R$.}

% [Source p. 377 / PDF 8]
Now one easily shows (cf.\
\hyperref[sec:ring-and-commutant-1]{\S\,II.1}) that an
\(A\in\mathfrak B\) is normal if
and only if $(A)''$ is Abelian. We extend this definition to arbitrary
conjugate operator pairs: they are to be called normal if $(R)''$ is
Abelian. In \hyperref[app:laurent-matrices]{Appendix~II} we shall see that this definition applies to broad
classes of operators. In particular, the theory to be developed below is
immediately applicable to the (complex, unbounded) Laurent matrices and
their generalizations.\footnote{\label{fn:24}Cf.\ Toeplitz, \emph{Mathematische Annalen}
\textbf{70} (1911), pp.~351--376.}

We shall then show that these normal operators, and only these, possess one
(and indeed only one) Hilbert spectral form with complex
eigenvalues.\footnote{\label{fn:25}\label{fn:normal-spectral-form}For a precise
formulation of this concept, which generalizes the usual (real) Hilbert
spectral form, cf.\ \Eref, Appendix~II.

The concept of normality for finite-dimensional matrices is due to
Frobenius [\emph{Journal f\"ur die reine und angewandte Mathematik}
\textbf{84} (1877), pp.~51--54]; he proved that precisely these matrices can
be transformed unitarily to diagonal form (with complex diagonal entries).
The same result was proved for the completely continuous operators of
Hilbert space in the encyclopedia article by Hellinger and Toeplitz (II C
13, pp.~1562--1563).

In \Eref, Appendix~I, the author proved that all bounded normal operators can
be put into the complex Hilbert spectral form (which is exactly the
equivalent of diagonal form). Here the theory is to be completed and
extended to unbounded operators. Independently of the author, and using a
different method, A.~Wintner has since also put the unitary operators (a
special case of bounded normal operators) into spectral form. Cf.\
\emph{Mathematische Zeitschrift} \textbf{30}, nos.~1/2 (1929),
pp.~228--282.} These investigations therefore have an essentially
different direction from the preceding sections and are to be regarded
rather as a continuation of E.

As one sees, the situation for normal operators is somewhat more favorable
than for H.\ operators: for the latter the Hilbert spectral form was not
always attainable,\footnote{\label{fn:26}Cf.\ \Eref, Introduction~VII.} and this is all the
more remarkable because, in finite-dimensional spaces and even for bounded
operators, H.\ was a special case of normal
($AA^*=A^*A$ follows from $A=A^*$). The explanation is that we shall prove
the following: an H.\ operator is, in essence, normal if and only if it is
hypermaximal, and precisely these H.\ operators have a Hilbert spectral form
(cf.\ \Eref, Theorem~36). Thus the fact that H.\ appears as a special case of
normal in the situations just named is due simply to the fact that
everything there is hypermaximal.

\medskip
\phantomsection
\noindent\textbf{5.}\label{intro:5}
The organization of this paper is as follows. In
\hyperref[sec:topologies]{Section~I} we present and
discuss the different ways of topologizing \(\mathfrak B\). In
\hyperref[sec:ring-and-commutant]{Section~II} we
introduce the operations $\mathfrak M'$ and $R(\mathfrak M)$ (formation of
a ring) and prove various properties of them (including the already
announced analogue of the Burnside--Frobenius--Schur theorem).

% [Source p. 378 / PDF 9]
In \hyperref[sec:iii]{Section~III} we discuss the Abelian rings (cf.\ what was said above about
them). \hyperref[sec:iv-normal-operators]{Sections~IV} and
\hyperref[sec:v-spectral-form]{V} are almost independent of the preceding material
(apart from the definitions); their content is the general theory of normal
operators (without an assumption of reality or boundedness) and their
spectral representation.

\hyperref[app:bounded-spectral-families]{Appendix~I} treats the bounded
H.\ operators; \hyperref[app:laurent-matrices]{Appendix~II} gives an
application of the theory of normal operators to Laurent matrices and more
general matrices. \hyperref[app:commutativity-unbounded]{Appendix~III}
shows the inapplicability, in the unbounded
case, of the usual definition of commutability of matrices.

\section{Topologizations of \texorpdfstring{$\mathfrak H$ and
$\mathfrak B$}{H and B}}\label{sec:topologies}

\subsection*{1.}\phantomsection\label{sec:topologies-1}
Before considering the topology of $\mathfrak B$, we shall describe the
somewhat simpler topological relations in $\mathfrak H$---less for the sake
of their later applications than in order to illustrate the situation in
$\mathfrak B$.

Up to now (as in \Eref) we have considered only the following topology in
$\mathfrak H$, which is called ``strong,''\footnote{\label{fn:27}Cf.\ Weyl, doctoral
dissertation, G\"ottingen, 1908, pp.~8--9.} and which, following Hausdorff,
we characterize by specifying all neighborhoods of an arbitrary point
$f_0\in\mathfrak H$.

\medskip
\noindent\textbf{\textit{Strong topology on $\mathfrak H$.}}
Let
\[
  \mathcal U_1(f_0;\epsilon)
    =\{f\in\mathfrak H:|f-f_0|<\epsilon\}.
\]
All $\mathcal U_1(f_0;\epsilon)$ with $\epsilon>0$ are the neighborhoods of
$f_0$.

These neighborhoods are therefore the interiors of balls about $f_0$. Since
$|f-f_0|$ is a distance (cf.\ \Eref, Definition~4 and note~27), it follows at
once that the four Hausdorff neighborhood axioms (which are postulated in
every topology) are satisfied, namely:\footnote{\label{fn:28}Cf.\ Hausdorff,
\emph{Grundz\"uge der Mengenlehre}, 1st ed., Leipzig and Berlin, 1914; the
entire topology is developed there from this concept of neighborhood.

We recall how accumulation points and limits are to be defined. A point is
an accumulation point of a set if each of its neighborhoods contains points
of the set; it is the limit of a sequence if each of its neighborhoods
contains all but finitely many terms of that sequence. It is easily seen
that being a limit means being an accumulation point of every infinite
subset of the sequence.}
\begin{enumerate}
  \item[$\alpha)$] A point belongs to each of its neighborhoods.
  \item[$\beta)$] The intersection of two neighborhoods of a point contains
  a neighborhood of that point.
  \item[$\gamma)$] Every point in a neighborhood itself has a neighborhood
  that is a subset of the first.
  \item[$\delta)$] Any two distinct points have disjoint neighborhoods.
\end{enumerate}

% [Source p. 379 / PDF 10]
We also know already from E.\ that the operations $af$, $f\pm g$, $(f,g)$,
and $|f|$ are continuous in this sense (and $f\pm g$ and $(f,g)$ are
continuous in both variables simultaneously).

In the strong topology the so-called first Hausdorff axiom of countability
holds:\footnote{\label{fn:29}Cf.\ the work cited in \NoteRef{28}. Because $\mathfrak H$ is
separable (\Eref, \S~I, Axiom~C), the second axiom of countability also holds;
cf.\ the same source.} for every $f_0$ there is a decreasing sequence of
neighborhoods such that every neighborhood of $f_0$ contains one from this
sequence; take
\[
  \mathcal U_1\left(f_0;\frac1n\right),\qquad n=1,2,\ldots.
\]
From this one easily infers that, for every accumulation point of a set in
$\mathfrak H$, there is a sequence of points of the set that has it as its
limit.

It is also customary to introduce another topology in $\mathfrak H$, the
so-called ``weak'' topology.\footnote{\label{fn:30}Weyl (cf.\ \NoteRef{27}) defines it
somewhat differently from us; cf.\ p.~9 of the work cited there.} Ordinarily
this means that only weak limits, and not neighborhoods, are defined:
$f_1,f_2,\ldots$ converges weakly to $f$ if, for every fixed $\varphi$,
\[
  (f_n,\varphi)\longrightarrow(f,\varphi).
\]
It is easy, however, to reduce this to a notion of neighborhood.

\medskip
\noindent\textbf{\textit{Weak topology on $\mathfrak H$.}}
Let
\[
 \mathcal U_2(f_0;\varphi_1,\ldots,\varphi_s,\epsilon)
 =
 \left\{
 f\in\mathfrak H:
 \begin{array}{c}
 |(f-f_0,\varphi_1)|<\epsilon,\ldots,\\[-2pt]
 |(f-f_0,\varphi_s)|<\epsilon
 \end{array}
 \right\}.
\]
All
$\mathcal U_2(f_0;\varphi_1,\ldots,\varphi_s,\epsilon)$, with arbitrary
$\varphi_1,\ldots,\varphi_s\in\mathfrak H$, arbitrary $s$, and
$\epsilon>0$, are the neighborhoods of $f_0$.

Again, one easily verifies Hausdorff's axioms $\alpha)$--$\delta)$ and the
fact that a limit in the sense of this topology is identical with the weak
limit defined above. The continuity of $af$ and $f\pm g$ (the latter in both
variables simultaneously) is evident. The continuity of $(f,g)$ in $f$,
with $g$ fixed, follows from the definition; because
\[
  (f,g)=\overline{(g,f)},
\]
it is likewise continuous in $g$, with $f$ fixed. It is not, however,
continuous in both variables simultaneously, for then
$|f|=\sqrt{(f,f)}$ would also be continuous, which is
false.\footnote{\label{fn:31}Indeed, every
 $\mathcal U_2(0;
 \varphi_1,\ldots,\varphi_s,\epsilon)$ still contains
 $f$ with $|f|=1$: it suffices to choose $f$ orthogonal to all
 $\varphi_1,\ldots,\varphi_s$ and of absolute value~1\TranslatorNoteMark.}
\TranslatorNoteText{The source prints \(f_0\) in
\(\mathcal U_2(f_0;\ldots)\); the orthogonality argument requires the
center \(0\).}

Since
\[
  \mathcal U_1(f_0;\epsilon)
  \subset
  \mathcal U_2(f_0;\varphi_1,\ldots,\varphi_s,\eta)
\]
for example when
\[
  \epsilon=
  \frac{\eta}{\max(|\varphi_1|,\ldots,|\varphi_s|)},
\]
every accumulation point or limit in the strong sense is also one in the
weak sense. The converse cannot hold, already because $|f|$ has different
continuity properties in the two topologies.

The first Hausdorff axiom of countability does not hold in the weak topology
of $\mathfrak H$, for its most important consequence fails: there is a set
and an accumulation point of it which is not the limit of any sequence of
points of that set. Thus, in contrast to what one is otherwise accustomed
to,
% [Source p. 380 / PDF 11]
weak closedness cannot be reduced to convergence properties.

Before giving the counterexample, we note one further fact: if
$f_1,f_2,\ldots$ converges weakly to $f$, then, by a remark of E.~Schmidt,
$|f_1|,|f_2|,\ldots$ are bounded.\footnote{\label{fn:32}%
\fontsize{9.5pt}{11.2pt}\selectfont
This also follows from a general
theorem of St.~Banach, \emph{Fundamenta Mathematicae} \textbf{3} (1922),
p.~157. For completeness we reproduce his elegant proof, specialized to
the present case.

Suppose, then, that $f_1,f_2,\ldots$ converges weakly to $f$; we must prove
that $|f_n|\leq c$ for some fixed $c$. It suffices to prove
\[
  |(f_n,\varphi)|\leq c|\varphi|
  \qquad\text{for all $\varphi$},
\]
for putting $\varphi=f_n$ gives the assertion. The displayed relation
certainly always holds if it holds for the $\varphi$ with
$|\varphi|=\epsilon$ ($\epsilon>0$), i.e.\ if $|(f_n,\varphi)|$ is bounded
for all $n$ and all $|\varphi|=\epsilon$; all the more so if this is true
for every $|\varphi|\leq\epsilon$, that is, throughout an arbitrary ball
about $0$.

If all $|(f_n,\varphi)|$ are bounded while $\varphi$ varies in a fixed ball
$\mathfrak K$ given by $|\varphi-\varphi_0|<\delta$, then they are also
bounded in a ball about $0$ (namely, for $|\varphi|\leq\delta$), since such
a $\varphi$ is certainly the difference of two points of $\mathfrak K$
(for example, $\varphi_0\pm\frac12\varphi$). Thus, if the theorem were
false, the quantities $|(f_n,\varphi)|$ would be unbounded in every ball
$\mathfrak K$: for every $c$ there would be a $\varphi_0\in\mathfrak K$
and an $n_0$ such that $|(f_{n_0},\varphi_0)|>c$. By continuity,
$|(f_{n_0},\varphi)|>c$ would then hold throughout a suitable ball about
$\varphi_0$, which may be chosen to lie inside $\mathfrak K$.

Now let $\mathfrak K$ be any ball; let $\mathfrak K'$ be a ball in
$\mathfrak K$ throughout which $|(f_{n'},\varphi)|>1$ ($n'$ fixed), with
radius less than $1$; let $\mathfrak K''$ be a ball in $\mathfrak K'$
throughout which $|(f_{n''},\varphi)|>2$ ($n''$ fixed), with radius less
than $\frac12$; and so forth. Since $\mathfrak H$ is complete (cf.\ \Eref,
\S~I, Axiom~E), all the balls
$\mathfrak K,\mathfrak K',\mathfrak K'',\ldots$ have a common point
$\overline\varphi$, and for this point
\[
 |(f_{n'},\overline\varphi)|>1,\quad
 |(f_{n''},\overline\varphi)|>2,\quad\ldots;
\]
hence
\[
  \limsup_{n\to\infty}|(f_n,\overline\varphi)|=\infty,
\]
contrary to the hypothesis.

We add the following observation. Let
$\varphi_1,\varphi_2,\ldots$ be a complete normalized orthogonal system. If
$f_1,f_2,\ldots$ converges weakly to $f$, then, for every
$\nu=1,2,\ldots$,
\[
  \lim_{n\to\infty}(f_n,\varphi_\nu)=(f,\varphi_\nu),
\]
and the $|f_n|$ are bounded. This necessary condition is also sufficient:
the set of all $\varphi$ for which
\[
  \lim_{n\to\infty}(f_n,\varphi)=(f,\varphi)
\]
contains $\varphi_1,\varphi_2,\ldots$ and is plainly a linear manifold.
Because the $|f_n|$ are bounded (so the $(f_n,\varphi)$ are equicontinuous
in $\varphi$), it is closed in the strong sense; hence it comprises all of
$\mathfrak H$. This proves everything asserted.}

The example is the set $\mathfrak A$ (we realize $\mathfrak H$ as the
sequence space of the $\{x_1,x_2,\ldots\}$) of all
$\{x_1,x_2,\ldots\}$ for which only two coordinates $x_m,x_n$ are
nonzero, with
\[
  x_m=1,\qquad x_n=m
  \qquad(m\neq n;\ m,n\ \text{otherwise arbitrary}).
\]
The point $0$ is a weak accumulation point of $\mathfrak A$. Indeed, given
\[
 \{u_1^1,u_2^1,\ldots\},\ldots,
 \{u_1^s,u_2^s,\ldots\},\epsilon,
\]
we can force
\[
 \left|\sum_{p=1}^{\infty}u_p^1x_p\right|<\epsilon,
 \ldots,
 \left|\sum_{p=1}^{\infty}u_p^sx_p\right|<\epsilon
 \qquad(\{x_1,x_2,\ldots\}\in\mathfrak A)
\]
as follows. First choose $m$ so that
\[
  |u_m^1|<\frac{\epsilon}{2},\ldots,
  |u_m^s|<\frac{\epsilon}{2}
\]
(since $u_p^1\to0,\ldots,u_p^s\to0$), and then choose $n$ so that
\[
  |u_n^1|<\frac{\epsilon}{2m},\ldots,
  |u_n^s|<\frac{\epsilon}{2m}.
\]
On the other hand, no sequence of points of $\mathfrak A$ converges weakly
to $0$. In such a sequence the absolute values, namely
$\sqrt{1+m^2}$, would be bounded, so only finitely many values of $m$
could occur. Consequently, for some $m_0$ one would have

% [Source p. 381 / PDF 12]
$m=m_0$ infinitely often, i.e.\ $x_{m_0}=1$, whereas weak convergence would
permit $x_{m_0}=1$ only finitely often.

Finally, attention should be drawn to the following fact. Inside any fixed
ball $\mathfrak K$, all $f$ are bounded, say $|f|\leq c$. Let
$\overline\varphi_1,\overline\varphi_2,\ldots$ be a complete normalized
orthogonal system. Then every
\[
  \mathcal U_2(f_0;\varphi_1,\ldots,\varphi_s,\epsilon)
\]
contains a
\[
  \mathcal U_2
  (f_0;\overline\varphi_1,\ldots,\overline\varphi_n,\delta),
\]
insofar as both are restricted to $\mathfrak K$. Indeed, write
\[
  \varphi_1=\sum_{\nu=1}^{\infty}a_\nu^1\overline\varphi_\nu,
  \quad\ldots,\quad
  \varphi_s=\sum_{\nu=1}^{\infty}a_\nu^s\overline\varphi_\nu,
\]
choose $n$ so that
\[
  \sum_{\nu=n+1}^{\infty}|a_\nu^1|^2,\ldots,
  \sum_{\nu=n+1}^{\infty}|a_\nu^s|^2
  <\frac{\epsilon^2}{16c^2},
\]
and then choose
\[
  \delta\leq
  \frac{\epsilon}
  {2\max\left(
      \sum_{\nu=1}^{n}|a_\nu^1|,\ldots,
      \sum_{\nu=1}^{n}|a_\nu^s|
    \right)}.
\]
For $r=1,\ldots,s$ and every $f$ in the latter neighborhood, we then have
\begin{align*}
 |(f-f_0,\varphi_r)|
 &\leq
 \sum_{\nu=1}^{n}|a_\nu^r|\,
       |(f-f_0,\overline\varphi_\nu)|
 +\left|
   \left(f-f_0,\sum_{\nu=n+1}^{\infty}
       a_\nu^r\overline\varphi_\nu\right)
 \right|\\
 &\leq
 \sum_{\nu=1}^{n}|a_\nu^r|\delta
 +|f-f_0|
 \sqrt{\sum_{\nu=n+1}^{\infty}|a_\nu^r|^2}
 <\frac{\epsilon}{2}+\frac{\epsilon}{2}
 =\epsilon.
\end{align*}
Thus these $f$ belong to the first neighborhood. Since we may plainly take
$\delta=1/p$ ($p=1,2,\ldots$), the first Hausdorff axiom of countability is
therefore satisfied.\footnote{\label{fn:33}Since there is a sequence dense everywhere in
$\mathfrak H$ in the strong, and hence also in the weak, sense, one easily
proves the second axiom of countability as well; cf.\ \NoteRef{29}.}

It is well known that the weak, but not the strong, topology is even compact
inside $\mathfrak K$.\footnote{\label{fn:34}All methods of selection in convergence
proofs in Hilbert space amount to this fact.} Thus the strong topology
always satisfies the axiom of countability, but is compact neither in all
of $\mathfrak H$ nor in $\mathfrak K$; the weak topology has neither
property in all of $\mathfrak H$, but has both in $\mathfrak K$.

\subsection*{2.}\phantomsection\label{sec:topologies-2}
After this preliminary orientation we turn to the relations in
$\mathfrak B$, which are our actual concern.

To define convergence of a sequence of operators $A_1,A_2,\ldots$ to an
operator $A$, one ordinarily requires either strong convergence of every
$A_nf$ to $Af$, or weak convergence; the latter means
\[
  (A_nf,g)\longrightarrow(Af,g)
  \qquad\text{for every $f,g$}.
\]
Thus one requires either
\[
  |(A_n-A)f|\longrightarrow0
  \quad\text{for every $f$},
\]
or
\[
  |((A_n-A)f,g)|\longrightarrow0
  \quad\text{for every $f,g$}.
\]
These are, respectively, strong and weak convergence of operators in
$\mathfrak B$. As in \hyperref[sec:topologies-1]{subsection 1}, we reduce
these notions of convergence to corresponding topologies in
$\mathfrak B$.

\medskip
\noindent\textbf{\textit{Strong topology on $\mathfrak B$.}}
Let
\[
 \mathcal U_3(A_0;\varphi_1,\ldots,\varphi_s,\epsilon)
 =
 \left\{
 A\in\mathfrak B:
 \begin{array}{c}
 |(A-A_0)\varphi_1|<\epsilon,\ldots,\\[-2pt]
 |(A-A_0)\varphi_s|<\epsilon
 \end{array}
 \right\}.
\]
All
$\mathcal U_3(A_0;\varphi_1,\ldots,\varphi_s,\epsilon)$

% [Source p. 382 / PDF 13]
with arbitrary $\varphi_1,\ldots,\varphi_s\in\mathfrak H$, arbitrary $s$,
and $\epsilon>0$ are the neighborhoods of $A_0$.

\medskip
\noindent\textbf{\textit{Weak topology on $\mathfrak B$.}}
Let
\[
 \mathcal U_4
 (A_0;\varphi_1,\psi_1,\ldots,\varphi_s,\psi_s,\epsilon)
\]
be the set of all $A\in\mathfrak B$ such that
\[
 |((A-A_0)\varphi_1,\psi_1)|<\epsilon,\ldots,
 |((A-A_0)\varphi_s,\psi_s)|<\epsilon.
\]
All these sets, with arbitrary
$\varphi_1,\psi_1,\ldots,\varphi_s,\psi_s\in\mathfrak H$, arbitrary $s$,
and $\epsilon>0$, are the neighborhoods of $A_0$.

One again readily verifies Hausdorff's axioms $\alpha)$--$\delta)$ and that
these topologies yield the notions of convergence just described.
Furthermore, it is clear that the operations $aA$ and $A+B$ are continuous
in both topologies (the latter in both variables simultaneously). The
operation $A^*$ is continuous in the weak topology, because
\[
  ((A-A_0)\varphi,\psi)
  =
  \overline{((A^*-A_0^*)\psi,\varphi)},
\]
but not in the strong topology (we omit the easily supplied
counterexamples). The product $AB$ is continuous in either topology in
each variable separately when the other is held fixed. Its strong
continuity in $A$ is seen by replacing
$\varphi_1,\ldots,\varphi_s$ with
$B\varphi_1,\ldots,B\varphi_s$; its strong continuity in $B$ follows from
the continuity of the fixed operator $A$. Its weak continuity in $A$
follows by replacing
\[
  \varphi_1,\psi_1,\ldots,\varphi_s,\psi_s
\]
with
\[
  B\varphi_1,\psi_1,\ldots,B\varphi_s,\psi_s;
\]
its weak continuity in $B$ follows from the preceding result and
\[
  AB=(B^*A^*)^*.
\]
In neither topology, however, is the product $AB$ continuous in both
variables simultaneously (here too the counterexamples are so simple that
we omit them).

Since
\[
  \mathcal U_3(A_0;\varphi_1,\ldots,\varphi_s,\epsilon)
  \subset
  \mathcal U_4
  (A_0;\varphi_1,\psi_1,\ldots,\varphi_s,\psi_s,\eta)
\]
for example when
\[
  \epsilon=
  \frac{\eta}{\max(|\psi_1|,\ldots,|\psi_s|)},
\]
every accumulation point or limit in the strong sense is also one in the
weak sense. The converse cannot hold, already because $A^*$ has different
continuity properties in the two topologies.

The first axiom of countability holds in neither of these topologies. In
fact, for each there are sets with an accumulation point which is not the
limit of any sequence of their points. We shall prove both assertions by
exhibiting a set $\mathfrak A$ of which $0$ is a strong accumulation point,
whereas no sequence of its points has $0$ even as a weak limit.

First we establish once more that if $A_1,A_2,\ldots$ converges weakly
(and hence, a fortiori, if it converges strongly) to $A$, then
$A_1,A_2,\ldots$ are uniformly bounded: there is a fixed $c$ such that
\[
  |A_nf|\leq c|f|
  \qquad\text{for every $n$ and $f$}.
\]
It suffices to prove this for weak convergence, where the proof is
analogous to that of the corresponding theorem in
$\mathfrak H$.\footnote{\label{fn:35}By \Eref, Theorem~12$\beta$, it suffices to prove
\[
  |(A_nf,g)|\leq C|f|\,|g|.
\]
We prove this by St.~Banach's method exactly as for the analogous relation
of E.~Schmidt in \NoteRef{32}, except that the function $(f_n,\varphi)$ of
$n,\varphi$ is replaced by the function $(A_nf,g)$ of $n,f,g$, and the
balls $\mathfrak K,\mathfrak K',\mathfrak K'',\ldots$ for $\varphi$ by
pairs of balls
\[
 (\mathfrak K,\mathfrak L),\
 (\mathfrak K',\mathfrak L'),\
 (\mathfrak K'',\mathfrak L''),\ldots
\]
for $f$ and $g$, respectively.}

% [Source p. 383 / PDF 14]
It follows incidentally that multiplication $AB$ is continuous in both
variables simultaneously with respect to strong \emph{convergence} (not the
topology). Indeed, from $A_n\to A$ and $B_n\to B$ it follows that
\[
  A_nBf\longrightarrow ABf.
\]
Moreover, $B_nf-Bf\to0$, and the $A_n$ are equicontinuous, so
\[
  A_nB_nf-A_nBf\longrightarrow0.
\]
Thus $A_nB_nf\to ABf$, i.e.\ $A_nB_n\to AB$ (all arrows being in the sense
of strong convergence in $\mathfrak H$ or $\mathfrak B$, as appropriate).
We now give the promised counterexample.

Let $\overline A_{m,n}$ be the following operator in $\mathfrak B$
($\mathfrak H$ again being realized as the sequence space of the
$\{x_1,x_2,\ldots\}$):
\[
  \overline A_{m,n}\{x_1,x_2,\ldots\}
  =\{y_1,y_2,\ldots\},
\]
where
\[
  y_m=x_m,\qquad y_n=mx_n,\qquad
  y_p=0\quad(p\neq m,n),
  \qquad m\neq n.
\]
Let $\mathfrak A$ be the set of all $\overline A_{m,n}$. The zero operator
is a strong accumulation point of $\mathfrak A$. Indeed, given
\[
  \{u_1^1,u_2^1,\ldots\},\ldots,
  \{u_1^s,u_2^s,\ldots\},\epsilon,
\]
we can force
\[
 |\overline A_{m,n}\{u_1^1,u_2^1,\ldots\}|<\epsilon,\ldots,
 |\overline A_{m,n}\{u_1^s,u_2^s,\ldots\}|<\epsilon,
\]
that is,\TranslatorNote{In the last radical in this estimate the source prints
\(m^2|u_m^s|^2\).  The correction to \(m^2|u_n^s|^2\) is required by the
subsequent choice of \(n\).}
\[
 \sqrt{|u_m^1|^2+m^2|u_n^1|^2}<\epsilon,\ldots,
 \sqrt{|u_m^s|^2+m^2|u_n^s|^2}<\epsilon,
\]
as follows. First choose $m$ so that
\[
 |u_m^1|<\frac{\epsilon}{\sqrt2},\ldots,
 |u_m^s|<\frac{\epsilon}{\sqrt2}
\]
(since $u_p^1\to0,\ldots,u_p^s\to0$), and then choose $n$ so that
\[
 |u_n^1|<\frac{\epsilon}{\sqrt2\,m},\ldots,
 |u_n^s|<\frac{\epsilon}{\sqrt2\,m}.
\]
On the other hand, no sequence of elements of $\mathfrak A$ converges
weakly to $0$. In such a sequence one would have uniformly
\[
  |\overline A_{m,n}f|\leq c|f|,
\]
so that $m$ would be bounded and only finitely many values of $m$ could
occur. Hence some $m_0$ would occur infinitely often. If
\[
  \varphi=\psi=\{x_1,x_2,\ldots\},
  \qquad x_{m_0}=1,\quad x_p=0\ (p\neq m_0),
\]
then $(\overline A_{m,n}\varphi,\psi)=1$ infinitely often, whereas weak
convergence would permit this equality only finitely often.

The discontinuity of $A^*$ in the strong topology is sometimes
inconvenient for our purposes. We therefore introduce \emph{strong double
convergence}: it obtains when
\[
  A_1,A_2,\ldots\longrightarrow A
  \quad\text{and}\quad
  A_1^*,A_2^*,\ldots\longrightarrow A^*
\]
strongly. It follows from what was said above that, with respect to strong
double convergence, the operations $aA,A^*,A+B,AB$ are continuous (the
last two in both variables simultaneously).

Finally we show the following.\TranslatorNote{The source
interchanges the two indices in this first convergence statement.  The
form printed here is the one used by the diagonal construction below.}
Suppose that $A_{m,n}$ converges strongly to
$A_n$ as $m\to\infty$, with $n$ fixed,
and that $A_n$ converges strongly to
$A$ as $n\to\infty$. If all $A_{m,n}$ are equicontinuous, i.e.\ there is a
fixed $c$ such that
\[
  |A_{m,n}f|\leq c|f|,
\]
then a suitable subsequence $A_{M_p,N_p}$ converges strongly to $A$ as
$p\to\infty$. Let $f_1,f_2,\ldots$ be a sequence everywhere dense in
$\mathfrak H$ in the strong sense (cf.\ \Eref, \S~I, Axiom~C). Choose $N_p$ so
that

% [Source p. 384 / PDF 15]
\[
 |A_{N_p}f_1-Af_1|<\frac1p,\ldots,
 |A_{N_p}f_p-Af_p|<\frac1p
\]
(since $A_nf\to Af$ always), and then choose $M_p$ so that
\[
 |A_{M_p,N_p}f_1-A_{N_p}f_1|<\frac1p,\ldots,
 |A_{M_p,N_p}f_p-A_{N_p}f_p|<\frac1p
\]
(since $A_{m,n}f\to A_nf$ with $n$ fixed and $m\to\infty$). It follows that
\[
  A_{M_p,N_p}f\longrightarrow Af
\]
for $f=f_1,f_2,\ldots$, i.e.\ on an everywhere dense subset of
$\mathfrak H$. But all the $A_{M_p,N_p}$, as well as $A$, are
equicontinuous; the conclusion therefore holds for every $f$, as claimed.
The same is plainly true for strong double convergence.

\subsection*{3.}\phantomsection\label{sec:topologies-3}
There is a third way to topologize $\mathfrak B$. Define convergence of
$A_1,A_2,\ldots$ to $A$ by requiring $|(A_n-A)f|$ to tend to $0$
uniformly, as $n\to\infty$, for all $f$; that is,
\[
  |(A_n-A)f|\leq c_n|f|,
  \qquad c_n\longrightarrow0.
\]
Equivalently, one may require $|((A_n-A)f,g)|$ to tend to $0$ uniformly,
as $n\to\infty$, for all $f,g$; that is,
\[
  |((A_n-A)f,g)|\leq c_n|f|\,|g|,
  \qquad c_n\longrightarrow0.
\]
By \Eref, Theorem~12$\beta$, the two conditions are identical. We call this,
with the natural terminology, \emph{uniform convergence}; as we see, it is
immaterial whether strong or weak convergence is sharpened in this way.
The corresponding topology is plainly the following.

\medskip
\noindent\textbf{\textit{Uniform topology on $\mathfrak B$.}}
Let $\mathcal U_5(A_0;\epsilon)$ be the set of all $A\in\mathfrak B$ for
which there is a fixed $\epsilon'<\epsilon$ such that\TranslatorNote{The
source omits the factor \(|f|\) from this displayed inequality; it is
restored from the definition of uniform convergence immediately above.}
\[
  |(A-A_0)f|<\epsilon'|f|
  \qquad\text{for every $f$}.
\]

All $\mathcal U_5(A_0;\epsilon)$ with $\epsilon>0$ are the neighborhoods of
$A_0$.

It is clear that Hausdorff's axioms $\alpha)$--$\delta)$ are satisfied and
that the resulting notion of convergence is the one just given. One also
sees that the operations $aA,A^*,A+B,AB$ are continuous (the last two in
both variables simultaneously).\footnote{\label{fn:36}For $A^*$, this follows because
\[
  |Af|\leq c|f|\ \text{for every $f$}
  \quad\Longrightarrow\quad
  |A^*f|\leq c|f|;
\]
these assertions are equivalent, respectively, to
\[
  |(Af,g)|\leq c|f|\,|g|,
  \qquad
  |(A^*f,g)|\leq c|f|\,|g|
\]
for every $f,g$ (cf.\ \Eref, Theorem~12$\beta$), and
$(A^*f,g)=\overline{(Ag,f)}$.}

Every strong neighborhood
$\mathcal U_3(A_0;\varphi_1,\ldots,\varphi_s,\epsilon)$ contains a
sufficiently small uniform neighborhood $\mathcal U_5(A_0;\eta)$.
\TranslatorNote{The displayed set inclusion in the source is reversed;
the verbal statement here gives the direction required by the two
topologies.}

Consequently every uniform accumulation point or limit is also one in the
strong (and hence, a fortiori, in the weak) sense.\footnote{\label{fn:37}Uniform
convergence also implies strong double convergence, because $*$ is
continuous for the former.} The converse is false, since the first axiom of
countability fails in both the strong and weak topologies, whereas, as we
shall now see, it holds in the uniform topology.

Since among the $\mathcal U_5(A_0;\epsilon)$ we may restrict ourselves to
\[
  \mathcal U_5\left(A_0;\frac1n\right),
  \qquad n=1,2,\ldots,
\]
the first axiom of countability holds. This topology, moreover, arises from
a metric on the linear space $\mathfrak B$: if we
% [Source p. 385 / PDF 16]
define as the modulus of \(A\) the upper bound of \(\lvert Af\rvert\)
for \(\lvert f\rvert=1\), that is, the smallest \(c\) for which
\(\lvert Af\rvert\leq c\lvert f\rvert\) always holds---it shall be
denoted by \(\lvert A\rvert\)---then, first, the following relations are
immediate:
\[
  \lvert aA\rvert=\lvert a\rvert\,\lvert A\rvert,\qquad
  \lvert A+B\rvert\leq\lvert A\rvert+\lvert B\rvert,\qquad
  \lvert AB\rvert\leq\lvert A\rvert\,\lvert B\rvert.
\]
Thus \(\lvert A-B\rvert\) is a distance in the linear space
\(\mathfrak B\) (cf.\ \Eref, \S~I, Axiom~B, and also Definition~4 and
\NoteRef{27}); and, second,
\(\mathcal U_5(A_0;\epsilon)\) is the set of all \(A\) for which
\(\lvert A-A_0\rvert<\epsilon\).  In other words, the uniform topology
arises from this metric.

As in \(\mathfrak H\), so also in \(\mathfrak B\), within every fixed
ball \(\mathfrak K\) (where all \(A\) satisfy \(\lvert A\rvert\leq c\)),
the character of the strong and weak topologies changes essentially:
there both satisfy the first axiom of countability.  Indeed, let
\(\overline\varphi_1,\overline\varphi_2,\ldots\) be a complete normalized
orthogonal system.  Then every
\(\mathcal U_3(A_0;\varphi_1,\ldots,\varphi_s,\epsilon)\), respectively
\[
  \mathcal U_4
  (A_0;\varphi_1,\psi_1,\ldots,\varphi_s,\psi_s,\epsilon),
\]
contains a
\[
  \mathcal U_3
  (A_0;\overline\varphi_1,\ldots,\overline\varphi_n,\delta),
\]
respectively a\TranslatorNote{Throughout this argument the source calls the
strong and weak neighborhoods \(\mathcal U_4\) and \(\mathcal U_5\),
respectively. The definitions given earlier require \(\mathcal U_3\) and
\(\mathcal U_4\), which are used here.}
\[
  \mathcal U_4
  (A_0;\overline\varphi_{k_1},\overline\varphi_{l_1},
  \ldots,\overline\varphi_{k_m},\overline\varphi_{l_m},\delta).
\]
In the first case put
\[
  \varphi_1=\sum_{\nu=1}^{\infty}a_\nu^1\overline\varphi_\nu,
  \ldots,\quad
  \varphi_s=\sum_{\nu=1}^{\infty}a_\nu^s\overline\varphi_\nu,
\]
and choose \(n\) so that
\[
  \sum_{\nu=n+1}^{\infty}\lvert a_\nu^1\rvert^2
  <\frac{\epsilon^2}{16c^2},\ldots,\quad
  \sum_{\nu=n+1}^{\infty}\lvert a_\nu^s\rvert^2
  <\frac{\epsilon^2}{16c^2},
\]
and choose\TranslatorNote{In this displayed choice of \(\delta\) for the
strong topology the source writes the upper summation limits as \(\infty\).
They have been corrected to \(n\), as required by the estimate below.}
\[
  \delta\leq
  \frac{\epsilon}{
    2\max\left(
      \sum_{\nu=1}^{n}\lvert a_\nu^1\rvert,\ldots,
      \sum_{\nu=1}^{n}\lvert a_\nu^s\rvert
    \right)}.
\]
Then, for \(r=1,\ldots,s\) and every \(A\) in the latter
\(\mathcal U_3\),
\[
\begin{aligned}
  \lvert(A-A_0)\varphi_r\rvert
  &\leq
  \sum_{\nu=1}^{n}\lvert a_\nu^r\rvert
    \lvert(A-A_0)\overline\varphi_\nu\rvert
  +\left\lvert
    (A-A_0)\sum_{\nu=n+1}^{\infty}
      a_\nu^r\overline\varphi_\nu
   \right\rvert                                                   \\
  &\leq
  \sum_{\nu=1}^{n}\lvert a_\nu^r\rvert\delta
  +2c\sqrt{\sum_{\nu=n+1}^{\infty}\lvert a_\nu^r\rvert^2}
  <\frac{\epsilon}{2}+\frac{\epsilon}{2}=\epsilon.
\end{aligned}
\]
Hence these \(A\) belong to the first-mentioned \(\mathcal U_3\).

In the second case put
\[
\begin{aligned}
  \varphi_1&=\sum_{\nu=1}^{\infty}a_\nu^1\overline\varphi_\nu,
  &\psi_1&=\sum_{\nu=1}^{\infty}b_\nu^1\overline\varphi_\nu,
  &\ldots\qquad
  \varphi_s&=\sum_{\nu=1}^{\infty}a_\nu^s\overline\varphi_\nu,
  &\psi_s&=\sum_{\nu=1}^{\infty}b_\nu^s\overline\varphi_\nu,
\end{aligned}
\]
and choose \(n\) so that
\[
\begin{gathered}
  \sum_{\nu=n+1}^{\infty}\lvert a_\nu^1\rvert^2<\eta^2,\qquad
  \sum_{\nu=n+1}^{\infty}\lvert b_\nu^1\rvert^2<\eta^2,\quad\ldots,\\
  \sum_{\nu=n+1}^{\infty}\lvert a_\nu^s\rvert^2<\eta^2,\qquad
  \sum_{\nu=n+1}^{\infty}\lvert b_\nu^s\rvert^2<\eta^2,
\end{gathered}
\]
where \(\eta>0\) has been chosen so that\TranslatorNote{In this choice and
in the estimate that follows, the source prints \(\eta^2\) in place of
\(2c\eta^2\).  Since \(\lvert A\rvert,\lvert A_0\rvert\leq c\), one has
\(\lvert A-A_0\rvert\leq2c\); the corrected coefficient \(2c\) is
therefore used here.}
\[
  4c\max(\lvert\varphi_1\rvert,\lvert\psi_1\rvert,\ldots,
          \lvert\varphi_s\rvert,\lvert\psi_s\rvert)\eta
          +2c\eta^2
  =\frac{\epsilon}{2}.
\]
Choose further\TranslatorNote{In this displayed choice of \(\delta\) for
the weak topology the source prints only the squared sums of the
\(a_\nu^r\). They have been corrected to the products
\((\sum_{\nu=1}^n\lvert a_\nu^r\rvert)
 (\sum_{\nu=1}^n\lvert b_\nu^r\rvert)\).}
\[
  \delta\leq
  \frac{\epsilon}{
    2\max\left(
      \left(\sum_{\nu=1}^{n}\lvert a_\nu^1\rvert\right)
      \left(\sum_{\nu=1}^{n}\lvert b_\nu^1\rvert\right),\ldots,
      \left(\sum_{\nu=1}^{n}\lvert a_\nu^s\rvert\right)
      \left(\sum_{\nu=1}^{n}\lvert b_\nu^s\rvert\right)
    \right)}.
\]
Finally, let \(m=n^2\), and let
\((k_1,l_1),\ldots,(k_m,l_m)\) be all pairs \((\rho,\sigma)\) with
\(\rho\leq n\), \(\sigma\leq n\).  Then, for \(r=1,\ldots,s\) and every
\(A\) in the latter \(\mathcal U_4\),\TranslatorNote{In the second tail
term of this estimate the source prints \(\psi_1\).  It has been corrected
to \(\psi_r\), as required by the subsequent bound.}

% [Source p. 386 / PDF 17]
\[
\begin{aligned}
&\left|
 ((A-A_0)\varphi_r,\psi_r)
 -
 \left(
   (A-A_0)\sum_{\nu=1}^{n}a_\nu^r\overline\varphi_\nu,
   \sum_{\nu=1}^{n}b_\nu^r\overline\varphi_\nu
 \right)
\right|                                                           \\
&=\left|
 -\left(
   (A-A_0)\varphi_r,
   \sum_{\nu=n+1}^{\infty}b_\nu^r\overline\varphi_\nu
 \right)
 -\left(
   (A-A_0)\sum_{\nu=n+1}^{\infty}a_\nu^r\overline\varphi_\nu,
   \psi_r
 \right)\right.                                                   \\
&\hspace{8em}\left.
 +\left(
   (A-A_0)\sum_{\nu=n+1}^{\infty}a_\nu^r\overline\varphi_\nu,
   \sum_{\nu=n+1}^{\infty}b_\nu^r\overline\varphi_\nu
 \right)
\right|                                                           \\
&\leq
 2c\lvert\varphi_r\rvert
 \sqrt{\sum_{\nu=n+1}^{\infty}\lvert b_\nu^r\rvert^2}
 +2c
 \sqrt{\sum_{\nu=n+1}^{\infty}\lvert a_\nu^r\rvert^2}
 \lvert\psi_r\rvert                                               \\
&\hspace{8em}
 +2c\sqrt{
   \left(\sum_{\nu=n+1}^{\infty}\lvert a_\nu^r\rvert^2\right)
   \left(\sum_{\nu=n+1}^{\infty}\lvert b_\nu^r\rvert^2\right)}\\
&<4c\max(\lvert\varphi_r\rvert,\lvert\psi_r\rvert)\eta+2c\eta^2
 \leq\frac{\epsilon}{2},
\end{aligned}
\]
and
\[
\begin{aligned}
 \left|
 \left(
   (A-A_0)\sum_{\nu=1}^{n}a_\nu^r\overline\varphi_\nu,
   \sum_{\nu=1}^{n}b_\nu^r\overline\varphi_\nu
 \right)
 \right|
 &\leq
 \sum_{\mu,\nu=1}^{n}
   \lvert a_\mu^r\rvert\,\lvert b_\nu^r\rvert
   \lvert((A-A_0)\overline\varphi_\mu,\overline\varphi_\nu)\rvert \\
 &\leq
 \sum_{\mu,\nu=1}^{n}
   \lvert a_\mu^r\rvert\,\lvert b_\nu^r\rvert\,\delta
 \leq\frac{\epsilon}{2}.
\end{aligned}
\]
Consequently
\[
  \lvert((A-A_0)\varphi_r,\psi_r)\rvert
  \leq\frac{\epsilon}{2}+\frac{\epsilon}{2}=\epsilon;
\]
that is, these \(A\) belong to the first-mentioned \(\mathcal U_4\).
Since we can plainly choose \(\delta=1/p\)
\((p=1,2,\ldots)\), the assertion above is proved.

\subsection*{4.}\phantomsection\label{sec:topologies-4}
In conclusion we note: every subset \(\mathfrak M\) of \(\mathfrak B\)
has a countable subset \(\mathfrak N\) such that every element \(A\) of
\(\mathfrak M\) is the limit of a strongly doubly convergent sequence
from \(\mathfrak N\).\footnote{\label{fn:38}Thus \(\mathfrak N\) is
everywhere dense in \(\mathfrak M\), both strongly and weakly; in
particular, we may put \(\mathfrak M=\mathfrak B\).  This is certainly
false for the uniform topology, since there one can readily give a set
of continuum many \(A\)'s whose pairwise distances are \(1\).  If we
take \(\mathfrak M\) to be a ball \(\mathfrak K\), then from the validity
of the first axiom of countability in \(\mathfrak K\), in both the strong
and weak senses, we can also readily infer that of the second
(cf.\ \NoteRef{29} and \NoteRef{33}).  We mention further that the weak
topology in \(\mathfrak K\) is compact
(cf.\ the end of \hyperref[sec:topologies-1]{subsection~1} and
\NoteRef{34}).}

Let \(\mathfrak M_c\) be the set of all \(A\in\mathfrak M\) with
\(\lvert A\rvert<c\).  It is enough to prove the assertion for all
\(\mathfrak M_1,\mathfrak M_2,\ldots\) (and take
\(\mathfrak N\) successively equal to
\(\mathfrak N_1,\mathfrak N_2,\ldots\)), since we can then put
\(\mathfrak N=\mathfrak N_1+\mathfrak N_2+\cdots\).  We may therefore
assume from the beginning that \(\lvert A\rvert<c\) throughout
\(\mathfrak M\).

We realize \(\mathfrak H\) as sequence space (the
\(\{x_1,x_2,\ldots\}\) with finite
\(\sum_{n=1}^{\infty}\lvert x_n\rvert^2\)).  Then \(A\) has a matrix
\(\{a_{\mu,\nu}\}\), since in general

% [Source p. 387 / PDF 18]
\[
  A\{x_1,x_2,\ldots\}=\{y_1,y_2,\ldots\},\qquad
  y_\nu=\sum_{\mu=1}^{\infty}a_{\mu,\nu}x_\mu
\]
(\(A\) is linear and continuous).  Let \(p=1,2,\ldots\), and let
 \(a_{\mu,\nu}^{(p)}\) (\(\mu,\nu=1,\ldots,p\)) be arbitrary complex
 numbers with rational real and imaginary parts (Gaussian rationals)
\TranslatorNote{The source says only ``rational numbers.''  In the
complex Hilbert space considered here this must mean Gaussian rationals;
real rational numbers alone cannot approximate arbitrary complex matrix
entries.}
 such that
\[
  \lvert a_{\mu,\nu}^{(p)}-a_{\mu,\nu}\rvert<\frac1p
  \qquad(\mu,\nu=1,\ldots,p).
\]
We define \(A^{(p)}\) by
\[
  A^{(p)}\{x_1,x_2,\ldots\}=\{y_1,y_2,\ldots\},\qquad
  y_\nu=
  \begin{cases}
    \displaystyle\sum_{\mu=1}^{p}a_{\mu,\nu}^{(p)}x_\mu,
      &\nu=1,\ldots,p,\\[0.7em]
    0,&\text{for all other \(\nu\).}
  \end{cases}
\]
For \(p\geq m\) we then have\TranslatorNote{In the middle line of this
estimate the source prints \(\mu=\nu+1\) as the lower limit of the final
sum.  Since the construction is truncated at \(p\), and the following
line has the same tail, this has been corrected to \(\mu=p+1\).}
\[
\begin{aligned}
  \lvert A^{(p)}\varphi_m-A\varphi_m\rvert^2
  &=
  \left\lvert
    \sum_{\mu=1}^{p}a_{m\mu}^{(p)}\varphi_\mu
    -\sum_{\mu=1}^{\infty}a_{m\mu}\varphi_\mu
  \right\rvert^2                                                   \\
  &=
  \left\lvert
    \sum_{\mu=1}^{p}
      (a_{m\mu}^{(p)}-a_{m\mu})\varphi_\mu
    -\sum_{\mu=p+1}^{\infty}a_{m\mu}\varphi_\mu
  \right\rvert^2                                                   \\
  &=
  \sum_{\mu=1}^{p}\lvert a_{m\mu}^{(p)}-a_{m\mu}\rvert^2
  +\sum_{\mu=p+1}^{\infty}\lvert a_{m\mu}\rvert^2
  <\frac1p+\sum_{\mu=p+1}^{\infty}\lvert a_{m\mu}\rvert^2.
\end{aligned}
\]
As \(p\to\infty\), this plainly tends to
\(0\),\footnote{\label{fn:39}Because
\(\sum_{\mu=1}^{\infty}\lvert a_{m\mu}\rvert^2\) converges.  The
convergence of all the series
\[
  \sum_{m=1}^{\infty}\lvert a_{mn}\rvert^2,\qquad
  \sum_{n=1}^{\infty}\lvert a_{mn}\rvert^2
\]
may be seen as follows.  The second converges because
\[
  A\{0,\ldots,0,\underset{m}{1},0,\ldots\}
  =\{a_{m1},a_{m2},\ldots\};
\]
the first then does as well if \(A^*\) is considered in place of \(A\)
(where \(a_{mn}\) is replaced by \(\overline a_{nm}\)).}
and this holds for every \(m=1,2,\ldots\).  Thus
\[
  A^{(p)}\varphi_m\longrightarrow A\varphi_m;
\]
in precisely the same way one shows
\[
  A^{(p)*}\varphi_m\longrightarrow A^*\varphi_m.
\]

Now consider all operators \(B\in\mathfrak B\) whose matrix
\(\{b_{\mu,\nu}\}\) consists entirely of Gaussian rationals, only finitely
many of which are nonzero.  They plainly form a countable set
\(B_1,B_2,\ldots\).  The \(A^{(p)}\) above are such \(B\)'s.  Thus, for
every \(A\in\mathfrak B\), we have a sequence
\(B_{\nu_1},B_{\nu_2},\ldots\) such that
\[
  B_{\nu_p}f\longrightarrow Af,\qquad
  B_{\nu_p}^*f\longrightarrow A^*f
\]
for all \(f=\varphi_1,\varphi_2,\ldots\).

We now pass to \(\mathfrak M\).  For each pair
\(\nu,p=1,2,\ldots\), we ask whether there is an
\(A\in\mathfrak M\) such that
\[
  \lvert B_\nu f-Af\rvert<\frac1p,\qquad
  \lvert B_\nu^*f-A^*f\rvert<\frac1p
  \quad\text{for }f=\varphi_1,\ldots,\varphi_p.
\]
If there is one, choose one such \(A\) and call it \(A_{\nu,p}\).  The
\(A_{\nu,p}\) form a countable subset \(\mathfrak N\) of
\(\mathfrak M\); we shall show that it has the required property.

% [Source p. 388 / PDF 19]
Let \(A\) be any element\TranslatorNote{At this point
the source prints
\(A\in\mathfrak B\).  The correction to \(A\in\mathfrak M\) is required
both by the assertion being proved and by the sentence below that uses
\(A\) itself to establish the existence of \(A_{n,p}\).} of
\(\mathfrak M\),
and let \(p=1,2,\ldots\).
Form the sequence \(B_{\nu_1},B_{\nu_2},\ldots\) just described.  For
this sequence all
\[
  \lvert B_{\nu_q}f-Af\rvert,\qquad
  \lvert B_{\nu_q}^*f-A^*f\rvert
\]
tend to \(0\) as \(q\to\infty\)
\((f=\varphi_1,\varphi_2,\ldots)\).  Thus, if we choose \(q\)
sufficiently large and set \(n=\nu_q\), then
\[
  \lvert B_nf-Af\rvert<\frac1p,\qquad
  \lvert B_n^*f-A^*f\rvert<\frac1p
\]
for all \(f=\varphi_1,\ldots,\varphi_p\).  Hence \(A_{n,p}\) can be
formed (since \(A\), for example, has the required property), and, for
all \(f=\varphi_1,\ldots,\varphi_p\),
\[
  \lvert B_nf-A_{n,p}f\rvert<\frac1p,\qquad
  \lvert B_n^*f-A_{n,p}^*f\rvert<\frac1p.
\]
Consequently\TranslatorNote{In the second \(2/p\) estimate the source
prints \(\lvert A_{n,p}^*f-Af\rvert\).  The second \(A\) has been corrected
to \(A^*\), as required by the preceding adjoint estimates and by the
conclusion \(\overline A_p^*f\to A^*f\).}
\[
  \lvert A_{n,p}f-Af\rvert<\frac2p,\qquad
  \lvert A_{n,p}^*f-A^*f\rvert<\frac2p.
\]
Write \(\overline A_p\) for \(A_{n,p}\) (where \(n\), too, depends on
\(p\)).  It is then plain that
\[
  \overline A_pf\longrightarrow Af,\qquad
  \overline A_p^*f\longrightarrow A^*f
\]
for all \(f=\varphi_1,\varphi_2,\ldots\), hence also for every linear
combination of finitely many of these vectors, that is, on a subset
everywhere dense in \(\mathfrak H\) in the strong sense.  Since all the
\(\overline A_p\) belong to \(\mathfrak N\), hence to \(\mathfrak M\),
and are therefore equicontinuous (uniformly in \(p\)), the convergence
holds for every \(f\).  Thus
\(\overline A_1,\overline A_2,\ldots\) converges strongly doubly to
\(A\), while it is a sequence from \(\mathfrak N\).

\section{Properties of
  \texorpdfstring{\(R(\mathfrak M)\) and \(\mathfrak M'\)}
  {R(M) and M'}}\label{sec:ring-and-commutant}

\subsection*{1.}\phantomsection\label{sec:ring-and-commutant-1}
We now turn to the definition of the operations \(R(\mathfrak M)\) and
\(\mathfrak M'\) already mentioned in
\hyperref[intro:2]{Introduction~2}.

\begin{definition}[Definition 1]\label{def:1}
A subset \(\mathfrak M\) of \(\mathfrak B\) is called a ring if, first,
with \(A,B\) it also contains
\[
  aA,\quad A^*,\quad A+B,\quad AB,
\]
and, second, it is weakly closed.\footnote{\label{fn:40}Observe that
weak closedness is more than strong closedness.}
\end{definition}

\begin{definition}[Definition 2]\label{def:2}
Since the intersection of arbitrarily many rings is again a ring, the
intersection of all rings containing a given subset
\(\mathfrak M\subset\mathfrak B\) is in particular a ring---the smallest
ring containing \(\mathfrak M\).  It shall be denoted by
\(R(\mathfrak M)\), the ring generated by \(\mathfrak M\).
\end{definition}

\begin{definition}[Definition 3]\label{def:3}
Let \(\mathfrak M\) be a subset of \(\mathfrak B\).  The set of all
\(A\in\mathfrak B\) for which \(A\) and \(A^*\) commute with every
\(B\in\mathfrak M\) (cf.\ \NoteRef{16}) shall be denoted by
\(\mathfrak M'\).  Iterates
\(\mathfrak M'',\mathfrak M''',\ldots\) may be formed in the same way.
\end{definition}

The set \(\mathfrak M'\) is always a ring.  That it contains
\(aA,A^*,A+B,AB\) along with \(A,B\) is clear; but it is also weakly
closed.  For if \(A\) is a weak accumulation point of
\(\mathfrak M'\) and \(B\in\mathfrak M\), then for every
\(A'\in\mathfrak M'\)
\[
  A'B=BA',\qquad A'^*B=BA'^*,
\]
and these four expressions are all continuous in \(A\) (weakly, with
\(B\) fixed).  Hence
\[
  AB=BA,\qquad A^*B=BA^*.
\]

If \(A\in\mathfrak M\) and \(B\in\mathfrak M'\), then \(B,B^*\)
commute with \(A\), and therefore \(A,A^*\) commute with \(B\); hence
\(A\in\mathfrak M''\).  Thus \(\mathfrak M\subset\mathfrak M''\).
It is further clear that
\[
  \mathfrak M\subset\mathfrak N
  \quad\Longrightarrow\quad
  \mathfrak M'\supset\mathfrak N'.
\]
Applying this to \(\mathfrak M\subset\mathfrak M''\) gives
\(\mathfrak M'\supset\mathfrak M'''\).  But if in the same inclusion
\(\mathfrak M\) is merely replaced by \(\mathfrak M'\), then
\(\mathfrak M'\subset\mathfrak M'''\).  Hence
\(\mathfrak M'=\mathfrak M'''\).  Replacing \(\mathfrak M\) successively
by \(\mathfrak M',\mathfrak M'',\ldots\), we obtain

% [Source p. 389 / PDF 20]
\[
  \mathfrak M'=\mathfrak M'''=\mathfrak M^{\mathrm V}=\cdots,
  \qquad
  \mathfrak M''=\mathfrak M^{\mathrm{IV}}
  =\mathfrak M^{\mathrm{VI}}=\cdots.
\]

Since \(\mathfrak M'\) is always a ring and plainly contains \(1\), the
same is true of \(\mathfrak M''\).  Because
\(\mathfrak M\subset\mathfrak M''\), it follows that
\[
  (\mathfrak M,1)\footnote{\label{fn:41}The union of
  \(\mathfrak M\) with \(1\).}\subset\mathfrak M'',
  \qquad R(\mathfrak M,1)\subset\mathfrak M''
\]
(and hence also \(R(\mathfrak M)\subset\mathfrak M''\)).  We shall soon
show (\TheoremRef{5}) that in fact
\[
  R(\mathfrak M,1)=\mathfrak M''.
\]

We call a set \(\mathfrak M\) Abelian if all its elements and their
adjoints commute with one another.  Thus, for sets which contain \(A^*\)
together with \(A\), for example rings, the commutativity of all the
elements suffices.  To be Abelian plainly means
\(\mathfrak M\subset\mathfrak M'\).  Since this implies
\(\mathfrak M''\subset\mathfrak M'''\), \(\mathfrak M''\) is Abelian
whenever \(\mathfrak M\) is; and because
\(\mathfrak M\subset\mathfrak M''\), the converse also holds.  From
\[
  \mathfrak M\subset R(\mathfrak M)\subset\mathfrak M''
\]
it follows that \(R(\mathfrak M)\) is Abelian simultaneously with these
sets.  Thus the Abelian character is equivalent for all three sets
\(\mathfrak M\), \(R(\mathfrak M)\), and \(\mathfrak M''\).

Let \(\mathfrak M\) be a ring.  If it is Abelian, then every
\(A\in\mathfrak M\) commutes with \(A^*\), that is, \(A\) is normal
(cf.\ \hyperref[intro:2]{Introduction~2}, at \NoteRef{16}).  If it is not Abelian, it
contains two noncommuting operators \(A,B\).  Put
\[
  A_1=\frac{A+A^*}{2},\qquad
  A_2=\frac{A-A^*}{2i},\qquad
  B_1=\frac{B+B^*}{2},\qquad
  B_2=\frac{B-B^*}{2i}.
\]
The \(A_1,A_2,B_1,B_2\) are Hermitian operators in \(\mathfrak M\) and
certainly do not all commute (since \(A=A_1+iA_2\),
\(B=B_1+iB_2\)).  We may therefore take \(A,B\) themselves to be
Hermitian.  Then \(C=A+iB\) belongs to \(\mathfrak M\) and fails to
commute with \(C^*=A-iB\); it is therefore not normal.  Consequently:
a ring \(\mathfrak M\) is Abelian if and only if all its elements are
normal.  An arbitrary \(\mathfrak M\) has the corresponding property if
and only if \(R(\mathfrak M)\) does.

\subsection*{2.}\phantomsection\label{sec:ring-and-commutant-2}
We shall now discuss the relations of the preceding notions to projection
operators and unitary operators.

\begingroup
\interfootnotelinepenalty=0
\begin{theorem}[Theorem 1]\label{thm:1}
Let \(\mathfrak M\) be a ring, let \(A\) be a bounded Hermitian operator,
and let \(E(\lambda)\) be its resolution of the
identity.\footnote{\label{fn:42}In \Eref\ (Definition~17), a resolution
of the identity was defined as a family of projection operators
\(E(\lambda)\), \(-\infty<\lambda<\infty\), with the following
properties:
\begin{enumerate}
  \item[\(\alpha)\)] If \(\lambda\leq\mu\), then \(E(\lambda)\) is a
  part of \(E(\mu)\).
  \item[\(\beta)\)] If \(\lambda\geq\lambda_0\) and
  \(\lambda\to\lambda_0\), then
  \(E(\lambda)f\to E(\lambda_0)f\) for every \(f\).
  \item[\(\gamma)\)] If \(\lambda\to-\infty\), respectively
  \(\lambda\to+\infty\), then always
  \(E(\lambda)f\to0\), respectively \(E(\lambda)f\to f\).
\end{enumerate}
All convergences, as always in \Eref, are strong.  We would now say in
\(\beta)\) that \(E(\lambda)\) is a right-continuous function of
\(\lambda\) in the strong topology of \(\mathfrak B\); and in
\(\gamma)\), that at \(\lambda=\pm\infty\) it can be defined
continuously in the strong topology of \(\mathfrak B\) as \(0\) and
\(1\), respectively.  \emph{The source reverses \(0\) and \(1\) in
this bracketed restatement; the order here is required by
\(\gamma)\) itself.}

The resolution of the identity \(E(\lambda)\) belongs to the (not
necessarily bounded) Hermitian operator \(A\) (cf.\ \Eref,
Theorem~36) if, in addition, the following holds.

\medskip
\noindent\textit{Spectral form.}
The expression \(Af\) has meaning if and only if
\[
  \int_{-\infty}^{\infty}
    \lambda^2\,d\lvert E(\lambda)f\rvert^2
\]
is finite.  This is a Stieltjes integral; since the integrand is always
nonnegative and, as can be shown, the expression following \(d\) never
decreases, the integral is either finite (convergent) or
\(+\infty\) (properly divergent), and here the former is required.  If
this is the case, then, for every \(g\),
\[
  (Af,g)=\int_{-\infty}^{\infty}
    \lambda\,d(E(\lambda)f,g),
\]
the integral being absolutely convergent.

It was established in \Eref\ that the so-called hypermaximal Hermitian
operators have such resolutions of the identity, and only they do;
indeed, all hypermaximal Hermitian operators and all resolutions of the
identity correspond one-to-one.  This was the generalization of
Hilbert's spectral form to the unbounded case
(cf.\ \Eref, Theorem~36).

For a bounded Hermitian operator \(A\), the resolution of the identity
always exists.  This is the result of Hilbert's theory; cf.\ also
\Eref, where it is shown that all bounded \(A\) are hypermaximal
(Theorem~48, \(\alpha\), \(\beta\)).  The resolution of the
identity of such a Hermitian operator is characterized by the fact that
condition \(\gamma)\) may be strengthened to say that \(E(\lambda)\) is
actually constant for all sufficiently small, respectively large,
\(\lambda\), and hence is \(0\), respectively \(1\).  If this occurs,
say, for \(\lambda\leq-c\), respectively \(\lambda\geq c\), then all
\(\int_{-\infty}^{\infty}\) above may be replaced by
\(\int_{-c}^{c}\) (and \(\gamma)\) is automatically satisfied).

This characterization of boundedness is one of the results of Hilbert's
theory.  For completeness only, we shall give a short proof of it in
\hyperref[app:bounded-spectral-families]{Appendix~I}.}

% [Source p. 390 / PDF 21]
\(A\) belongs to \(\mathfrak M\) if and only if all \(E(\lambda)\) with
\(\lambda<0\), and all \(1-E(\lambda)\) with \(\lambda\geq0\), belong
to \(\mathfrak M\).  If \(\mathfrak M\) contains \(1\), we may say more
simply: if and only if all \(E(\lambda)\) belong to \(\mathfrak M\).
\end{theorem}
\endgroup

\begin{proof}
The second assertion follows from the first, so we consider the first.

First, the condition is sufficient.  Suppose that all
\(E(\lambda)\), \(\lambda<0\), and all
\(1-E(\lambda)\), \(\lambda\geq0\), belong to \(\mathfrak M\), and let
\(-c\leq\lambda\leq c\) be the interval mentioned in \NoteRef{42},
outside which \(E(\lambda)\) is \(0\), respectively \(1\).  Let
\(K_\epsilon\) be a chain of numbers from \(-c\) to \(c\):
\[
  -c=\lambda_0<\lambda_1<\cdots<\lambda_{n-1}<\lambda_n=c,
\]
with mesh at most \(\epsilon\), that is,
\[
  \lambda_\nu-\lambda_{\nu-1}\leq\epsilon
  \quad(\nu=1,2,\ldots,n;\ \epsilon>0),
\]
and containing \(0\), say \(\lambda_m=0\).  Then
\[
\begin{aligned}
  A^{K_\epsilon}
  &=\sum_{\nu=1}^{n}
    \lambda_\nu\bigl(E(\lambda_\nu)-E(\lambda_{\nu-1})\bigr)\\
  &=\sum_{\nu=1}^{m-1}
    \lambda_\nu\bigl(E(\lambda_\nu)-E(\lambda_{\nu-1})\bigr)\\
  &\quad+
    \sum_{\nu=m+1}^{n}
    \lambda_\nu
    \bigl((1-E(\lambda_{\nu-1}))-(1-E(\lambda_\nu))\bigr)
\end{aligned}
\]
also belongs to \(\mathfrak M\), and
\[
  (A^{K_\epsilon}f,g)
  =\sum_{\nu=1}^{n}\lambda_\nu
   \bigl((E(\lambda_\nu)f,g)
        -(E(\lambda_{\nu-1})f,g)\bigr).
\]

% [Source p. 391 / PDF 22]
As \(\epsilon\to0\), this tends, for all \(f,g\), to
\[
  \int_{-c}^{c}\lambda\,d(E(\lambda)f,g)=(Af,g).
\]
Thus \(A\) is a weak limit of the \(A^{K_\epsilon}\),\footnote{
\label{fn:43}It is even a uniform limit of them.  Indeed, if
\(f_\epsilon(\lambda)\) is defined by
\[
  f_\epsilon(\lambda)=\lambda_\nu
  \qquad\text{for }\lambda_{\nu-1}<\lambda\leq\lambda_\nu,
\]
then
\[
\begin{aligned}
  (A^{K_\epsilon}f,g)
  &=\int_{-c}^{c}f_\epsilon(\lambda)\,d(E(\lambda)f,g),\\
  (Af,g)
  &=\int_{-c}^{c}\lambda\,d(E(\lambda)f,g),\\
  ((A^{K_\epsilon}-A)f,g)
  &=\int_{-c}^{c}(f_\epsilon(\lambda)-\lambda)\,
    d(E(\lambda)f,g).
\end{aligned}
\]
Now \(\lvert f_\epsilon(\lambda)-\lambda\rvert\leq\epsilon\) always.
Moreover,
\[
\begin{aligned}
&\lvert(E(\mu)f,g)-(E(\lambda)f,g)\rvert\\
&\qquad\leq
 \sqrt{\lvert E(\mu)f\rvert^2-\lvert E(\lambda)f\rvert^2}\,
 \sqrt{\lvert E(\mu)g\rvert^2-\lvert E(\lambda)g\rvert^2}
\end{aligned}
\]
(cf.\ the estimate carried out in \Eref\ in the proof of
Theorem~36).  Hence
\[
  \lvert((A^{K_\epsilon}-A)f,g)\rvert
  \leq\epsilon\lvert f\rvert\,\lvert g\rvert,
\]
and consequently (cf.\ \hyperref[sec:topologies-3]{\S~I, 3})
\(\lvert A^{K_\epsilon}-A\rvert\leq\epsilon\).  This proves the
assertion.}
and hence also a weak accumulation point of \(\mathfrak M\).
Consequently \(A\in\mathfrak M\).

Next, the condition is necessary.  Suppose now that
\(A\in\mathfrak M\).  Then \(A,A^2,A^3,\ldots\), and all their linear
combinations, also belong to \(\mathfrak M\); that is, so does every
\(p(A)\), where \(p(x)\) is any polynomial with \(p(0)=0\).  One proves
easily\footnote{\label{fn:44}In every detail this calculation is the same
as the one carried out for unitary operators in \Eref, Appendix~I, in
the proof of Theorem \(14^*\).  The device goes back to F.~Riesz.} that
\[
  (p(A)f,g)
  =\int_{-c}^{c}p(\lambda)\,d(E(\lambda)f,g).
\]

For \(\lambda_0<0\), choose \(\epsilon>0\) and a polynomial
\(p_\epsilon(x)\), with \(p_\epsilon(0)=0\), satisfying
\[
  \left.
  \begin{array}{rll}
    1-\epsilon<p_\epsilon(x)<1+\epsilon
      &\text{for}&-c\leq x\leq\lambda_0,\\
    -\epsilon<p_\epsilon(x)<1+\epsilon
      &\text{for}&\lambda_0\leq x\leq\lambda_0+\epsilon,\\
    -\epsilon<p_\epsilon(x)<\epsilon
      &\text{for}&\lambda_0+\epsilon\leq x\leq c.
  \end{array}
  \right\}
\]
(one sees without difficulty that this is possible).  Then, for every
\(f,g\),
\[
\begin{aligned}
 &(p_\epsilon(A)f-E(\lambda_0)f,g)\\
 &\quad=
 \int_{-c}^{c}p_\epsilon(\lambda)\,d(E(\lambda)f,g)
 -\int_{-c}^{\lambda_0}d(E(\lambda)f,g)\\
 &\quad=
 \int_{-c}^{\lambda_0}(p_\epsilon(\lambda)-1)\,
   d(E(\lambda)f,g)
 +\int_{\lambda_0}^{\lambda_0+\epsilon}p_\epsilon(\lambda)\,
   d(E(\lambda)f,g)\\
 &\qquad+
  \int_{\lambda_0+\epsilon}^{c}p_\epsilon(\lambda)\,
    d(E(\lambda)f,g).
\end{aligned}
\]

In these three integrals\TranslatorNote{The source prints \(-c\), rather
than \(c\), as the upper limit of the third integral.  The interval in
question and the polynomial bounds immediately above require \(c\).} the
absolute values of the integrands are at
most \(\epsilon\), \(1+\epsilon\), and \(\epsilon\), respectively.
Using the method of estimation from \NoteRef{43}, we may therefore infer

% [Source p. 392 / PDF 23]
\[
\begin{aligned}
&\lvert(p_\epsilon(A)f-E(\lambda_0)f,g)\rvert\\
&\leq
 \epsilon\lvert E(\lambda_0)f\rvert\,
          \lvert E(\lambda_0)g\rvert\\
&\quad +(1+\epsilon)
 \sqrt{\lvert E(\lambda_0+\epsilon)f\rvert^2
       -\lvert E(\lambda_0)f\rvert^2}\,
 \sqrt{\lvert E(\lambda_0+\epsilon)g\rvert^2
       -\lvert E(\lambda_0)g\rvert^2}\\
&\quad +\epsilon
 \sqrt{\lvert f\rvert^2-\lvert E(\lambda_0+\epsilon)f\rvert^2}\,
 \sqrt{\lvert g\rvert^2-\lvert E(\lambda_0+\epsilon)g\rvert^2}
 \footnotemark\\
&\leq
 \epsilon\lvert f\rvert\,\lvert g\rvert
 {}+
 \sqrt{\lvert E(\lambda_0+\epsilon)f\rvert^2
       -\lvert E(\lambda_0)f\rvert^2}\,
 \sqrt{\lvert E(\lambda_0+\epsilon)g\rvert^2
       -\lvert E(\lambda_0)g\rvert^2}\\
&\leq
 \epsilon\lvert f\rvert\,\lvert g\rvert
 {}+
 \sqrt{\lvert E(\lambda_0+\epsilon)f\rvert^2
       -\lvert E(\lambda_0)f\rvert^2}\,\lvert g\rvert.
\end{aligned}
\]
\footnotetext{\label{fn:45}As in \Eref, proof of Theorem~36
(cf.\ \NoteRef{43} above), the inequality
\[
  \sqrt{a_1}\sqrt{b_1}+\cdots+\sqrt{a_p}\sqrt{b_p}
  \leq
  \sqrt{(a_1+\cdots+a_p)(b_1+\cdots+b_p)}
\]
is used here.}
Since this holds for every \(g\), it follows that\footnote{
\label{fn:46}Put \(g=p_\epsilon(A)f-E(\lambda_0)f\).}
\[
  \lvert p_\epsilon(A)f-E(\lambda_0)f\rvert
  \leq\epsilon\lvert f\rvert
  +\sqrt{\lvert E(\lambda_0+\epsilon)f\rvert^2
        -\lvert E(\lambda_0)f\rvert^2}.
\]
This tends to \(0\) with \(\epsilon\to0\).  Thus, as
\(\epsilon\to0\), for example through the sequence
\(\epsilon=1,\frac12,\frac13,\ldots\),
\[
  p_\epsilon(A)f\longrightarrow E(\lambda_0)f
\]
strongly in \(\mathfrak H\), for every \(f\); hence
\(p_\epsilon(A)\) converges strongly in \(\mathfrak B\) to
\(E(\lambda_0)\).  Therefore \(E(\lambda_0)\in\mathfrak M\).

For \(\lambda_0\geq0\), choose polynomials \(q_\epsilon(x)\), with
\(q_\epsilon(0)=0\), satisfying
\[
  \left.
  \begin{array}{rll}
    -\epsilon<q_\epsilon(x)<\epsilon
      &\text{for}&-c\leq x\leq\lambda_0,\\
    -\epsilon<q_\epsilon(x)<1+\epsilon
      &\text{for}&\lambda_0\leq x\leq\lambda_0+\epsilon,\\
    1-\epsilon<q_\epsilon(x)<1+\epsilon
      &\text{for}&\lambda_0+\epsilon\leq x\leq c.
  \end{array}
  \right\}
\]
This, too, is plainly possible.  Entirely analogous considerations show
that, as \(\epsilon\to0\), the \(q_\epsilon(A)\) converge strongly in
\(\mathfrak B\) to \(1-E(\lambda_0)\).  In this case, therefore,
\(1-E(\lambda_0)\in\mathfrak M\).

This proves everything.
\end{proof}

\begin{theorem}[Theorem 2]\label{thm:2}
Let \(\mathfrak M\) be a ring, and let
\(\mathfrak M^P,\mathfrak M^U\) denote the sets of all projection
operators and all unitary operators, respectively, in \(\mathfrak M\).
Then always
\[
  R(\mathfrak M^P)=\mathfrak M.
\]
The set \(\mathfrak M^U\) is empty if \(1\notin\mathfrak M\); if
\(1\in\mathfrak M\), then
\[
  R(\mathfrak M^U)=\mathfrak M.
\]
\end{theorem}

\begin{proof}
By \TheoremRef{1}, two rings \(\mathfrak M,\mathfrak N\) are identical
if
\[
  \mathfrak M^P=\mathfrak N^P.
\]
For then they contain the same Hermitian operators, and hence the same
operators altogether, since membership of an arbitrary operator \(A\)
is characterized by that of the two Hermitian operators
\[
  \frac{A+A^*}{2},\qquad \frac{A-A^*}{2i}.
\]
Now
\[
  \mathfrak M^P\subset\mathfrak M,\qquad
  R(\mathfrak M^P)\subset R(\mathfrak M)=\mathfrak M,\qquad
  \bigl(R(\mathfrak M^P)\bigr)^P\subset\mathfrak M^P,
\]
while, on the other hand,
\[
  R(\mathfrak M^P)\supset\mathfrak M^P,\qquad
  \bigl(R(\mathfrak M^P)\bigr)^P\supset\mathfrak M^P.
\]
Thus
\[
  \bigl(R(\mathfrak M^P)\bigr)^P=\mathfrak M^P,
\]
and hence \(R(\mathfrak M^P)=\mathfrak M\).

% [Source p. 393 / PDF 24]
If \(U\) is a unitary element of \(\mathfrak M\), then \(U^*\) and
\(UU^*=1\) also belong to \(\mathfrak M\); hence, if \(1\) does not belong
to \(\mathfrak M\), then \(\mathfrak M^{\,U}\) is empty.  Suppose now that
\(1\in\mathfrak M\).  First,
\[
  \mathfrak M^{\,U}\subset\mathfrak M,
  \qquad
  R(\mathfrak M^{\,U})\subset R(\mathfrak M)=\mathfrak M.
\]
Secondly, let \(E\in\mathfrak M^{\,P}\).  Then \(E,1\), and hence
\(2E-1\), belong to \(\mathfrak M\); but the last operator is unitary,
because \(E\) is a projection operator,\footnote{\label{fn:47}Indeed,
\[
  (2E-1)^*=2E-1,
  \qquad
  (2E-1)^2=4E^2-4E+1=1.
\]}
and therefore belongs to \(\mathfrak M^{\,U}\).  Consequently
\[
  E=\tfrac12\bigl((2E-1)+1\bigr)\in R(\mathfrak M^{\,U}),
\]
and it follows that
\[
  \mathfrak M^{\,P}\subset R(\mathfrak M^{\,U}),
  \qquad
  R(\mathfrak M^{\,P})\subset R(\mathfrak M^{\,U}),
\]
so that \(R(\mathfrak M^{\,U})\supset\mathfrak M\).  Thus
\(R(\mathfrak M^{\,U})=\mathfrak M\), as was to be proved.
\end{proof}

\subsection*{3.}\phantomsection\label{sec:ii-3}
Let \(\mathfrak M\) be an arbitrary subset of \(\mathfrak B\).  Consider
all \(f\) for which
\[
  Af=0,\qquad A^*f=0
\]
for every \(A\in\mathfrak M\).  These \(f\) plainly form a closed linear
manifold.  Let \(E_0\) be the projection operator of its orthogonal
complement (cf.\ \Eref, Definition~12).  Membership in the original
manifold is therefore characterized by \(E_0f=0\).  We define:

\begin{definition}[Definition 4]\label{def:principal-unit}
Let \(\mathfrak M\subset\mathfrak B\) be arbitrary.  The projection
operator \(E_0\) just constructed, for which \(E_0f=0\) is equivalent to
\[
  Af=0,\qquad A^*f=0
  \quad\text{for every \(A\in\mathfrak M\),}
\]
is called the \emph{principal unit} of \(\mathfrak M\).
\end{definition}

We first show:

\begin{theorem}[Theorem 3]\label{thm:3}
For every \(A\in\mathfrak M\), and indeed for every
\(A\in R(\mathfrak M)\),
\[
  E_0A=AE_0=A.
\]
\end{theorem}

\begin{proof}
Let \(A\in\mathfrak M\).  Since identically
\(E_0(1-E_0)f=0\), we have
\[
  A(1-E_0)f=0,
\]
and therefore \(A(1-E_0)=0\), \(A=AE_0\).  Applying the same argument
to \(A^*\) gives \(A^*=A^*E_0\); taking adjoints then gives \(A=E_0A\).

Now consider all \(A\) satisfying \(E_0A=AE_0=A\).  They plainly form a
ring, and, as we have just seen, this ring contains \(\mathfrak M\).
It therefore contains \(R(\mathfrak M)\), which proves the assertion.
\end{proof}

\begin{theorem}[Theorem 4]\label{thm:4}
The principal unit \(E_0\) belongs to both \(\mathfrak M'\) and
\(\mathfrak M''\).
\end{theorem}

\begin{proof}
We have already seen that every \(A\in\mathfrak M\) commutes with \(E_0\);
moreover \(E_0^*=E_0\).  Hence \(E_0\in\mathfrak M'\).

Now let \(B\in\mathfrak M'\).  If \(E_0f=0\), then \(Af=0\) for every
\(A\in\mathfrak M\), and hence
\[
  BAf=ABf=0.
\]
Likewise \(A^*f=0\), and therefore
\[
  BA^*f=A^*Bf=0.
\]
Thus \(E_0Bf=0\).  Since \(E_0(1-E_0)f=0\) identically, it follows that
\[
  E_0B(1-E_0)f=0,
\]
that is,
\[
  E_0B(1-E_0)=0,
  \qquad
  E_0B=E_0BE_0.
\]
Take adjoints here and replace \(B\) by \(B^*\), which also belongs to
\(\mathfrak M'\).  We obtain
\[
  BE_0=E_0BE_0,
\]
and consequently \(E_0B=BE_0\).  This holds for every
\(B\in\mathfrak M'\), so \(E_0\in\mathfrak M''\).
\end{proof}

We now prove the theorem announced in paragraph~3 of the Introduction.

\begin{theorem}[Theorem 5]\label{thm:5}
\(R(\mathfrak M)\) is the set of all \(A\in\mathfrak M''\) for which
\[
  E_0A=AE_0=A.
\]
\end{theorem}

% [Source p. 394 / PDF 25]
\begin{proof}
We already know that \(R(\mathfrak M)\) is a subset of the indicated set;
it remains only to show that every such \(A\) actually belongs to
\(R(\mathfrak M)\).

Let \(\varphi_1,\ldots,\varphi_k\) be arbitrary elements of
\(\mathfrak H\), which we keep fixed for the moment, \(k=1,2,\ldots\).
At the same time, take another Hilbert space
\(\overline{\mathfrak H}\), and choose in it \(k\) infinite-dimensional
closed linear manifolds
\[
  \overline{\mathfrak H}_1,\ldots,\overline{\mathfrak H}_k
\]
which are mutually orthogonal and together span
\(\overline{\mathfrak H}\).\footnote{\label{fn:48}For the terminology,
cf., for example, \Eref, \S I.  The desired construction is as follows.
Let
\(\overline\varphi_1,\overline\varphi_2,\ldots\) be a complete normalized
orthogonal system, and introduce in it a double indexing
\[
  \overline\varphi_{ln},
  \qquad l=1,\ldots,k,\quad n=1,2,\ldots.
\]
Let \(\overline{\mathfrak H}_l\) be the closed linear manifold spanned by
\(\overline\varphi_{l1},\overline\varphi_{l2},\ldots\)
\((l=1,\ldots,k)\).  These
\(\overline{\mathfrak H}_1,\ldots,\overline{\mathfrak H}_k\) have all
the required properties.}
The spaces
\(\overline{\mathfrak H}_1,\ldots,\overline{\mathfrak H}_k\) are
themselves Hilbert spaces and hence are isomorphic to
\(\mathfrak H\).\footnote{\label{fn:49}Cf.\ \Eref, Introduction~III.}
Let
\[
  \Gamma_1,\ldots,\Gamma_k
\]
be isomorphic mappings of \(\mathfrak H\) onto, respectively,
\(\overline{\mathfrak H}_1,\ldots,\overline{\mathfrak H}_k\).

To each \(A\in\mathfrak B\) we now assign the point\footnotemark
\[
  \Gamma_1(A\varphi_1)+\cdots+\Gamma_k(A\varphi_k)
  \in\overline{\mathfrak H};
\]
\footnotetext{\label{fn:50}The expression
\(\Gamma_1f_1+\cdots+\Gamma_kf_k\) plainly runs through all of
\(\overline{\mathfrak H}\), and does so exactly once, when
\(f_1,\ldots,f_k\) range independently over \(\mathfrak H\).}
call it the \emph{image point} of \(A\).  Let
\(\overline{\mathfrak L}\) be the set of image points of the operators
\(A\in R(\mathfrak M)\), plainly a linear manifold, and let
\(\overline{\mathfrak M}\) be its closed hull (strongly in
\(\overline{\mathfrak H}\)), hence a closed linear manifold.  Let
\(\overline E\) denote the projection operator onto
\(\overline{\mathfrak M}\).

If \(A\) is a bounded operator in \(\mathfrak H\), there plainly exists
one and only one bounded operator \(\overline A\) in
\(\overline{\mathfrak H}\) such that
\[
  \overline A\bigl(\Gamma_1f_1+\cdots+\Gamma_kf_k\bigr)
  =
  \Gamma_1(Af_1)+\cdots+\Gamma_k(Af_k)
\]
(cf.\ \NoteRef{50}).  Plainly
\[
  \overline A^{\,*}=\overline{A^*}.
\]

Now let \(B\in R(\mathfrak M)\).  If
\(\overline f\in\overline{\mathfrak L}\), then
\[
  \overline f
  =
  \Gamma_1(A\varphi_1)+\cdots+\Gamma_k(A\varphi_k)
\]
for some \(A\in R(\mathfrak M)\), and
\[
  \overline B\,\overline f
  =
  \Gamma_1(BA\varphi_1)+\cdots+\Gamma_k(BA\varphi_k).
\]
Since \(BA\in R(\mathfrak M)\), this last point also belongs to
\(\overline{\mathfrak L}\).  Thus \(\overline B\) maps
\(\overline{\mathfrak L}\) into itself, and therefore maps
\(\overline{\mathfrak M}\) into itself.  Every
\(\overline E\overline f\) lies in \(\overline{\mathfrak M}\), hence so
does every \(\overline B\overline E\overline f\); identically,
\[
  \overline E\,\overline B\,\overline E\,\overline f
  =
  \overline B\,\overline E\,\overline f,
\]
and consequently
\[
  \overline E\,\overline B\,\overline E
  =
  \overline B\,\overline E.
\]
Replace \(B\) by \(B^*\), which likewise belongs to \(R(\mathfrak M)\),
and take adjoints.  Since
\(\overline B^{\,*}=\overline{B^*}\), this gives
\[
  \overline E\,\overline B\,\overline E
  =
  \overline E\,\overline B.
\]
Therefore
\[
  \overline E\,\overline B
  =
  \overline B\,\overline E.
\]

Let us examine \(\overline E\) more closely.  Plainly,
\[
  \overline E
  \bigl(\Gamma_1(f_1)+\cdots+\Gamma_k(f_k)\bigr)
  =
  \Gamma_1(g_1)+\cdots+\Gamma_k(g_k),
\]
where \(g_1,\ldots,g_k\) depend linearly and continuously upon
\(f_1,\ldots,f_k\).  Hence
\[
  g_l=E_{1l}f_1+\cdots+E_{kl}f_k
  \qquad(l=1,\ldots,k),
\]
where all the \(E_{jl}\)
\((j,l=1,\ldots,k)\) are bounded operators in
\(\mathfrak H\).\footnote{\label{fn:51}As is easily
calculated\TranslatorNoteMark,
\[
  E_{jl}
  =
  \Gamma_l^{-1}P_{\overline{\mathfrak H}_l}\,
  \overline E\,\Gamma_j,
\]
and conversely
\[
  \overline E
  =
  \sum_{j,l=1}^{k}
  \Gamma_lE_{jl}\Gamma_j^{-1}
  P_{\overline{\mathfrak H}_j}.
\]}
\TranslatorNoteText{In footnote~\ref{fn:51}, the source interchanges the
indices \(j\) and \(l\) in both displayed formulas.  The indices have
been corrected here in accordance with the preceding identity
\(g_l=E_{1l}f_1+\cdots+E_{kl}f_k\).}

% [Source p. 395 / PDF 26]
It follows that
\begin{align*}
  \overline E\,\overline B\,\overline f
  &=
  \overline E\,\overline B
  \bigl(\Gamma_1(f_1)+\cdots+\Gamma_k(f_k)\bigr)\\
  &=
  \overline E
  \bigl(\Gamma_1(Bf_1)+\cdots+\Gamma_k(Bf_k)\bigr)\\
  &=
  \Gamma_1(E_{11}Bf_1+\cdots+E_{k1}Bf_k)
  +\cdots+
  \Gamma_k(E_{1k}Bf_1+\cdots+E_{kk}Bf_k),
\end{align*}
whereas
\begin{align*}
  \overline B\,\overline E\,\overline f
  &=
  \overline B\,\overline E
  \bigl(\Gamma_1(f_1)+\cdots+\Gamma_k(f_k)\bigr)\\
  &=
  \overline B\Bigl(
    \Gamma_1(E_{11}f_1+\cdots+E_{k1}f_k)
    +\cdots+
    \Gamma_k(E_{1k}f_1+\cdots+E_{kk}f_k)
  \Bigr)\\
  &=
  \Gamma_1(BE_{11}f_1+\cdots+BE_{k1}f_k)
  +\cdots+
  \Gamma_k(BE_{1k}f_1+\cdots+BE_{kk}f_k).
\end{align*}
Thus we must have
\[
  E_{11}Bf_1+\cdots+E_{k1}Bf_k
  =
  BE_{11}f_1+\cdots+BE_{k1}f_k,
  \quad\ldots,
\]
\[
  E_{1k}Bf_1+\cdots+E_{kk}Bf_k
  =
  BE_{1k}f_1+\cdots+BE_{kk}f_k.
\]
That is, in general,
\[
  E_{jl}B=BE_{jl}.
\]
Since \(B\) may be any element of \(\mathfrak M\), while
\(B,B^*\in R(\mathfrak M)\), all the \(E_{jl}\) commute with \(B\);
therefore every \(E_{jl}\) belongs to \(\mathfrak M'\).

Membership of
\[
  \overline f=\Gamma_1(f_1)+\cdots+\Gamma_k(f_k)
\]
in \(\overline{\mathfrak M}\) is characterized by
\(\overline E\overline f=\overline f\), that is,
\[
  \overline E
  \bigl(\Gamma_1(f_1)+\cdots+\Gamma_k(f_k)\bigr)
  =
  \Gamma_1(f_1)+\cdots+\Gamma_k(f_k),
\]
\[
  \Gamma_1(E_{11}f_1+\cdots+E_{k1}f_k)
  +\cdots+
  \Gamma_k(E_{1k}f_1+\cdots+E_{kk}f_k)
  =
  \Gamma_1(f_1)+\cdots+\Gamma_k(f_k),
\]
and hence by
\[
  (E_{11}-1)f_1+\cdots+E_{k1}f_k=0,
  \quad\ldots,\quad
  E_{1k}f_1+\cdots+(E_{kk}-1)f_k=0.
\]
If \(\overline f\) is the image point of \(A\), so that
\[
  \overline f
  =
  \Gamma_1(A\varphi_1)+\cdots+\Gamma_k(A\varphi_k),
\]
these conditions become
\[
  (E_{11}-1)A\varphi_1+\cdots+E_{k1}A\varphi_k=0,
  \quad\ldots,\quad
  E_{1k}A\varphi_1+\cdots+(E_{kk}-1)A\varphi_k=0.
\]
In particular, if \(A\in\mathfrak M''\), so that it commutes with all
\(E_{jl}\) (since these belong to \(\mathfrak M'\)), these equations
become
\[
  A\bigl[(E_{11}-1)\varphi_1+\cdots+E_{k1}\varphi_k\bigr]=0,
  \quad\ldots,\quad
  A\bigl[E_{1k}\varphi_1+\cdots+(E_{kk}-1)\varphi_k\bigr]=0.
\]
Thus they are conditions of the form
\[
  A\omega_1=0,\quad\ldots,\quad A\omega_k=0,
\]
with fixed \(\omega_1,\ldots,\omega_k\).

For every \(A\in\mathfrak M\), both \(A\) and \(A^*\) belong to
\(R(\mathfrak M)\), and hence their image points belong to
\(\overline{\mathfrak M}\).  Moreover \(A,A^*\in\mathfrak M''\).
Consequently
\[
  A\omega_1=A^*\omega_1=\cdots
  =A\omega_k=A^*\omega_k=0.
\]
By \hyperref[def:principal-unit]{Definition~4}, it follows that
\[
  E_0\omega_1=\cdots=E_0\omega_k=0.
\]
Therefore, if \(A\in\mathfrak M''\) and at the same time
\(E_0A=AE_0=A\), then
\[
  A\omega_1=AE_0\omega_1=0,
  \quad\ldots,\quad
  A\omega_k=AE_0\omega_k=0.
\]

% [Source p. 396 / PDF 27]
In other words, the image point of \(A\) belongs to
\(\overline{\mathfrak M}\).  Hence, for every \(\varepsilon>0\), there
is an \(A'\in R(\mathfrak M)\) whose image point (which is the general
point of \(\overline{\mathfrak L}\)) has distance less than
\(\varepsilon\) from that of \(A\).  But
\begin{align*}
 &\bigl|
   \Gamma_1(A'\varphi_1)+\cdots+\Gamma_k(A'\varphi_k)
   -\Gamma_1(A\varphi_1)-\cdots-\Gamma_k(A\varphi_k)
 \bigr|^2\\
 &\qquad=
 \bigl|
   \Gamma_1((A'-A)\varphi_1)+\cdots+
   \Gamma_k((A'-A)\varphi_k)
 \bigr|^2\\
 &\qquad=
 \bigl|\Gamma_1((A'-A)\varphi_1)\bigr|^2+\cdots+
 \bigl|\Gamma_k((A'-A)\varphi_k)\bigr|^2\\
 &\qquad=
 \bigl|(A'-A)\varphi_1\bigr|^2+\cdots+
 \bigl|(A'-A)\varphi_k\bigr|^2,
\end{align*}
and consequently
\[
  \bigl|(A'-A)\varphi_1\bigr|<\varepsilon,
  \quad\ldots,\quad
  \bigl|(A'-A)\varphi_k\bigr|<\varepsilon.
\]
Thus \(A'\) belongs to the strong neighborhood
\(\mathcal U_3(A;\varphi_1,\ldots,\varphi_k;\varepsilon)\) of \(A\).

Since \(\varphi_1,\ldots,\varphi_k\) and \(\varepsilon>0\) were wholly
arbitrary, \(A\) is a strong accumulation point of \(R(\mathfrak M)\).
As this ring is even weakly closed, \(A\in R(\mathfrak M)\).  This proves
the theorem.
\end{proof}

The proof of \TheoremRef{5} also yields the following:

\begin{theorem}[Addendum]\label{thm:5-addendum}
Suppose that a set \(\mathfrak M\) has, of the ring properties
\(\alpha)\), \(\beta)\) in \DefinitionRef{1}, property \(\alpha)\), but in
place of \(\beta)\) only the following, less extensive property:
if \(A\) and \(A^*\) are strong accumulation points of \(\mathfrak M\),
then \(A\in\mathfrak M\).  This is weaker than strong closure, and still
more so than weak closure.  Then \(\mathfrak M\) is nevertheless a ring.
\end{theorem}

\begin{proof}
In the proof of \TheoremRef{5}, only the following properties of
\(R(\mathfrak M)\) were used:
\[
  R(\mathfrak M)\supset\mathfrak M,
  \qquad
  R(\mathfrak M)\subset\mathfrak M'',
\]
every \(A\in R(\mathfrak M)\) satisfies \(E_0A=AE_0=A\), and together
with \(A,B\), the operators
\[
  aA,\qquad A^*,\qquad A+B,\qquad AB
\]
also belong to it.  Under these assumptions it was shown that every
\(A\in\mathfrak M''\) with \(E_0A=AE_0=A\) is a strong accumulation
point of \(R(\mathfrak M)\).  The strong closure of \(R(\mathfrak M)\),
which was needed only to infer from this that \(A\) belongs to it, is
not to be invoked now.

For the present \(\mathfrak M\), we may therefore replace
\(R(\mathfrak M)\) by \(\mathfrak M\) itself.  We then obtain: every
\(A\in\mathfrak M''\) with \(AE_0=E_0A=A\) is a strong accumulation
point of \(\mathfrak M\).  Since \(A^*\) is such a point whenever \(A\)
is, \(A^*\) too is a strong accumulation point of \(\mathfrak M\).  By
the second half of the hypothesis, not used until now, \(A\in\mathfrak M\).

It follows from \TheoremRef{5} that
\[
  R(\mathfrak M)\subset\mathfrak M,
\]
and hence \(R(\mathfrak M)=\mathfrak M\).  Thus \(\mathfrak M\) is a ring,
as asserted.
\end{proof}

For sets having ring property \(\alpha)\) (\DefinitionRef{1}), strong and
weak closure are therefore equivalent.  E.~Schmidt proved the same
thing in \(\mathfrak H\) for linear manifolds;\footnote{\label{fn:52}Cf.\
\emph{Rendiconti del Circolo Matematico di Palermo}, \textbf{25} (1908),
pp.~57--73.}
our present result is its analogue in \(\mathfrak B\).

% [Source p. 397 / PDF 28]
\subsection*{4.}\phantomsection\label{sec:ii-4}
We have just characterized \(R(\mathfrak M)\) with the aid of \(E_0\)
and \(\mathfrak M''\), where \(\mathfrak M\) is again an arbitrary subset
of \(\mathfrak B\).  We now wish to reduce \(E_0,\mathfrak M''\) to
\(R(\mathfrak M)\).

\begin{theorem}[Theorem 6]\label{thm:6}
The principal unit \(E_0\) belongs to \(R(\mathfrak M)\); indeed, it is
the largest projection operator in \(R(\mathfrak M)\), in the sense that
every projection operator in \(R(\mathfrak M)\) is a part of \(E_0\).
\end{theorem}

\begin{proof}
Since \(E_0\in\mathfrak M''\) and \(E_0E_0=E_0\), it belongs to
\(R(\mathfrak M)\) by \TheoremRef{5}.  If \(E\) is a projection operator
in \(R(\mathfrak M)\), then \(EE_0=E\); that is, \(E\) is a part of
\(E_0\) (\Eref, Theorem~17).
\end{proof}

\begin{theorem}[Theorem 7]\label{thm:7}
\(\mathfrak M''\) is the set of all operators
\[
  A+a\cdot1,
\]
or, equivalently, of all operators
\[
  A+a(1-E_0),
\]
where \(A\in R(\mathfrak M)\) and \(a\) is complex.
\end{theorem}

\begin{proof}
Since \(E_0\in R(\mathfrak M)\), the two descriptions mean the same
thing.  Because
\[
  R(\mathfrak M)\subset\mathfrak M'',
  \qquad
  1\in\mathfrak M'',
\]
the indicated set is plainly contained in \(\mathfrak M''\).  It remains
to show that every \(B\in\mathfrak M''\) has the form
\[
  A+a(1-E_0),
  \qquad A\in R(\mathfrak M).
\]

Along with \(E_0,B\), the operator
\[
  E_0B=BE_0=A
\]
belongs to \(\mathfrak M''\), since \(E_0\in\mathfrak M'\) and
\(B\in\mathfrak M''\), so they commute.  Moreover,
\[
  E_0A=E_0^2B=E_0B=A,
  \qquad
  AE_0=BE_0^2=BE_0=A.
\]
By \TheoremRef{5}, \(A\in R(\mathfrak M)\).  Therefore
\[
  C=B-A
\]
also belongs to \(\mathfrak M''\), and
\[
  E_0C=E_0B-E_0A=A-A=0,
  \qquad
  CE_0=BE_0-AE_0=A-A=0.
\]
If we can infer from this that \(C=a(1-E_0)\), the proof is complete.

Let \(F\) be a projection operator which is a part of \(1-E_0\).  Then
\(F\) is orthogonal to \(E_0\) (cf.\ \Eref, Theorem~17), so
\[
  E_0F=FE_0=0.
\]
For every \(A\in R(\mathfrak M)\), therefore,
\[
  FA=FE_0A=0,
  \qquad
  AF=AE_0F=0.
\]
Thus \(A\) and \(F\) commute.  In particular, for every
\(A\in\mathfrak M\), both \(A\) and \(A^*\) commute with \(F\), so
\(F\in\mathfrak M'\).  Hence \(F\) and \(C\) commute.  Together with
\(E_0C=CE_0=0\), this implies the desired relation\footnotemark
\[
  C=a(1-E_0).
\]
\footnotetext{\label{fn:53}Here is the proof.  Let \(E_0f=0\), \(f\ne0\).
As is immediately seen, there corresponds to \(f\), by
\[
  Fg=\frac{1}{|f|^2}(g,f)f,
\]
a projection operator \(F\).  From \(E_0f=0\) it follows that
\(E_0F=0\); hence \(E_0,F\) are orthogonal, and \(F,C\) commute.  Thus
\[
  Cf=CFf=FCf
  =\frac{1}{|f|^2}(Cf,f)f
  =a_f f.
\]
We now show that the number \(a_f\) is the same for every \(f\).
Clearly \(a_f=a_{cf}\) for \(c\ne0\), so \(a_f=a_g\) when \(f,g\) are
linearly dependent.  If they are independent, then
\[
  C(f+g)=a_{f+g}(f+g),
  \qquad
  C(f+g)=Cf+Cg=a_ff+a_gg,
\]
and therefore
\[
  (a_f-a_{f+g})f+(a_g-a_{f+g})g=0.
\]
It follows that \(a_f=a_g=a_{f+g}\).  Hence \(a_f=a\) is indeed
constant, and \(Cf=af\) whenever \(E_0f=0\) (the case \(f=0\) is
plainly immaterial).

Since \(E_0(1-E_0)f=0\) identically and \(CE_0=0\), for every \(f\)
\[
  Cf=Cf-CE_0f=C(1-E_0)f=a(1-E_0)f.
\]
Thus \(C=a(1-E_0)\).}
\end{proof}

Finally we show:

\begin{theorem}[Theorem 8]\label{thm:8}
The rings containing \(1\), the sets \(\mathfrak M\) satisfying
\(\mathfrak M=\mathfrak M''\), and the sets \(\mathfrak N'\) are identical
classes.
\end{theorem}

% [Source p. 398 / PDF 29]
\begin{proof}
If a ring \(\mathfrak M\) contains \(1\), then its principal unit satisfies
\[
  E_0=E_0\,1=1.
\]
For every \(A\) we consequently have
\[
  E_0A=1A=A,
  \qquad
  AE_0=A1=A.
\]
By \TheoremRef{5},
\[
  R(\mathfrak M)=\mathfrak M'',
\]
and, since \(\mathfrak M\) is a ring, \(R(\mathfrak M)=\mathfrak M\).
Thus \(\mathfrak M=\mathfrak M''\).

If a set \(\mathfrak M\) satisfies \(\mathfrak M=\mathfrak M''\), then
\(\mathfrak M=\mathfrak N'\) with \(\mathfrak N=\mathfrak M'\).  Every set
\(\mathfrak N'\) is a ring and contains \(1\) (cf.\ the beginning of
paragraph~1).  This proves the equivalence of the three criteria.
\end{proof}

\section{The Abelian Rings}
\label{sec:iii}

\subsection*{1.}\phantomsection\label{sec:iii-1}
We have already seen that, under ring property \(\alpha)\)
(\DefinitionRef{1})\TranslatorNote{The source cites Definition~4 here;
\DefinitionRef{1} is the definition that introduces ring property
\(\alpha)\).}, weak and strong closure are the same.  We now consider
Abelian sets and shall show that for them the closure requirements may be
weakened still further.

Let \(\mathfrak M\) first be an arbitrary subset of \(\mathfrak B\).
Apply, arbitrarily but only finitely often, the operations
\[
  aA,\qquad A^*,\qquad A+B,\qquad AB
\]
to its elements.  The resulting set will be denoted by
\(r(\mathfrak M)\); plainly
\[
  \mathfrak M\subset r(\mathfrak M)\subset R(\mathfrak M).
\]
If to \(r(\mathfrak M)\) we adjoin all \(A\) for which both \(A\) and
\(A^*\) are strong accumulation points of \(r(\mathfrak M)\), we obtain
a set \(r_1(\mathfrak M)\).  It is contained in \(R(\mathfrak M)\), contains
\(\mathfrak M\), and, by the symmetry of the definition, contains \(A^*\)
whenever it contains \(A\).  Together with \(A,B\), it also contains
\[
  aA,\qquad A+B,\qquad AB,
\]
because \(r(\mathfrak M)\) does so and these operations are continuous in
the strong topology.  For \(AB\), continuity holds only in each variable
separately, but this suffices here, as is easily seen.

If \(B\) and \(B^*\) are strong accumulation points of
\(r_1(\mathfrak M)\), then they are also strong accumulation points of
\(r(\mathfrak M)\), since \(r_1(\mathfrak M)\) consists of accumulation
points of \(r(\mathfrak M)\).  Hence \(B,B^*\in r_1(\mathfrak M)\).
The addendum to \TheoremRef{5} is therefore applicable:
\(r_1(\mathfrak M)\) is a ring, and consequently
\[
  r_1(\mathfrak M)=R(\mathfrak M).
\]
Let \(r_2(\mathfrak M)\), respectively \(r_3(\mathfrak M)\), be obtained
from \(r(\mathfrak M)\) by adjoining all its strong, respectively all its
weak, accumulation points.  Then plainly
\[
  r_1(\mathfrak M)
  \subset r_2(\mathfrak M)
  \subset r_3(\mathfrak M)
  \subset R(\mathfrak M),
\]
and therefore
\[
  r_1(\mathfrak M)
  =
  r_2(\mathfrak M)
  =
  r_3(\mathfrak M)
  =
  R(\mathfrak M).
\]

If now \(\mathfrak M\) is Abelian, meaning that all elements of
\(\mathfrak M\), together with their adjoints, commute with one another,
then a still smaller enlargement of \(r(\mathfrak M)\) suffices to obtain
\(R(\mathfrak M)\), namely:

\begin{theorem}[Theorem 9]\label{thm:9}
If \(\mathfrak M\) is Abelian, then \(R(\mathfrak M)\) is obtained from
\(r(\mathfrak M)\) merely by adjoining all limits of strongly doubly
convergent sequences of its elements.\footnote{\label{fn:54}This is already
a substantial sharpening, since the first axiom of countability holds
neither in the strong nor in the weak topology; that is, not every
accumulation point is a limit\TranslatorNoteMark
(cf.\ \hyperref[sec:topologies-2]{\S\,I, 2}).  It could not be decided
whether this theorem also holds independently of the Abelian character of
\(\mathfrak M\).}
\end{theorem}
\TranslatorNoteText{In \NoteRef{54} the source omits the word ``not'' before
``every accumulation point is a limit.''  The first axiom of countability
fails here, so the negation is required.}

\begin{proof}
Call the resulting set \(\overline r(\mathfrak M)\).  Plainly
\[
  \mathfrak M\subset\overline r(\mathfrak M)\subset R(\mathfrak M).
\]
It is also plain that, together with \(A,B\), this set contains
\[
  aA,\qquad A^*,\qquad A+B,\qquad AB,
\]

% [Source p. 399 / PDF 30]
because \(r(\mathfrak M)\) has this property and these operations are
continuous with respect to strong double convergence.  Hence, if
\(\overline r(\mathfrak M)\) is strongly closed, it is a ring by the
addendum to \TheoremRef{5}, and therefore equals \(R(\mathfrak M)\).  It
remains only to prove strong closure.\footnote{\label{fn:55}Observe that
even the closure of \(\overline r(\mathfrak M)\) with respect to strong
double convergence is not self-evident.}

Let \(A\) be a strong accumulation point of
\(\overline r(\mathfrak M)\).  Then \(A\in R(\mathfrak M)\), and is
therefore normal (\hyperref[sec:topologies-1]{\S\,I, 1}).  Given arbitrary
\(\varphi_1,\ldots,\varphi_k\) and \(\varepsilon>0\), there is an
\(A'\in\overline r(\mathfrak M)\) in
\[
  \mathcal U_3(A;\varphi_1,\ldots,\varphi_k;\varepsilon);
\]
that is,
\[
  \lvert(A'-A)\varphi_1\rvert<\varepsilon,
  \quad\ldots,\quad
  \lvert(A'-A)\varphi_k\rvert<\varepsilon.
\]
Since \(A'\in\overline r(\mathfrak M)\), it too is normal.  Moreover,
\(A,A^*,A',A'^{*}\), as elements of the Abelian ring \(R(\mathfrak M)\),
all commute.  Put
\[
  \frac12(A+A^*)=B,
  \qquad
  \frac{1}{2i}(A-A^*)=C,
\]
\[
  \frac12(A'+A'^{*})=B',
  \qquad
  \frac{1}{2i}(A'-A'^{*})=C'.
\]
These are all commuting bounded Hermitian operators, with
\(B',C'\in\overline r(\mathfrak M)\); furthermore,
\[
  A'-A=(B'-B)+i(C'-C).
\]

If \(B_1,C_1\) are two commuting bounded Hermitian operators, then
\begin{align*}
  \lvert(B_1+iC_1)f\rvert^2
  &=
  \lvert B_1f\rvert^2
  +(B_1f,iC_1f)
  +(iC_1f,B_1f)
  +\lvert iC_1f\rvert^2\\
  &=
  \lvert B_1f\rvert^2
  +(-iC_1B_1f,f)
  +(iB_1C_1f,f)
  +\lvert C_1f\rvert^2\\
  &=
  \lvert B_1f\rvert^2+\lvert C_1f\rvert^2.
\end{align*}
The preceding inequalities therefore imply
\[
  \lvert(B'-B)\varphi_1\rvert<\varepsilon,
  \quad\ldots,\quad
  \lvert(B'-B)\varphi_k\rvert<\varepsilon,
\]
and
\[
  \lvert(C'-C)\varphi_1\rvert<\varepsilon,
  \quad\ldots,\quad
  \lvert(C'-C)\varphi_k\rvert<\varepsilon.
\]

Write
\[
  c=\max(\lvert B\rvert,\lvert C\rvert),
  \qquad
  c'=\max(\lvert B'\rvert,\lvert C'\rvert).
\]
Thus, for every \(f\),
\[
  \lvert Bf\rvert\leq c\lvert f\rvert,
  \quad
  \lvert Cf\rvert\leq c\lvert f\rvert,
  \quad
  \lvert B'f\rvert\leq c'\lvert f\rvert,
  \quad
  \lvert C'f\rvert\leq c'\lvert f\rvert.
\]
While \(c\) is fixed, \(c'\) depends upon
\(\varphi_1,\ldots,\varphi_k,\varepsilon\).  By increasing both if
necessary, we assume
\[
  c'\geq c>0.
\]

We now examine the relation between \(B\) and \(B'\) more closely.  Let
\(p(x)\) be a polynomial with the following properties:
\[
\left.
\begin{aligned}
  -c
  &<p(x)<\min\{c,-c+\delta\}
  &&\text{for } -c'\leq x\leq -c,\\
  \max\{x-\delta,-c\}
  &<p(x)<\min\{x+\delta,c\}
  &&\text{for } -c\leq x\leq c,\\
  \max\{-c,c-\delta\}
  &<p(x)<c
  &&\text{for } c\leq x\leq c',
\end{aligned}
\right\}
\]
% [Source p. 400 / PDF 31]
where \(\delta>0\) is still at our disposal.  We see that, for
\(\lvert x\rvert\leq c\),
\[
  \lvert p(x)-x\rvert<\delta;
\]
for \(\lvert x\rvert\leq c'\),
\[
  \lvert p(x)\rvert<c;
\]
and, for \(\lvert x\rvert,\lvert y\rvert\leq c'\),
\[
  \lvert p(x)-p(y)\rvert<\lvert x-y\rvert+2\delta,
\]
hence
\[
  \bigl(p(x)-p(y)\bigr)^2-(x-y)^2<8c'\delta+4\delta^2.
\]
We now choose \(\delta\) so that in the first inequality
\(\lvert p(x)-x\rvert<\varepsilon\), and in the last one
\[
  \bigl(p(x)-p(y)\bigr)^2-(x-y)^2<\varepsilon^2
\]
holds.  (The polynomial \(p(x)\) depends upon \(c'\), \(\delta\), and
\(\varepsilon\), and hence, directly and indirectly, upon
\(\varphi_1,\ldots,\varphi_k,\varepsilon\).)

A theorem of F.~Riesz (cf.\ \Eref, Appendix~II, Theorems \(3^*\),
\(4^*\)) gives, for every \(f\),
\[
  \lvert p(B)f-Bf\rvert\leq\varepsilon\lvert f\rvert,
  \qquad
  \lvert p(B')f\rvert\leq c\lvert f\rvert.
\]
Furthermore, the polynomial
\[
  \varepsilon^2+(x-y)^2-\bigl(p(x)-p(y)\bigr)^2=q(x,y)
\]
is everywhere nonnegative in the square
\(\lvert x\rvert,\lvert y\rvert\leq c'\).  From this, from the
commutativity of \(B,B'\), and from
\(\lvert B\rvert,\lvert B'\rvert\leq c'\), one may infer, by an analogue
of the theorem just used---to which we shall return presently---that the
Hermitian operator \(q(B,B')\) is definite.  That is,
\[
  \bigl(q(B,B')f,f\bigr)\geq0,
\]
\[
  \bigl((p(B')-p(B))^2f,f\bigr)
  \leq\bigl((B'-B)^2f,f\bigr)+\varepsilon^2(f,f),
\]
and therefore
\[
  \lvert(p(B')-p(B))f\rvert^2
  \leq\lvert(B'-B)f\rvert^2+\varepsilon^2\lvert f\rvert^2.
\]
For \(f=\varphi_1,\ldots,\varphi_k\) we consequently have
\[
  \lvert(p(B')-p(B))\varphi_i\rvert
  \leq\varepsilon\sqrt{1+\lvert\varphi_i\rvert^2}.
\]
Together with the preceding inequality this gives
\[
  \lvert(p(B')-B)\varphi_i\rvert
  \leq\varepsilon\Bigl(\lvert\varphi_i\rvert
  +\sqrt{1+\lvert\varphi_i\rvert^2}\Bigr).
\]
Thus, if we put
\[
  \varepsilon=
  \frac{\eta}{
    \max_{i=1,\ldots,k}
    \bigl(\lvert\varphi_i\rvert+\sqrt{1+\lvert\varphi_i\rvert^2}\bigr)},
  \qquad \eta>0\ \text{arbitrary},
\]
then \(p(B')\) lies in
\(\mathcal U_3(B;\varphi_1,\ldots,\varphi_k;\eta)\).  At the same time,
as we have seen, \(\lvert p(B')\rvert\leq c\).  Thus \(B\) is a strong
accumulation point of that part of \(r(\mathfrak M)\) which lies in the
ball \(\lvert B^*\rvert\leq c\), and, more precisely, of the Hermitian
operators in it.  By \hyperref[sec:topologies-3]{\S\,I, 3}, it is therefore also the limit of a strongly
convergent sequence in that set; and, since only Hermitian operators are
involved, this sequence is even doubly convergent.  In other words,
\[
  B\in\overline r(\mathfrak M).
\]
In exactly the same way one shows that \(C\) belongs to it as well, and
hence so does \(A=B+iC\).  This proves everything.

It remains to supply the proof, omitted above, of the generalization of
F.~Riesz's theorem.  Let \(B,B'\) be two commuting Hermitian operators
with \(\lvert B\rvert,\lvert B'\rvert\leq c'\).  By \Eref,
Appendix~II, Theorem \(11^*\) (put \(A=B+iB'\)), there are two
resolutions of the identity \(E(\lambda),F(\lambda)\) such that
\[
  (Bf,g)=\int_{-c'}^{c'}\lambda\,d(E(\lambda)f,g),
  \qquad
  (B'f,g)=\int_{-c'}^{c'}\lambda\,d(F(\lambda)f,g),
\]
and all the \(E(\lambda)\) and \(F(\mu)\) commute.  As there, we put
\[
  \Delta(\lambda,\mu)=E(\lambda)F(\mu)=F(\mu)E(\lambda)
\]

% [Source p. 401 / PDF 32]
and note that
\[
  \Delta(\lambda',\mu')\Delta(\lambda'',\mu'')
  =
  \Delta\bigl(\min(\lambda',\lambda''),\min(\mu',\mu'')\bigr).
\]
It follows from the preceding formulas (all double integrals are to be
extended over
\(\int_{-c'}^{c'}\int_{-c'}^{c'}\)) that
\[
  (Bf,g)=\iint\lambda\,d(\Delta(\lambda,\mu)f,g),
  \qquad
  (B'f,g)=\iint\mu\,d(\Delta(\lambda,\mu)f,g).
\]
Exactly as in the proof of \TheoremRef{1} (cf.\ the reference in
\NoteRef{44}), one then has
\[
  \bigl(q(B,B')f,g\bigr)
  =\iint q(\lambda,\mu)\,d(\Delta(\lambda,\mu)f,g),
\]
\[
  \bigl(q(B,B')f,f\bigr)
  =\iint q(\lambda,\mu)\,d(\Delta(\lambda,\mu)f,f)
  =\iint q(\lambda,\mu)\,d\lvert\Delta(\lambda,\mu)f\rvert^2.
\]
Now \(q(\lambda,\mu)\geq0\) throughout the region of integration, and
\[
  \iint_{\mathfrak R}d\lvert\Delta(\lambda,\mu)f\rvert^2\geq0
\]
for every rectangle \(\mathfrak R\) (cf.\ \Eref, Appendix~II,
\S5); hence the last integral on the right is nonnegative.  Since
\(\bigl(q(B,B')f,f\bigr)\geq0\) for every \(f\), the operator
\(q(B,B')\) is definite, as asserted.\footnote{\label{fn:56}Observe that in this
proof it was necessary to call upon the resolutions of the identity of
given Hermitian operators, whereas, conversely, their existence can be
proved from F.~Riesz's theorem.  For the Stieltjes--Radon double
integrals used here in \Eref, Appendix~II, see
\emph{Sitzungsberichte der Wiener Akademie}, \textbf{122}, IIa (1913),
pp.~1295--1438, especially \S\S I--II.}
\end{proof}

\subsection*{2.}\phantomsection\label{sec:iii-2}
A ring \(\mathfrak M=R(A)\) with one generator \(A\) is Abelian if and
only if it is the set \((A)\), that is, if and only if \(A,A^*\)
commute: hence this is so for normal \(A\).  We shall now prove the
converse announced in \hyperref[intro:3]{Introduction~3}: every Abelian ring (that is, every
ring having only normal elements) can be put in the form \(R(A)\), where
the normal \(A\) may even be chosen Hermitian.

\begin{theorem}[Theorem 10]\label{thm:10}
As \(A\) ranges over all normal operators, \(R(A)\) ranges over all
Abelian rings, and no others.  Moreover, for a given
\(\mathfrak M=R(A)\), \(A\) can always be chosen to be a Hermitian
operator.
\end{theorem}

\begin{proof}
In view of what was just said, it is enough to show that for every
Abelian ring \(\mathfrak M\) there exists a Hermitian operator \(A\) with
\(\mathfrak M=R(A)\).

First,
\[
  \mathfrak M=R(\mathfrak M^P) \qquad\text{(\TheoremRef{2}).}
\]
By \hyperref[sec:topologies-4]{\S\,I, 4} there is in
\(\mathfrak M^P\) a countable set
\(\mathfrak N\subset\mathfrak M^P\) which is strongly dense everywhere.
Consequently
\[
  \mathfrak M^P\subset R(\mathfrak N),\qquad
  \mathfrak M=R(\mathfrak M^P)\subset R(\mathfrak N),
\]
while, on the other hand,
\[
  R(\mathfrak N)\subset R(\mathfrak M^P)=\mathfrak M.
\]
Thus \(\mathfrak M=R(\mathfrak N)\).  Here \(\mathfrak N\) is a sequence
of projection operators \(E_1,E_2,\ldots\) belonging to \(\mathfrak M\);
these therefore commute.

In what follows we write \(E\leq F\), more briefly, to say that \(E\)
is a part of \(F\).  Two projection operators \(E,F\) are
\emph{comparable}\TranslatorNote{In the
second alternative the source prints
\(F\geq E\), which merely repeats \(E\leq F\).  It has been corrected
to \(F\leq E\), the reverse ordering required by comparability.} if
\(E\leq F\) or \(F\leq E\).
We have \(\mathfrak M=R(E_1,E_2,\ldots)\), where all the \(E_n\) commute.  We
shall now alter them, while preserving the preceding property, so that
they are even comparable.

% [Source p. 402 / PDF 33]
More precisely, we shall give a sequence \(F_1,F_2,\ldots\) of
projection operators with the following properties:
\[
  F_1=E_1;\qquad F_2\leq F_1\leq F_3,
\]
and \(F_1,F_2,F_3\) are obtainable from \(E_1,E_2\), and conversely, by
the operations \(+,-,\cdot\); further,
\[
  \cdots,\quad
  F_{2^{n-1}}\leq F'_1\leq F_{2^{n-1}+1}\leq F'_2\leq\cdots
  \leq F'_{2^{n-1}-2}\leq F_{2^n-2}
  \leq F'_{2^{n-1}-1}\leq F_{2^n-1}.
\]
Here \(F'_1,\ldots,F'_{2^{n-1}-1}\) are
\(F_1,\ldots,F_{2^{n-1}-1}\), arranged in the order of magnitude
determined by \(\leq\); and \(F_1,\ldots,F_{2^n-1}\) are obtainable
from \(E_1,\ldots,E_n\), and conversely, by the operations
\(+,-,\cdot\); and so on.  If we have such \(F_p\), it is clear that
each \(F_p\) is formed by \(+,-,\cdot\) from finitely many \(E_n\), and
conversely.  Hence
\[
  R(F_1,F_2,\ldots)=R(E_1,E_2,\ldots)=\mathfrak M,
\]
and all the \(F_p\) are manifestly comparable.  We shall now construct
them.

The construction is inductive.  For \(n=1\) it is trivial.  Suppose it
has already been carried out for \(n=m-1\); we carry it out for
\(n=m\).  Let \(F'_1,\ldots,F'_{2^{m-1}-1}\) again be
\(F_1,\ldots,F_{2^{m-1}-1}\) in the \(\leq\)-ordering, and put
\[
\begin{aligned}
  F_{2^{m-1}}&=E_mF'_1,\\
  F_{2^{m-1}+1}&=F'_1+E_m(F'_2-F'_1),\\
  F_{2^{m-1}+2}&=F'_2+E_m(F'_3-F'_2),\quad\ldots,\\
  F_{2^m-2}&=F'_{2^{m-1}-2}
    +E_m(F'_{2^{m-1}-1}-F'_{2^{m-1}-2}),\\
  F_{2^m-1}&=F'_{2^{m-1}-1}
    +E_m(1-F'_{2^{m-1}-1}).
\end{aligned}
\]
The preceding formulas express
\(F_{2^{m-1}},\ldots,F_{2^m-1}\) in terms of
\(F'_1,\ldots,F'_{2^{m-1}-1},E_m\), that is, in terms of
\(F_1,\ldots,F_{2^{m-1}-1},E_m\), and hence, by the induction hypothesis,
in terms of \(E_1,\ldots,E_{m-1},E_m\).  Conversely, among
the \(E_n\) only \(E_m\) remains to be considered, and for it we have
\[
  E_m=F_{2^{m-1}}+(F_{2^{m-1}+1}-F'_1)+\cdots+
  (F_{2^m-1}-F'_{2^{m-1}-1}).
\]
Finally, it is also easy to verify that
\[
  F_{2^{m-1}}\leq F'_1\leq F_{2^{m-1}+1}\leq F'_2\leq\cdots
  \leq F'_{2^{m-1}-2}\leq F_{2^m-2}
  \leq F'_{2^{m-1}-1}\leq F_{2^m-1}.
\]

Let us now consider \(F_1,F_2,\ldots\).  Their order is
\[
  F_2\leq F_1\leq F_3,
\]
\[
  F_4\leq F_2\leq F_5\leq F_1\leq F_6\leq F_3\leq F_7,
\]
\[
  \begin{aligned}
  F_8\leq F_4\leq F_9&\leq F_2\\
  &\leq F_{10}\leq F_5\leq F_{11}\leq F_1\\
  &\leq F_{12}\leq F_6\leq F_{13}\leq F_3
  \leq F_{14}\leq F_7\leq F_{15},\ldots .
  \end{aligned}
\]
This order can be characterized as follows.  We renumber the \(F_p\) by
writing \(F_\rho\) for
\[
  F_p,\qquad
  p=2^{l_1}+2^{l_2}+\cdots+2^{l_m},
  \qquad l_1>l_2>\cdots>l_m\geq0
\]
(the dyadic expansion), where
\[
  \rho=\frac{1}{2\cdot3^{l_1}}
       +\frac{2}{3^{l_1-l_2}}+\cdots
       +\frac{2}{3^{l_1-l_m}}.
\]
They then stand in their order of magnitude: \(\rho<\sigma\) implies
\(F_\rho\leq F_\sigma\).  As is readily seen, the index \(\rho\) ranges
over the set of all midpoints of the intervals omitted from the familiar
Cantor ternary set\footnote{\label{fn:57}The set in question comprises all numbers
\(\sum_{s=1}^{\infty}\alpha_s/3^s\), where every
\(\alpha_s\) is \(0\) or \(2\).  It is well known to be perfect and
nowhere dense, and it omits the intervals
\[
  \sum_{s=1}^{k}\frac{\alpha_s}{3^s}+\frac{1}{3^{k+1}}
  <x<
  \sum_{s=1}^{k}\frac{\alpha_s}{3^s}+\frac{2}{3^{k+1}}.
\]
Their midpoints are
\(\sum_{s=1}^{k}\alpha_s/3^s+1/(2\cdot3^k)\), another way of writing
our \(\rho\).}

% [Source p. 403 / PDF 34]
---that is, over a countable set lying in the interval \(0<x<1\) and
containing none of its own accumulation points.

We now construct a resolution of the identity \(G(\lambda)\) having the
following properties:
\[
  G(\lambda)=0\quad(\lambda<-1),\qquad
  G(\lambda)=1\quad(\lambda\geq0);
\]
and, for \(-1\leq\lambda<0\), every \(G(\lambda)\) is a strong
accumulation point of the family \(F_\rho\), with, in particular,
\[
  G(\rho-1)=F_\rho.
\]
Let \(\mathfrak N\) be the set of all \(G(\lambda)\) with \(\lambda<0\)
and all \(1-G(\lambda)\) with \(\lambda\geq0\).  Then
\(\mathfrak N\subset\mathfrak M^P\), and hence
\[
  R(\mathfrak N)\subset R(\mathfrak M)=\mathfrak M.
\]
On the other hand, \(\mathfrak N\) contains all the \(F_\rho\), that is,
\(F_1,F_2,\ldots\), and therefore
\[
  R(\mathfrak N)\supset R(F_1,F_2,\ldots)=\mathfrak M.
\]
Thus \(R(\mathfrak N)=\mathfrak M\).  If \(A\) is the bounded Hermitian
operator having the resolution of the identity \(G(\lambda)\) just
described (cf.\ \hyperref[app:bounded-spectral-families]{Appendix~I}), then \TheoremRef{1} gives
\[
  R(A)=R(\mathfrak N)=\mathfrak M,
\]
and everything is proved.  It remains, therefore, to construct
\(G(\lambda)\), and it is enough to define it for
\(-1\leq\lambda<0\).

In each open interval omitted from the Cantor set (cf.\ \NoteRef{57}), let
\[
  \overline G(\lambda-1)=F_\rho
\]
when \(\rho\) is the midpoint of that interval.  Thus
\(\overline G(\lambda)\) is defined on an open set dense everywhere in
\(-1\leq\lambda<0\); moreover,
\[
  \lambda<\mu\quad\Longrightarrow\quad
  \overline G(\lambda)\leq\overline G(\mu),
\]
and
\[
  \lambda\longrightarrow\lambda_0
  \quad\Longrightarrow\quad
  \overline G(\lambda)f\longrightarrow
  \overline G(\lambda_0)f
\]
whenever both sides are defined
(\(\overline G(\lambda)\) is strongly semicontinuous in
\(\lambda\)).  We may therefore extend its definition to the whole of
\(-1\leq\lambda<0\) as follows.  Let
\(\lambda_1>\lambda_2>\cdots\), \(\lambda_n\to\lambda\), be a sequence
for which
\(\overline G(\lambda_1),\overline G(\lambda_2),\ldots\) are already
defined.  Since
\[
  \overline G(\lambda_1)\geq\overline G(\lambda_2)\geq\cdots,
\]
the \(\overline G(\lambda_n)\) have a projection operator \(G\) as
their strong limit (cf.\ \Eref, Theorem~19).  The operator \(G\)
depends only upon \(\lambda\).  Indeed, from two such sequences
\[
  \lambda'_1>\lambda'_2>\cdots,\qquad
  \lambda''_1>\lambda''_2>\cdots
\]
we can combine subsequences into a new sequence
\[
  \lambda'_{m_1}>\lambda''_{n_1}>
  \lambda'_{m_2}>\lambda''_{n_2}>\cdots,
\]
which must also converge; hence the two limits are the same.  If
\(\overline G(\lambda)\) was already defined, then the preceding
remarks show that \(G=\overline G(\lambda)\).  In general we denote
\(G\) by \(G(\lambda)\).

For \(\mu<\lambda\), choose
\[
  \lambda_1>\lambda_2>\cdots,\quad\lambda_n\to\lambda,
  \qquad
  \mu_1>\mu_2>\cdots,\quad\mu_n\to\mu,
\]
with \(\mu_n<\lambda_n\).  Then
\(\overline G(\mu_n)\leq\overline G(\lambda_n)\), and consequently, by
continuity, \(G(\mu)\leq G(\lambda)\).  Furthermore, for a given \(f\)
and a suitable \(\lambda_n\),
\[
  \lvert\overline G(\lambda_n)f\rvert^2
  \leq\lvert G(\lambda)f\rvert^2+\varepsilon.
\]
Hence, for \(\lambda\leq\lambda'\leq\lambda_n\)
(\(\lambda_n>\lambda\)), all the more
\[
  \lvert G(\lambda')f\rvert^2
  \leq\lvert\overline G(\lambda_n)f\rvert^2
  \leq\lvert G(\lambda)f\rvert^2+\varepsilon,
\]
\[
  \lvert G(\lambda')f-G(\lambda)f\rvert^2
  =\lvert G(\lambda')f\rvert^2-\lvert G(\lambda)f\rvert^2
  \leq\varepsilon
\]
(cf.\ \Eref, Theorem~16).  Thus, for
\(\lambda'\geq\lambda\), \(\lambda'\to\lambda\),
\(G(\lambda')f\to G(\lambda)f\).  Consequently \(G(\lambda)\) is the
desired resolution of the identity, and we have reached our goal.
\end{proof}

% [Source p. 404 / PDF 35]
If \(\mathfrak M\) is an Abelian ring, then by \TheoremRef{10} just
proved it equals \(R(A)\), where \(A\) is a Hermitian operator; and by
\TheoremRef{9}, \(R(A)\) is the set of strong double limits of \(r(A)\).
In the present case, however, \(r(A)\), because \(A=A^*\), is the set
of all \(p(A)\), where \(p(x)\) ranges over all polynomials satisfying
\(p(0)=0\).  Thus \(\mathfrak M\) is the set of the limits of all
strongly doubly convergent sequences
\[
  p_1(A),p_2(A),\ldots
\]
(\(p_1(x),p_2(x),\ldots\) being polynomials vanishing at \(x=0\)).
With a more general interpretation of the notion \(f(A)\) than the one
we have used hitherto (namely, by extending it to discontinuous
\(f(x)\)), one could show that \(\mathfrak M\) is the set of all \(f(A)\),
where \(f(x)\) is any function vanishing at \(x=0\) for which the
extended definition applies.  We shall not enter into these matters
here.

\section{General Properties of (Unbounded) Normal Operators}
\label{sec:iv-normal-operators}

\subsection*{1.}\phantomsection\label{sec:iv-1}
We pass to our second subject, the study of unbounded operators and the
definition of normality for them.  We must therefore begin by considering
arbitrary operators (not necessarily everywhere defined or continuous;
for the present we do not even require them to be linear and closed),
which we shall denote by \(R,S,\ldots\), in contrast with the
\(A,B,\ldots\) of \(\mathfrak B\).

\begin{definition}[Definition 5]\label{def:5-operations}
The operator \(aR\) is defined for those \(f\) for which \(Rf\) is
defined, and its value is then \(a\cdot Rf\).  The operator \(R\pm S\)
is defined for those \(f\) for which \(Rf\) and \(Sf\) are defined, and
its value is then \(Rf\pm Sf\).  The operator \(RS\) is defined for
those \(f\) for which \(Sf\) and \(R(Sf)\) are defined, and its value is
then \(R(Sf)\).\footnote{\label{fn:58}It is clear that, with this definition, the
commutative and associative laws for addition and the associative law
for multiplication hold.  Of the distributive laws, one always has
\((R\pm S)T=RT\pm ST\), whereas \(R(S\pm T)=RS\pm RT\) only when
\(R\) is linear and everywhere defined (if \(R\) is merely linear, the
left-hand side is an extension of the right-hand side).  We have
\(R+0=R\), \(1R=R1=R\), and \(R0=0\) (provided \(Rf\) vanishes for
\(f=0\)); on the other hand, \(0R\) is defined only where \(R\) is.
The relation between \(+\) and \(-\) is analogous.  Thus, in the
unbounded case, the operations \(+,-,\cdot\) are no longer as
transparent as in the bounded case.}
\end{definition}

The operators \(R,A\) \emph{commute} (it is essential that
\(A\in\mathfrak B\)) if \(RA\) is an extension of \(AR\); that is, if,
whenever \(Rf\) is defined, \(R(Af)\) is also defined---\(Af\) and
\(A(Rf)\) are defined in any case---and is equal to \(A(Rf)\).\footnote{\label{fn:59}
According to \Eref, Introduction~V, \(R\) is an extension of \(S\)
when \(Rf\) is defined wherever \(Sf\) is, and the two are equal there.
If \(R\) too is everywhere defined (for example, if \(R\in\mathfrak B\)),
the present definition of commutativity says \(RA=AR\), that is, it is
ordinary commutativity.  For arbitrary \(R\), it would scarcely be
possible to require \(RA=AR\), since then \(R\) would not commute with
\(0\) (cf.\ \NoteRef{21} and \NoteRef{58}).  The asymmetry of our definition
(between \(R\) and \(A\)) shows that a direct definition of
commutativity for arbitrary \(R,S\) must encounter difficulties.}

% [Source p. 405 / PDF 36]
We can now extend \DefinitionRef{3} to sets \(\mathfrak M\) of arbitrary
operators.

\begin{definition}[Definition \(3'\)]\label{def:3-prime}
If \(\mathfrak M\) is a set of arbitrary operators (no longer necessarily
a subset of \(\mathfrak B\)), let \(\mathfrak M'\) be the set of all
\(A\in\mathfrak B\) for which \(A\) and \(A^*\) commute with every
\(R\in\mathfrak M\).
\end{definition}

Thus in every case \(\mathfrak M'\subset\mathfrak B\), and
\(\mathfrak M'',\mathfrak M''',\ldots\) consequently always fall under
the old definition.  Which of the properties of \(\mathfrak M'\) are
still retained?

It is immediately seen that \(1\) always belongs to \(\mathfrak M'\),
and, together with \(A,B\), so do \(A^*\) and \(AB\).  If, in addition,
all the elements of \(\mathfrak M\) are linear, then \(aA\) and \(A+B\)
also belong to it whenever \(A,B\) do.  Now suppose that all the
elements of \(\mathfrak M\) are closed, that \(A\) is a strong
accumulation point of \(\mathfrak M'\), and that \(R\in\mathfrak M\).
If \(Rf\) is defined, there is an \(A'_n\in\mathfrak M'\) in
\[
  \mathcal U_3\left(A;f,Rf;\frac1n\right);
\]
that is,
\[
  A'_nf\longrightarrow Af,\qquad
  RA'_nf=A'_nRf\longrightarrow ARf.
\]
Since \(R\) is closed, \(R(Af)\) is therefore defined and equals
\(ARf\).  Thus \(A\) and \(R\) commute.

Consequently, if \(A\) and \(A^*\) are strong accumulation points of
\(\mathfrak M'\), then \(A\in\mathfrak M'\).  Let us now assume, once and
for all, that every element of \(\mathfrak M\) is closed and linear.
Then, by the foregoing, all the hypotheses of the supplement to
\TheoremRef{5} are satisfied, and \(\mathfrak M'\) is a ring.  Since \(1\)
plainly belongs to it,
\[
  \mathfrak M'=\mathfrak M''' \qquad\text{(\TheoremRef{8}).}
\]

For arbitrary \(\mathfrak M\), we may apply the results from the beginning
of \hyperref[sec:ring-and-commutant-1]{\S\,II, 1} to
\(\mathfrak M'\subset\mathfrak B\):
\[
  \mathfrak M'\subset\mathfrak M'''=\mathfrak M^{\mathrm V}
  =\mathfrak M^{\mathrm{VII}}=\cdots,
  \qquad
  \mathfrak M''=\mathfrak M^{\mathrm{IV}}
  =\mathfrak M^{\mathrm{VI}}=\cdots.
\]
We cannot, however, conclude that \(\mathfrak M'=\mathfrak M'''\), since
the relation \(\mathfrak M\subset\mathfrak M''\) is missing
(\(\mathfrak M'\subset\mathfrak B\), but not necessarily
\(\mathfrak M\subset\mathfrak B\)).  As we saw, the equality does hold
when \(\mathfrak M\) consists entirely of closed linear operators.

Finally, let us record which elements occur in
\(\mathfrak M'^{\,P}\) and \(\mathfrak M'^{\,U}\).  Every projection
operator \(E\) commuting with all \(R\in\mathfrak M\) belongs to
\(\mathfrak M'^{\,P}\) (since \(E=E^*\)); that is, every \(E\) which
reduces all the \(R\in\mathfrak M\) (cf.\ \Eref, Definition~13).
Every unitary \(U\) for which both \(U\) and \(U^*=U^{-1}\) commute with
every \(R\in\mathfrak M\) belongs to \(\mathfrak M'^{\,U}\).  In other
words, \(RU\) is an extension of \(UR\), and \(RU^{-1}\) is an extension
of \(U^{-1}R\).  It is immediately seen that this means
\[
  RU=UR
\]
(with identical domains as well), or, equivalently,
\[
  U^{-1}RU=R.
\]

\subsection*{2.}\phantomsection\label{sec:iv-2}
We now introduce the notion of a conjugate pair (briefly, a
\emph{conjugate pair}) of operators.

\begin{definition}[Definition 5]\label{def:conjugate-pair}
Two operators \(R,R^*\) form a conjugate pair if they have the same
domain, if that domain spans the closed linear manifold
\(\mathfrak H\), and if throughout it
\[
  (Rf,g)=(f,R^*g),\qquad
  (f,Rg)=(R^*f,g)
\]
holds (cf.\ the beginning of \NoteRef{12}).
\end{definition}

% [Source p. 406 / PDF 37]
It is immediately seen that, together with \(R,R^*\), the operators
\(R^*,R\) also form a conjugate pair.  One \(R\) cannot form conjugate
pairs with two different operators \(R_1^*,R_2^*\): the domains of
\(R_1^*,R_2^*\) are prescribed to be the same, as are
\((R_1^*f,g)\) and \((R_2^*f,g)\) for a set of \(g\)'s spanning the
closed linear manifold \(\mathfrak H\), and hence \(R_1^*f\) and
\(R_2^*f\) themselves.  Nor can two different \(R_1,R_2\) have the same
\(R^*\), by the first observation.  The symbol \(^*\) is justified by
noting that, for \(A\in\mathfrak B\), precisely \(A,A^*\) form a
conjugate pair.  The operator \(R\) is Hermitian if and only if \(R,R\)
form a conjugate pair.

A conjugate pair \(S,S^*\) is an extension of another pair \(R,R^*\)
if the common domain of \(S,S^*\) contains that of \(R,R^*\), and if on
the latter\footnote{\label{fn:60}Each of the two following relations follows from the
other: the restrictions of \(S,S^*\) to the domain of \(R,R^*\) still
form a conjugate pair.}
\[
  Rf=Sf,\qquad R^*f=S^*f.
\]
If the two pairs are not identical, we speak, exactly as for Hermitian
operators (cf.\ \Eref, \S II), of a \emph{proper extension}.  A
conjugate pair without a proper extension is \emph{maximal}.

We also note that every extension of a Hermitian operator is Hermitian:
if \(R,R^*\) has the extension \(S,S^*\), then
\[
  R=R^*\quad\Longrightarrow\quad S=S^*.
\]
Indeed, if \(Sf\) and \(Rg\) are defined, then
\[
  (Sf,g)=(f,S^*g)=(f,R^*g)=(f,Rg)
  =(f,Sg)=(S^*f,g).
\]
Since these \(g\)'s span the closed linear manifold \(\mathfrak H\), it
follows that \(Sf=S^*f\), and hence \(S=S^*\).

It is quite easy to see (exactly as for Hermitian operators; cf.\
\Eref, \S II, Theorem~9) that every conjugate pair can be extended
to one consisting of two linear operators.  It is not, however,
immediately clear whether they can also be made closed in this
way.\footnote{\label{fn:61}The methods in question yield only a kind of common
closure of the two operators \(R,R^*\) of the pair: from
\(f_n\to f\), \(Rf_n\to f^*\), and \(R^*f_n\to f^{**}\), one infers
that \(Rf,R^*f\) are defined and equal to \(f^*,f^{**}\),
respectively\TranslatorNoteMark.}
\TranslatorNoteText{At the end of \NoteRef{61} the source identifies
\(Rf,R^*f\) with \(f^{**},f^{***}\).  The convergence assumptions
immediately preceding instead give \(f^*,f^{**}\), as printed here.}
We shall not pursue this question here.

\subsection*{3.}\phantomsection\label{sec:iv-3}
An \(A\in\mathfrak B\) is normal if \(A,A^*\) commute, that is, if
\((A)\) is Abelian, or, equivalently
(cf.\ \hyperref[sec:ring-and-commutant-1]{\S\,II, 1}), if
\((A)''\) is Abelian.  We transfer this definition to arbitrary \(R\)
(since always \((R)''\subset\mathfrak B\)).

\begin{definition}[Definition 6]\label{def:normal-conjugate-pair}
A conjugate pair \(R,R^*\) is \emph{normal} if \((R)''\)
(a subset of \(\mathfrak B\)) is Abelian.
\end{definition}

By what was said above, it is clear that for operators in
\(\mathfrak B\) this is the old definition.

\begin{theorem}[Theorem 11]\label{thm:11}
The pair \(R,R^*\) is normal if and only if \(R^*,R\) is normal.
\end{theorem}

\begin{proof}
We prove
\[
  (R)'=(R^*)'.
\]
It then follows that \((R)''=(R^*)''\), and hence the assertion.  By
symmetry, it is enough

% [Source p. 407 / PDF 38]
to prove \((R)'\subset(R^*)'\); that is, if \(A,A^*\) commute with
\(R\), they must also commute with \(R^*\).

Suppose, then, that \(R^*f\) is defined.  Then \(Rf,RAf,RA^*f\) are
defined as well, and hence so are \(R^*Af,R^*A^*f\).  If \(R^*g\) is
also defined, we have
\[
\begin{aligned}
  (R^*Af,g)
  &=(Af,Rg)=(f,A^*Rg)=(f,RA^*g)\\
  &=(R^*f,A^*g)=(AR^*f,g),
\end{aligned}
\]
\[
\begin{aligned}
  (R^*A^*f,g)
  &=(A^*f,Rg)=(f,ARg)=(f,RAg)\\
  &=(R^*f,Ag)=(A^*R^*f,g).
\end{aligned}
\]
Since these \(g\)'s span the closed linear manifold \(\mathfrak H\),
\[
  R^*Af=AR^*f,\qquad
  R^*A^*f=A^*R^*f.
\]
Thus \(A,A^*\) commute with \(R^*\), as asserted.
\end{proof}

\begin{theorem}[Theorem 12]\label{thm:12}
If \(R,R\) (\(R\) a Hermitian operator) is assumed linear and closed,
then it is normal if and only if \(R\) is hypermaximal.\footnote{\label{fn:62}
``Hypermaximal'' is abbreviated \emph{hypermax.}; cf.\ \Eref,
Definition~9.  If \(R\) is not linear and closed, consider its closure
\(\widetilde R\) (\Eref, \S II, Theorem~9).  Since it is Hermitian,
\(\widetilde R,\widetilde R\) is a conjugate pair; and since everything
commuting with \(R\) also commutes with \(\widetilde R\), one has
\((R)'\subset(\widetilde R)'\) and
\((R)''\supset(\widetilde R)''\).  Thus \((\widetilde R)''\) is Abelian
as well.  Hence \(\widetilde R,\widetilde R\) is normal, and the present
theorem applies to the closed linear operator \(\widetilde R\).}
\end{theorem}

\begin{proof}
We begin with some general observations about closed linear Hermitian
operators.  Form the Cayley transform \(U\) of \(R\).  Let
\(\mathfrak E,\mathfrak F\) be its domain and range, respectively
(cf.\ \Eref, \S V, Theorem~2); both are closed linear manifolds.
Let \(E\) be the projection operator onto \(\mathfrak E\)
(cf.\ \Eref, \S III, Theorem~13).

Suppose \(A\) commutes with \(R\).  Every \(\varphi\) for which
\(U\varphi\) is defined has the form \(Rf+if\); therefore \(RAf\) is
defined, and
\[
  A\varphi=ARf+iAf=RAf+iAf.
\]
Thus \(U(A\varphi)\) is also defined, and
\[
  U(A\varphi)=RAf-iAf=ARf-iAf=A(U\varphi);
\]
that is, \(U,A\) commute.  Conversely, suppose \(A\) commutes with
\(U\).  Every \(f\) for which \(Rf\) is defined has the form
\(\varphi-U\varphi\); hence \(UA\varphi\) is defined and
\[
  Af=A\varphi-AU\varphi=A\varphi-UA\varphi.
\]
Thus \(R(Af)\) is defined as well, and
\[
  R(Af)=i(A\varphi+UA\varphi)
  =i(A\varphi+AU\varphi)=A(Rf);
\]
that is, \(R,A\) commute.  It follows that
\[
  (R)'=(U)',\qquad (R)''=(U)''.
\]

The sufficiency of the condition is now immediate.  If \(R\) is
hypermaximal, then \(U\) is unitary (cf.\ \Eref, Theorem~35), and
therefore belongs to \(\mathfrak B\).  Moreover, it commutes with
\(U^*=U^{-1}\); hence \(U,U^*\) is normal, and, because
\((R)''=(U)''\), so is \(R,R\).

It remains to prove necessity.  Let \(R,R\) be normal, and let
\(A\in(R)'\).  Then \(A\in(U)'\).  Along with \(Uf\), the vector \(UAf\)
is defined; that is, together with \(f\), \(Af\) belongs to
\(\mathfrak E\).  The vector \(Ef\) always belongs to \(\mathfrak E\),
and hence so does \(AEf\).  Therefore
\[
  EAEf=AEf,\qquad\text{or}\qquad EAE=AE.
\]
Taking adjoints and replacing \(A\) by \(A^*\) (which also belongs to
\((R)'\)) gives \(EAE=EA\).  Hence \(AE=EA\): \(A,E\) commute.  Thus
\(A\) commutes with both \(U\) and \(E\), and therefore with
\[
  UE=V.
\]
This is important because \(V\) is everywhere defined and hence belongs
to \(\mathfrak B\).  Since this holds for every \(A\in(R)'\) (which
contains \(A^*\) together with \(A\)), it follows that \(V\in(R)''\).

% [Source p. 408 / PDF 39]
Thus \(V^*\in(R)''\) as well; and, since \((R)''\) is Abelian,
\(V,V^*\) commute.  For all \(f,g\),
\[
  (V^*Vf,g)=(Vf,Vg)=(U(Ef),U(Eg))=(Ef,Eg)=(Ef,g).
\]
Consequently
\[
  V^*V=VV^*=E,
\]
and this operator too must belong to \((R)''\).  Hence \(V,E\) commute,
with the consequence
\[
  EV=VE=UE\cdot E=UE=V.
\]
If \(\varphi\in\mathfrak E\), then
\[
  \varphi=E\varphi,\qquad
  U\varphi=UE\varphi=V\varphi=EV\varphi.
\]
Thus \(U\varphi\in\mathfrak E\), and therefore
\(\varphi-U\varphi\in\mathfrak E\).  But the vectors
\(\varphi-U\varphi\) must be dense everywhere (\Eref, Theorem~24);
hence \(\mathfrak E\) is an everywhere dense closed linear manifold.
Thus
\[
  \mathfrak E=\mathfrak H,\qquad E=1.
\]
It follows that
\[
  V=U\cdot1=U,\qquad V^*V=VV^*=1,
\]
and hence
\[
  U^*U=UU^*=1.
\]
Thus \(U\) is unitary, and \(R\) is hypermaximal.
\end{proof}

We emphasize once more the fact found incidentally in the proof of this
theorem: for a hypermaximal Hermitian operator \(R\) and its unitary
Cayley transform \(U\),
\[
  (R)'=(U)',
\]
and consequently
\[
  (R)''=(U)'',\quad (R)'''=(U)''',\ldots.
\]
Now \((U)''\) is Abelian, so
\[
  (U)''\subset(U)'''=(U)',\qquad
  (R)''\subset(R)'.
\]
Moreover,
\[
  (R)'=(R)'''=\cdots,\qquad
  (R)''=(R)^{\mathrm{IV}}=\cdots,
\]
because this is so for \(U\), or by \hyperref[sec:iv-1]{\S\,IV, 1}.

\begin{theorem}[Theorem 13]\label{thm:13}
Let \(R\), as above, be a hypermaximal Hermitian operator, \(U\) its
Cayley transform, \(E(\lambda)\) its resolution of the identity, and
\(\mathcal E\) the set of all \(E(\lambda)\).  Then
\[
  (R)'=(U)'=\mathcal E',
\]
and hence
\[
  (R)''=(U)''=\mathcal E'',\qquad
  (R)'''=(U)'''=\mathcal E''',\ldots.
\]
\end{theorem}

\begin{proof}
Everything follows from the first equality; and since
\((R)'=(U)'\) has already been established, only
\((U)'=\mathcal E'\) remains to be proved.  This follows, in any case,
from\footnote{\label{fn:63}In general,
\[
  \mathfrak M\subset R(\mathfrak M)\subset\mathfrak M'',\qquad
  \mathfrak M'\supset(R(\mathfrak M))'\supset\mathfrak M'''=\mathfrak M';
\]
hence \((R(\mathfrak M))'=\mathfrak M'\): the ring \(R(\mathfrak M)\)
determines \(\mathfrak M'\).}
\[
  R(U)=R(\mathcal E).
\]
It is therefore enough to show that \(U\in R(\mathcal E)\), and that
every \(E(x)\)\footnote{\label{fn:64}We write \(E(x)\), \(0<x<1\), in place of
\(E(\lambda)\), \(-\infty<\lambda<\infty\), with the correspondence
\(\lambda=-\cot\pi x\),
\(x=-\frac1\pi\operatorname{arccot}\lambda\)
(cf.\ \Eref, proof of Theorem~36).} belongs to \(R(U)\).

From the relation
\[
  (Uf,g)=\int_0^1e^{2\pi ix}\,d(E(x)f,g),
\]
one proves literally as in the corresponding part of the first half of
the proof of \TheoremRef{1}, for the bounded Hermitian operator \(A\) used
there, that \(U\in R(\mathcal E)\).  To see, conversely, that
\(E(x)\in R(U)\), recall how these \(E(x)\) were constructed in
\Eref, Appendix~II.  The \(E(x)\) were differences

% [Source p. 409 / PDF 40]
of strong limits of certain \(F(\rho)\), \(0<\rho<1\)
(cf.\ \S6 there); it is therefore enough that these belong to \(R(U)\).
The \(F(\rho)\), however, were formed by addition, subtraction, and
multiplication from the \(E(U,a,b)\),
\(-\infty<a,b<\infty\) (cf.\ \S\S4 and 5 there), so only the latter
matter.  But \(E(U,a,b)\) was a product of projection operators from the
resolutions of the identity\footnote{\label{fn:65}\(E(U,a,b)\) (called \(U_A\) in the
passage cited) is
defined as
\[
  P_{\mathfrak M}(R,a)\,P_{\mathfrak N}(S,b),
  \qquad
  R=\frac{U+U^*}{2},\quad S=\frac{U-U^*}{2i},
\]
and only the properties of these projection operators stated in
Theorems \(8^*,9^*\) there are used; these properties plainly belong to
the corresponding resolutions of the identity.} of
\[
  \frac{U+U^*}{2}
  \qquad\text{and}\qquad
  \frac{U-U^*}{2i};
\]
and, by our \TheoremRef{2}, these belong to \(R(U)\) provided the last two
operators do.  This is clear: \(R(U)\) contains \(U,U^*\), and hence
also the two operators just displayed.
\end{proof}

\subsection*{4.}\phantomsection\label{sec:iv-4}
After these preparations we return to our proper subject.  Let
\(R,R^*\) again be a conjugate pair.

\begin{theorem}[Theorem 14]\label{thm:14}
Let \(R,R^*\) be normal.  Form the two Hermitian operators
\[
  S_1=\frac{R+R^*}{2},\qquad
  S_2=\frac{R-R^*}{2i},
\]
and then \(\widetilde S_1,\widetilde S_2\) (cf.\ the beginning of
\NoteRef{62}); the latter are hypermaximal.  The operators
\[
  \overline R=\widetilde S_1+i\widetilde S_2,\qquad
  \overline R^{\,*}=\widetilde S_1-i\widetilde S_2
\]
(both defined on the intersection of the domains of
\(\widetilde S_1,\widetilde S_2\)) form a conjugate pair, and in fact a
normal one.  The pair \(\overline R,\overline R^{\,*}\) is an extension
of \(R,R^*\).
\end{theorem}

\begin{proof}
Form \(S_1,S_2\) as prescribed; their Hermitian character is evident;
then form \(\widetilde S_1,\widetilde S_2\).  Let
\(A\in\mathfrak B\).  If \(A\) commutes with \(R,R^*\), it also commutes
with \(S_1,S_2\), and hence with
\(\widetilde S_1,\widetilde S_2\).  It follows that
\[
  (R,R^*)'\subset(\widetilde S_1)',\qquad
  (R,R^*)'\subset(\widetilde S_2)'.
\]
By \((R)'=(R^*)'\) (cf.\ the proof of \TheoremRef{11}),
\((R)'=(R,R^*)'\); hence
\[
  (R)'\subset(\widetilde S_1)',\qquad
  (R)'\subset(\widetilde S_2)',\qquad
  (R)''\supset(\widetilde S_1)'',\qquad
  (R)''\supset(\widetilde S_2)''.
\]
Along with \((R)''\), therefore,
\((\widetilde S_1)''\) and \((\widetilde S_2)''\) are Abelian.  Thus the
conjugate pairs
\(\widetilde S_1,\widetilde S_1\) and
\(\widetilde S_2,\widetilde S_2\) are normal, and by
\TheoremRef{12} the operators
\(\widetilde S_1,\widetilde S_2\) are hypermaximal.

Since \(\widetilde S_1,\widetilde S_2\) are extensions of \(S_1,S_2\),
respectively, the pair \(\overline R,\overline R^{\,*}\) is an extension
of \(R,R^*\).  It is clear, since
\(\widetilde S_1,\widetilde S_2\) are Hermitian, that
\(\overline R,\overline R^{\,*}\) is a conjugate pair.  Its normality is
seen as follows.  We have
\[
  (\overline R)'=(\overline R^{\,*})',
\]
both being equal to
\((\overline R,\overline R^{\,*})'\); and this set plainly contains
\((\widetilde S_1,\widetilde S_2)'\), which in turn, as we already know,
contains \((R)'\).  Thus
\[
  (\overline R)'\supset(R)',\qquad
  (\overline R)''\subset(R)''.
\]
Since \((R)''\) is Abelian, so is \((\overline R)''\); hence
\(\overline R,\overline R^{\,*}\) is normal.
\end{proof}

\begin{theorem}[Theorem 15]\label{thm:15}
The vectors \(\overline Rf\) and \(\overline R^{\,*}f\) are defined if
and only if there exist two vectors \(f^*,f^{**}\) such that, for every
\(g\) in the common domain of \(R\) and \(R^*\),
\[
  (f,Rg)=(f^{**},g),\qquad
  (f,R^*g)=(f^*,g).
\]
In that case\TranslatorNote{In the second formula of the conclusion the
source prints \(\overline R^{\,*}f=f^{***}\).  The defining identity
immediately above requires \(f^{**}\), as printed here.}
\[
  \overline Rf=f^*,\qquad
  \overline R^{\,*}f=f^{**}.
\]
\end{theorem}

% [Source p. 410 / PDF 41]
\begin{proof}
Necessity is clear:
\[
  (f,Rg)=(f,\overline Rg)=(\overline R^{\,*}f,g),
  \qquad
  (f,R^*g)=(f,\overline R^{\,*}g)=(\overline Rf,g).
\]
For sufficiency, observe that
\[
  \frac{f^*+f^{**}}2,\qquad
  \frac{f^*-f^{**}}{2i}
\]
are, by definition, the vectors assigned to \(f\) by \(S_1,S_2\),
respectively (cf.\ \Eref, Definitions~8 and 9), and hence also by
\(\widetilde S_1,\widetilde S_2\).  Since the latter are hypermaximal,
\(\widetilde S_1f,\widetilde S_2f\) are defined and equal,
respectively, to
\[
  \frac{f^*+f^{**}}2,\qquad
  \frac{f^*-f^{**}}{2i}.
\]
Consequently \(\overline Rf,\overline R^{\,*}f\) are defined as well,
and are equal to \(f^*,f^{**}\), respectively.
\end{proof}

\begin{theorem}[Theorem 16]\label{thm:16}
Every extension \(T,T^*\) of \(R,R^*\) (a conjugate pair, not
necessarily a normal one) has \(\overline R,\overline R^{\,*}\) as an
extension.  The pair \(\overline R,\overline R^{\,*}\), even regarded
merely as a conjugate pair, is maximal; indeed, it is the unique maximal
extension of \(R,R^*\).
\end{theorem}

\begin{proof}
The last two assertions follow from the first, so we need consider only
that one.  Suppose \(Tf,T^*f\) are defined, and put
\[
  Tf=f^*,\qquad T^*f=f^{**}.
\]
Then the hypotheses of \TheoremRef{15} are satisfied.  Indeed, when
\(Rg,R^*g\) are defined,
\[
  (f,Rg)=(f,Tg)=(T^*f,g),\qquad
  (f,R^*g)=(f,T^*g)=(Tf,g).
\]
Thus \(\overline Rf,\overline R^{\,*}f\) are also defined, and equal to
\(Tf,T^*f\), respectively.  Hence
\(\overline R,\overline R^{\,*}\) is an extension of \(T,T^*\).
\end{proof}

For normal conjugate pairs of operators, therefore, the situation is
much simpler than for Hermitian operators (cf.\ \Eref): a maximal
extension is always possible here in one and only one way.  In what
follows we shall consequently restrict ourselves to maximal pairs
\(R,R^*\); by \TheoremRef{16} these are characterized by
\[
  \overline R=R,\qquad \overline R^{\,*}=R^*.
\]
For such pairs we shall establish the analogue of Hilbert's spectral
representation.

\section{Spectral Form of Normal Operators}
\label{sec:v-spectral-form}

\subsection*{1.}\phantomsection\label{sec:v-1}
As announced, we investigate normal maximal conjugate pairs \(R,R^*\).

\begin{theorem}[Theorem 17]\label{thm:17}
Let \(R,R^*\) be a normal maximal conjugate pair.  There then exists a
family of projection operators \(\Delta(z)\), indexed by all complex
numbers \(z\), with the following properties.

\(\alpha\))
\[
  \Delta(z')\Delta(z'')=\Delta(z'')\Delta(z')=\Delta(z),
\]
where
\[
  \operatorname{Re}z=\min(\operatorname{Re}z',
  \operatorname{Re}z''),\qquad
  \operatorname{Im}z=\min(\operatorname{Im}z',
  \operatorname{Im}z'').
\]

\(\beta\)) For every \(f\),
\[
  \Delta(z)f\longrightarrow\Delta(z_0)f
\]
as \(z\to z_0\), while
\(\operatorname{Re}z\geq\operatorname{Re}z_0\) and
\(\operatorname{Im}z\geq\operatorname{Im}z_0\).  The convergence is
uniform both when \(f,\operatorname{Re}z\) are fixed and
\(\operatorname{Im}z\) is arbitrary, and when
\(f,\operatorname{Im}z\) are fixed and \(\operatorname{Re}z\) is
arbitrary.

% [Source p. 411 / PDF 42]
\(\gamma\)) For every \(f\),
\[
  \Delta(z)f\longrightarrow0
\]
when \(\operatorname{Re}z\to-\infty\) or
\(\operatorname{Im}z\to-\infty\), whereas
\[
  \Delta(z)f\longrightarrow f
\]
when \(\operatorname{Re}z\to+\infty\) and
\(\operatorname{Im}z\to+\infty\).

(A family of projection operators satisfying \(\alpha\))--\(\gamma\))
is called a \emph{complex resolution of the identity}.)

\(\delta\)) The vector \(Rf\) (and \(R^*f\)) is defined if and only if
\[
  \iint_{\mathbb C}|z|^2\,d\lvert\Delta(z)f\rvert^2
\]
is finite.  (The double integral extends over the whole complex plane.
For every rectangle \(\mathfrak R\)\footnotemark,
\[
  \iint_{\mathfrak R}d\lvert\Delta(z)f\rvert^2\geq0,
\]
\footnotetext{\label{fn:66}If the rectangle \(\mathfrak R\) has the vertices
\[
  a'+b'i,\quad a'+b''i,\quad a''+b''i,\quad a''+b'i
  \qquad(a'\leq a'',\ b'\leq b''),
\]
then
\[
\begin{aligned}
  \iint_{\mathfrak R}d\lvert\Delta(z)f\rvert^2
  ={}&\lvert\Delta(a'+b'i)f\rvert^2
     -\lvert\Delta(a'+b''i)f\rvert^2\\
    &+\lvert\Delta(a''+b''i)f\rvert^2
     -\lvert\Delta(a''+b'i)f\rvert^2.
\end{aligned}
\]
Since \(\lvert Ef\rvert^2=(f,Ef)\) for a projection operator \(E\),
this equals
\[
  \left|
  \left\{
  \Delta(a'+b'i)-\Delta(a'+b''i)
  +\Delta(a''+b''i)-\Delta(a''+b'i)
  \right\}f\right|^2\geq0,
\]
provided the operator in braces is a projection operator.  By
\(\alpha\)), it is indeed so, for
\[
\begin{aligned}
 &\Delta(a'+b'i)-\Delta(a'+b''i)
  +\Delta(a''+b''i)-\Delta(a''+b'i)\\
 &\qquad=
  \Delta(a''+b''i)
  \bigl(1-\Delta(a'+b''i)\bigr)
  \bigl(1-\Delta(a''+b'i)\bigr),
\end{aligned}
\]
which makes the assertion evident.}
and the integrand \(|z|^2\) is always nonnegative.  Thus the integral
displayed in \(\delta\)) is, by its nature, either convergent and
finite or properly divergent to \(+\infty\); the latter possibility is
to be excluded.)

\(\varepsilon\)) In the case just described, for every \(g\),
\[
  (Rf,g)=\iint_{\mathbb C}z\,d(\Delta(z)f,g),\qquad
  (R^*f,g)=\iint_{\mathbb C}\overline z\,d(\Delta(z)f,g).
\]
For such \(f\), and arbitrary \(g\), these integrals converge
absolutely.\footnote{\label{fn:67}All these arguments are closely analogous to those
of \Eref, Theorem~36, with which they may be compared.  For
completeness we give here a proof of the absolute convergence of
\(\iint z\,d(\Delta(z)f,g)\); the proof for
\(\iint\overline z\,d(\Delta(z)f,g)\) is the same.

Let \(\mathfrak R_1,\ldots,\mathfrak R_k\) be any rectangles with no
common interior points, together filling a finite part of the plane,
and let \(\zeta_1,\ldots,\zeta_k\) be arbitrary points in them.  If
\(E_h\) is the projection operator calculated for the rectangle
\(\mathfrak R_h\) as in \NoteRef{66}, then, for \(h=1,\ldots,k\),
\[
\begin{aligned}
 \left|\iint_{\mathfrak R_h}d(\Delta(z)f,g)\right|
 &=|(E_hf,g)|=|(E_hf,E_hg)|\\
 &\leq|E_hf|\,|E_hg|\\
 &=\sqrt{
   \iint_{\mathfrak R_h}d|\Delta(z)f|^2\,
   \iint_{\mathfrak R_h}d|\Delta(z)g|^2}.
\end{aligned}
\]

% [Source p. 412 / PDF 43]
It follows (cf.\ the estimates in \Eref, proof of Theorem~36, and
our \NoteRef{45}) that
\[
\begin{aligned}
 \left|\sum_{h=1}^k\zeta_h
   \iint_{\mathfrak R_h}d(\Delta(z)f,g)\right|
 &\leq
 \sum_{h=1}^k|\zeta_h|
 \sqrt{
   \iint_{\mathfrak R_h}d|\Delta(z)f|^2\,
   \iint_{\mathfrak R_h}d|\Delta(z)g|^2}\\
 &\leq
 \sqrt{\sum_{h=1}^k|\zeta_h|^2
   \iint_{\mathfrak R_h}d|\Delta(z)f|^2}\,
 \sqrt{\sum_{h=1}^k
   \iint_{\mathfrak R_h}d|\Delta(z)g|^2}.
\end{aligned}
\]
But
\[
  \iint_{\mathbb C}|z|^2\,d|\Delta(z)f|^2
  \quad\text{and}\quad
  \iint_{\mathbb C}d|\Delta(z)g|^2=|g|^2
\]
are finite and hence converge absolutely (cf.\ the observation following
\(\delta\))).  The same therefore follows for
\(\iint z\,d(\Delta(z)f,g)\).}

(On the basis of \(\delta\))--\(\varepsilon\)), \(R,R^*\) and
\(\Delta(z)\) are said to \emph{belong together}.)

\(\zeta\)) If the Hermitian operators
\(\widetilde S_1,\widetilde S_2\) are formed as in
\TheoremRef{14}, then \(\widetilde S_1f\), respectively
\(\widetilde S_2f\), is defined if and only if
\[
  \iint_{\mathbb C}(\operatorname{Re}z)^2\,
  d|\Delta(z)f|^2,
  \qquad\text{respectively}\qquad
  \iint_{\mathbb C}(\operatorname{Im}z)^2\,
  d|\Delta(z)f|^2
\]
is finite.  (For the nature of these integrals see
\hyperref[thm:17]{p.~411} and
\NoteRef{66}.)

\(\eta\)) In that case, for every \(g\),
\[
  (\widetilde S_1f,g)
  =\iint_{\mathbb C}\operatorname{Re}z\,
    d(\Delta(z)f,g),
\]
\[
  (\widetilde S_2f,g)
  =\iint_{\mathbb C}\operatorname{Im}z\,
    d(\Delta(z)f,g).
\]
The absolute convergence of these integrals is assured; cf.\
\(\varepsilon\)) and \NoteRef{67}.
\end{theorem}

\begin{proof}
Since \(\widetilde S_1,\widetilde S_2\) are hypermaximal, there belong
to them two resolutions of the identity \(E(\lambda),F(\lambda)\)
(cf.\ \Eref, Theorem~36, and our \NoteRef{42}).  By definition,
\(\widetilde S_1f\), respectively \(\widetilde S_2f\), is defined if
\[
  \int_{-\infty}^{\infty}\lambda^2\,d|E(\lambda)f|^2,
  \qquad\text{respectively}\qquad
  \int_{-\infty}^{\infty}\lambda^2\,d|F(\lambda)f|^2,
\]
is finite; and then, for every \(g\),
\[
  (\widetilde S_1f,g)
  =\int_{-\infty}^{\infty}\lambda\,d(E(\lambda)f,g),
\]
\[
  (\widetilde S_2f,g)
  =\int_{-\infty}^{\infty}\lambda\,d(F(\lambda)f,g).
\]

Let the set of all \(E(\lambda)\), respectively all \(F(\lambda)\), be
\(\mathcal E\), respectively \(\mathcal F\).  Every \(E(\lambda)\)
belongs to \(\mathcal E\), hence to \(\mathcal E''\), and therefore by
\TheoremRef{13} to \((\widetilde S_1)''\).  But already in the
proof of \TheoremRef{14} we saw that
\((\widetilde S_1)''\subset(R)''\); hence
\(\mathcal E\subset(R)''\).  Similarly,
\(\mathcal F\subset(R)''\).  Since \((R)''\) is Abelian, the elements of
\(\mathcal E\) commute with those of \(\mathcal F\): every \(E(a)\)
commutes with every \(F(b)\).  We may therefore define a family of
projection operators \(\Delta(z)\) by
\[
  \Delta(a+ib)=E(a)F(b)=F(b)E(a).
\]
It is readily verified that \(\alpha\))--\(\gamma\)) hold; that is,
this is a complex resolution of the identity.

Furthermore, the integrals
\[
  \int_{-\infty}^{\infty}\lambda^2\,d|E(\lambda)f|^2,\qquad
  \int_{-\infty}^{\infty}\lambda^2\,d|F(\lambda)f|^2,
\]
\[
  \int_{-\infty}^{\infty}\lambda\,d(E(\lambda)f,g),\qquad
  \int_{-\infty}^{\infty}\lambda\,d(F(\lambda)f,g)
\]

% [Source p. 413 / PDF 44]
pass respectively into
\[
  \int_{-\infty}^{\infty}\int_{-\infty}^{\infty}
  a^2\,d|\Delta(a+ib)f|^2,
\qquad
  \int_{-\infty}^{\infty}\int_{-\infty}^{\infty}
  b^2\,d|\Delta(a+ib)f|^2,
\]
\[
  \int_{-\infty}^{\infty}\int_{-\infty}^{\infty}
  a\,d(\Delta(a+ib)f,g),
\qquad
  \int_{-\infty}^{\infty}\int_{-\infty}^{\infty}
  b\,d(\Delta(a+ib)f,g).
\]
Thus, on putting \(z=a+ib\), they become precisely the integrals in
\(\zeta\))--\(\eta\)).  Those assertions are therefore established.

Since \(R=\overline R\) and
\(R^*=\overline R^{\,*}\), the vector \(Rf\) (and \(R^*f\)) is defined
if and only if \(\widetilde S_1f,\widetilde S_2f\) are both defined;
that is, if and only if both integrals in \(\zeta\)) are finite.  In
view of what was said there about their essentially nonnegative nature,
this means that their sum is finite.  Since
\[
  |z|^2=(\operatorname{Re}z)^2+(\operatorname{Im}z)^2,
\]
this is exactly assertion \(\delta\)).  Assertion \(\varepsilon\))
follows immediately from \(\eta\)).
\end{proof}

\begin{theorem}[Theorem 18]\label{thm:18}
The family \(\Delta(z)\) is already determined uniquely by conditions
\(\alpha\))--\(\varepsilon\)).  Thus, to every maximal conjugate pair
\(R,R^*\) there belongs one and only one complex resolution of the
identity \(\Delta(z)\).
\end{theorem}

\begin{proof}
Let \(\Delta(z)\) be the family constructed in the proof of
\TheoremRef{17}, and let \(\Delta'(z)\) be another family satisfying
\(\alpha\))--\(\varepsilon\)).  We must prove
\[
  \Delta(z)=\Delta'(z)
\]
for every \(z\).

First, \(\eta\)) follows from \(\delta\)), at least for those \(f\) for
which \(Rf\) is defined, that is, for which both
\(\widetilde S_1f,\widetilde S_2f\) are defined.  For these \(f\), it
therefore holds with \(\Delta'(z)\) as well.  But
\(\widetilde S_1\), respectively \(\widetilde S_2\), restricted to the
domain of \(R\), is \(S_1\), respectively \(S_2\) (both have the same
domain).  We are therefore dealing with the \(f\)'s for which
\(S_1f,S_2f\) are defined.

As \(b\) increases, with \(a\) fixed, \(\Delta'(a+ib)\) never decreases
in the ordering of projection operators (cf.\ \Eref, \S III,
Theorem~14).  Hence, as \(b\to\infty\), a strong limit exists
(\Eref, \S III, Theorem~19):
\[
  \Delta'(a+ib)\longrightarrow E'(a).
\]
Similarly, with \(b\) fixed, a strong limit exists as \(a\to\infty\):
\[
  \Delta'(a+ib)\longrightarrow F'(b).
\]
The operators \(E'(a),F'(b)\) are projection operators which, like
\(\Delta'(a+ib)\), never decrease as \(a,b\) increase.  For
\(a'\geq a,\ b'\geq b\), condition \(\alpha\)) gives
\[
  \Delta'(a'+ib)\Delta'(a+ib')
  =\Delta'(a+ib')\Delta'(a'+ib)
  =\Delta'(a+ib).
\]
On passing to the limit \(a',b'\to+\infty\), this yields
\[
  F'(b)E'(a)=E'(a)F'(b)=\Delta'(a+ib).
\]
From
\[
 \Delta'(a+ib)f\longrightarrow
 \begin{cases}
   E'(a)f,&b\to+\infty,\\
   0,&b\to-\infty,
 \end{cases}
\qquad
 \Delta'(a+ib)f\longrightarrow
 \begin{cases}
   F'(b)f,&a\to+\infty,\\
   0,&a\to-\infty,
 \end{cases}
\]

% [Source p. 414 / PDF 45]
one immediately obtains
\[
\begin{aligned}
 \iint_{\mathbb C}\operatorname{Re}z\,
 d(\Delta'(z)f,g)
 &=\int_{-\infty}^{\infty}\int_{-\infty}^{\infty}
 a\,d(E'(a)F'(b)f,g)\\
 &=\int_{-\infty}^{\infty}a\,d(E'(a)f,g),
\end{aligned}
\]
\[
\begin{aligned}
 \iint_{\mathbb C}\operatorname{Im}z\,
 d(\Delta'(z)f,g)
 &=\int_{-\infty}^{\infty}\int_{-\infty}^{\infty}
 b\,d(E'(a)F'(b)f,g)\\
 &=\int_{-\infty}^{\infty}b\,d(F'(b)f,g),
\end{aligned}
\]
in the sense that each right-hand side converges whenever the
corresponding left-hand side does.  The two equalities hold
independently of one another.

It follows without difficulty from \(\alpha\))--\(\gamma\)) that
\(E'(\lambda),F'(\lambda)\) are resolutions of the identity.  Let the
Hermitian operators belonging to them be \(T_1,T_2\), respectively
(cf.\ \Eref, Theorem~36).  By what has already been shown, whenever
\(T_1f\) and \(S_1g\) are both defined,
\[
  (f,S_1g)
  =\int_{-\infty}^{\infty}\lambda\,d(f,E'(\lambda)g)
  =\int_{-\infty}^{\infty}\lambda\,d(E'(\lambda)f,g)
  =(T_1f,g).
\]
Thus \(f\) is an extension element of \(S_1\), with assigned value
\(T_1f\), and hence also an extension element of \(\widetilde S_1\).  But
\(\widetilde S_1\) is hypermaximal; consequently
\(\widetilde S_1f\) is defined and equals \(T_1f\).  Thus
\(\widetilde S_1\) is an extension of \(T_1\); and, since \(T_1\) too is
hypermaximal,
\[
  \widetilde S_1=T_1.
\]
Likewise one proves
\[
  \widetilde S_2=T_2.
\]
Since the resolutions of the identity of
\(\widetilde S_1,\widetilde S_2\) are \(E(\lambda),F(\lambda)\),
respectively, it follows that
\[
  E'(\lambda)=E(\lambda),\qquad
  F'(\lambda)=F(\lambda).
\]
Therefore
\[
  \Delta'(a+ib)=E'(a)F'(b)=E(a)F(b)=\Delta(a+ib),
\]
that is, \(\Delta'(z)=\Delta(z)\), as was to be proved.
\end{proof}

\begin{theorem}[Theorem 19]\label{thm:19}
To every complex resolution of the identity, that is, to every family
\(\Delta(z)\) satisfying \(\alpha\))--\(\gamma\)), there belongs one and
only one conjugate pair \(R,R^*\), according to
\(\delta\))--\(\varepsilon\)); and this pair is maximal and normal.
\end{theorem}

\begin{proof}
If such a pair exists, it is unique: condition \(\delta\)) determines
its domain, and \(\varepsilon\)) determines
\((Rf,g),(R^*f,g)\), hence \(Rf,R^*f\) themselves.  As regards
existence, it remains to say this: if we have any pair of operators
\(R,R^*\) defined by \(\delta\)), \(\varepsilon\)) (and if the domain is
dense everywhere), then those two conditions themselves make them a
conjugate pair; only maximality and normality then remain to be checked.

The existence of \(R,R^*\) according to \(\delta\)),
\(\varepsilon\)) is established once we know the following.  If
\[
  \iint_{\mathbb C}|z|^2\,d|\Delta(z)f|^2
\]
is finite, then there exist two vectors \(f^*,f^{**}\) such that, for
every \(g\),
\[
  \iint_{\mathbb C}z\,d(\Delta(z)f,g)=(f^*,g),
  \qquad
  \iint_{\mathbb C}\overline z\,d(\Delta(z)f,g)=(f^{**},g).
\]
These are then \(Rf\) and \(R^*f\), respectively.  For the moment call
the two integrals on the left \(L^*(g),L^{**}(g)\), with \(f\) fixed.
Plainly,
\[
  L^*(a_1g_1+\cdots+a_ng_n)
  =\overline a_1L^*(g_1)+\cdots+\overline a_nL^*(g_n),
\]
\[
  L^{**}(a_1g_1+\cdots+a_ng_n)
  =\overline a_1L^{**}(g_1)+\cdots+\overline a_nL^{**}(g_n).
\]
% The proof continues on source p. 415 / PDF 46.
% [Source p. 415 / PDF 46]
By F.~Riesz's theorem (cf.\ \Eref, \NoteRef{52}) our goal is reached as soon as we find a constant \(c\) such that
\[
  \lvert L^{*}(g)\rvert\leq c\lvert g\rvert,
  \qquad
  \lvert L^{**}(g)\rvert\leq c\lvert g\rvert .
\]
But the estimate in \NoteRef{67} gives
\[
  \lvert L^{*}(g)\rvert
  \leq
  \sqrt{\iint \lvert z\rvert^{2}\,d\lvert\Delta(z)f\rvert^{2}}\,
  \sqrt{\iint d\lvert\Delta(z)g\rvert^{2}}
  =
  \sqrt{\iint \lvert z\rvert^{2}\,d\lvert\Delta(z)f\rvert^{2}}\,\lvert g\rvert ,
\]
and likewise
\[
  \lvert L^{**}(g)\rvert
  \leq
  \sqrt{\iint \lvert z\rvert^{2}\,d\lvert\Delta(z)f\rvert^{2}}\,\lvert g\rvert ,
\]
which proves everything required.

The domain, that is, the set of \(f\) for which
\(\iint \lvert z\rvert^{2}\,d\lvert\Delta(z)f\rvert^{2}\) is finite, is
everywhere dense.  To see this, proceed as follows.  Just as, in the proof
of \TheoremRef{18}, the resolutions of the identity \(E'(\lambda)\) and
\(F'(\lambda)\) were formed from \(\Delta'(z)\), we now form from
\(\Delta(z)\) two resolutions of the identity \(E(\lambda)\) and
\(F(\lambda)\).  The integral
\(\iint \lvert z\rvert^{2}\,d\lvert\Delta(z)f\rvert^{2}\) is finite whenever
\[
  \iint (\Re z)^{2}\,d\lvert\Delta(z)f\rvert^{2},
  \qquad
  \iint (\Im z)^{2}\,d\lvert\Delta(z)f\rvert^{2}
\]
are finite, that is, whenever
\[
  \int_{-\infty}^{\infty}\lambda^{2}\,d\lvert E(\lambda)f\rvert^{2},
  \qquad
  \int_{-\infty}^{\infty}\lambda^{2}\,d\lvert F(\lambda)f\rvert^{2}
\]
are finite (cf.\ the proof cited above).  This is certainly so for every
\(f\) for which
\[
  E(\lambda)f=F(\lambda)f=
  \begin{cases}
    f,& \lambda\geq c,\\
    0,& \lambda\leq -c
  \end{cases}
  \qquad
  \text{(\(c\) arbitrary but fixed)}
\]
holds,\footnote{\label{fn:68}For then the integrals from \(-\infty\) to \(\infty\) may
be replaced by integrals from \(-c\) to \(c\), and the integrand
\(\lambda^{2}\) is bounded.}
and hence for every
\[
  f=(E(c)-E(-c))(F(c)-F(-c))g .
\]
As this converges to \(g\) when \(c\to\infty\), while \(g\) is arbitrary,
the set in question is indeed everywhere dense.

It remains to show that \(R,R^{*}\) is maximal and normal.  First,
normality.  Let \(\mathcal D\) be the set of all \(\Delta(z)\).  Since these
operators are Hermitian and mutually commuting, \(\mathcal D\) is Abelian,
and consequently so is \(\mathcal D''\).  It therefore suffices to prove
\[
  (R)''\subset\mathcal D'',
\]
and a fortiori it suffices to prove \((R)'\supset\mathcal D'\).  This will
be established if we show that every \(A\) which commutes with all
\(\Delta(z)\) also commutes with \(R\).

Suppose, then, that \(Rf\) is defined.  We assert that \(R(Af)\) is
defined and equals \(A(Rf)\).  First we must infer the finiteness of
\[
  \iint\lvert z\rvert^{2}\,d\lvert\Delta(z)Af\rvert^{2}
\]
from that of
\[
  \iint\lvert z\rvert^{2}\,d\lvert\Delta(z)f\rvert^{2},
\]
and then prove, for example,
\[
  (R(Af),g)=(A(Rf),g)\qquad\text{for every \(g\)}.
\]
Equivalently, we must prove
\[
  (Af,R^{*}g)=(Rf,A^{*}g),
  \qquad\text{or}\qquad
  \iint z\,d(Af,\Delta(z)g)
  =
  \iint z\,d(\Delta(z)f,A^{*}g).
\]

The operator \(A\) is bounded, so that \(\lvert Af\rvert\leq c\lvert
f\rvert\) always.  Hence, for every rectangle \(\mathfrak R\), if
\(E_{\mathfrak R}\) denotes the projection operator constructed from it
as in \NoteRef{66},
\[
  \iint_{\mathfrak R}d\lvert\Delta(z)Af\rvert^{2}
  =
  \lvert E_{\mathfrak R}Af\rvert^{2}
  =
  \lvert AE_{\mathfrak R}f\rvert^{2}
  \leq
  c^{2}\lvert E_{\mathfrak R}f\rvert^{2}
  =
  c^{2}\iint_{\mathfrak R}d\lvert\Delta(z)f\rvert^{2}.
\]
Since \(\lvert z\rvert^{2}\geq0\), it follows
that\TranslatorNote{In this sentence and in both integrals immediately below,
the source prints \(\lambda^2\).  The integration is over the complex
\(z\)-plane and the preceding estimates require \(\lvert z\rvert^2\);
that expression is therefore used here.}
\[
  \iint\lvert z\rvert^{2}\,d\lvert\Delta(z)Af\rvert^{2}
  \leq
  c^{2}\iint\lvert z\rvert^{2}\,d\lvert\Delta(z)f\rvert^{2},
\]
which proves the first assertion.  The second is immediate from
\[
  (Af,\Delta(z)g)
  =
  (\Delta(z)Af,g)
  =
  (A\Delta(z)f,g)
  =
  (\Delta(z)f,A^{*}g).
\]

% [Source p. 416 / PDF 47]
We now consider maximality.  It states that
\[
  \overline R=R,\qquad \overline R^{*}=R^{*}.
\]
Since the intersection of the domains of \(\widetilde S_{1}\) and
\(\widetilde S_{2}\) (cf.\ \TheoremRef{14}) is the domain of
\(\overline R,\overline R^{*}\), we must require that \(R\), and hence
also \(R^{*}\), be defined everywhere on this intersection.  We again use
the resolutions of the identity \(E(\lambda)\) and \(F(\lambda)\) above,
and let the corresponding (hypermaximal) Hermitian operators be
\(T_{1},T_{2}\).  Evidently \(T_{1}\) is an extension of \(S_{1}\), hence,
being closed and linear, also of \(\widetilde S_{1}\); and since the latter
is hypermaximal, \(T_{1}=\widetilde S_{1}\).  In precisely the same way,
\(T_{2}=\widetilde S_{2}\).  Thus, in the intersection of the domains of
\(\widetilde S_{1}\) and \(\widetilde S_{2}\), both \(T_{1}f\) and
\(T_{2}f\) are defined; that is, the integrals
\[
  \int_{-\infty}^{\infty}\lambda^{2}\,
     d\lvert E(\lambda)f\rvert^{2}
  =
  \iint(\Re z)^{2}\,d\lvert\Delta(z)f\rvert^{2},
\]
\[
  \int_{-\infty}^{\infty}\lambda^{2}\,
     d\lvert F(\lambda)f\rvert^{2}
  =
  \iint(\Im z)^{2}\,d\lvert\Delta(z)f\rvert^{2}
\]
are finite.  Their sum
\(\iint\lvert z\rvert^{2}\,d\lvert\Delta(z)f\rvert^{2}\) is therefore
finite as well.  Hence \(Rf\) is defined, and the proof is complete.
\end{proof}

\TheoremRef{17}--\TheoremRef{19} show that maximal normal conjugate pairs of operators
correspond one-to-one to complex resolutions of the identity, just as
hypermaximal Hermitian operators correspond to real resolutions of the
identity (by \Eref, Theorem~36).  Because maximal extensions are unique in
general, and because there is no analogue here of the hypermaximality
condition, the situation is in fact somewhat simpler.

\subsection*{2. Some elementary properties of maximal normal conjugate pairs}
\phantomsection\label{sec:max-normal-properties}

\begin{theorem}[Theorem 20]\label{thm:20}
Let \(R,R^{*}\) be a maximal normal conjugate pair, let \(\Delta(z)\) be
its complex resolution of the identity, and let \(\mathcal D\) be the set
of all \(\Delta(z)\).  Then
\[
  (R)'=\mathcal D',\qquad
  (R)''=\mathcal D'',\qquad
  (R)'''=\mathcal D''',\quad\ldots .
\]
\end{theorem}

\begin{proof}
It plainly suffices to prove the first equality.  In the preceding proof
we already saw that \((R)'\supset\mathcal D'\); it remains only to show
\((R)'\subset\mathcal D'\).  In the proof of \TheoremRef{14} we obtained
\((R)'\subset(\widetilde S_{1})'\) and
\((R)'\subset(\widetilde S_{2})'\), hence
\[
  (R)'\subset(\widetilde S_{1},\widetilde S_{2})'.
\]
Thus it is enough to see that
\[
  (\widetilde S_{1},\widetilde S_{2})'\subset\mathcal D'.
\]
Let the resolutions of the identity of \(\widetilde S_{1}\) and
\(\widetilde S_{2}\) be \(E(\lambda)\) and \(F(\lambda)\), respectively,
and let \(\mathcal E\) and \(\mathcal F\) denote the sets of all
\(E(\lambda)\) and \(F(\lambda)\).  By \TheoremRef{13},
\[
  (\widetilde S_{1})'=\mathcal E',
  \qquad
  (\widetilde S_{2})'=\mathcal F',
\]
and therefore
\[
  (\widetilde S_{1},\widetilde S_{2})'=(\mathcal E,\mathcal F)'.
\]
If \(A\in\mathfrak B\) commutes with every element of both
\(\mathcal E\) and \(\mathcal F\), then it commutes with every
\[
  \Delta(a+ib)=E(a)F(b),
\]
and hence with every element of \(\mathcal D\).  Thus
\((\mathcal E,\mathcal F)'\subset\mathcal D'\), as required.
\end{proof}

\begin{theorem}[Theorem 21]\label{thm:21}
If \(R,R^{*}\) is a normal conjugate pair, then in general
\[
  \lvert Rf\rvert=\lvert R^{*}f\rvert.
\]
If \(R,R^{*}\) is also maximal, then both \(R\) and \(R^{*}\) are closed
linear operators.
\end{theorem}

\begin{proof}
Since \(R,R^{*}\) has at most one maximal extension (\TheoremRef{16}), it is
enough to prove the first assertion for maximal \(R,R^{*}\).  Let
\(\Delta(z)\) be its complex resolution of the identity.

% [Source p. 417 / PDF 48]
Whenever \(Rf\) (and \(R^{*}f\)) is defined, we have
\begin{align*}
  (Rf,\Delta(a+ib)g)
  &=
  \int_{-\infty}^{\infty}\!\!\int_{-\infty}^{\infty}
  (a'+ib')\,
  d\bigl(\Delta(a'+ib')f,\Delta(a+ib)g\bigr)\\
  &=
  \int_{-\infty}^{\infty}\!\!\int_{-\infty}^{\infty}
  (a'+ib')\,
  d\bigl(\Delta(a+ib)\Delta(a'+ib')f,g\bigr)\\
  &=
  \int_{-\infty}^{\infty}\!\!\int_{-\infty}^{\infty}
  (a'+ib')\,
  d\bigl(\Delta(\min(a,a')+i\min(b,b'))f,g\bigr)\\
  &=
  \int_{-\infty}^{a}\!\!\int_{-\infty}^{b}
  (a'+ib')\,d\bigl(\Delta(a'+ib')f,g\bigr).
\end{align*}
In particular,\TranslatorNote{In the last Stieltjes differentials above
the source prints the unprimed \(\Delta(a+ib)\).  The primed integration
variables shown here are required by the surrounding integrals.}
\begin{align*}
  (\Delta(a+ib)f,Rf)
  &=
  \int_{-\infty}^{a}\!\!\int_{-\infty}^{b}
  (a'-ib')\,d\bigl(f,\Delta(a'+ib')f\bigr)\\
  &=
  \int_{-\infty}^{a}\!\!\int_{-\infty}^{b}
  (a'-ib')\,d\lvert\Delta(a'+ib')f\rvert^{2}.
\end{align*}

It follows that\footnote{\label{fn:69}For this transformation of the integral,
cf.\ \Eref, \NoteRef{55}.}
\begin{align*}
  \lvert Rf\rvert^{2}
  &=(Rf,Rf)\\
  &=
  \int_{-\infty}^{\infty}\!\!\int_{-\infty}^{\infty}
  (a+ib)\,d\bigl(\Delta(a+ib)f,Rf\bigr)\\
  &=
  \int_{-\infty}^{\infty}\!\!\int_{-\infty}^{\infty}
  (a+ib)\,
  d\left[
    \int_{-\infty}^{a}\!\!\int_{-\infty}^{b}
    (a'-ib')\,d\lvert\Delta(a'+ib')f\rvert^{2}
  \right]\\
  &=
  \int_{-\infty}^{\infty}\!\!\int_{-\infty}^{\infty}
  (a+ib)(a-ib)\,d\lvert\Delta(a+ib)f\rvert^{2}\\
  &=
  \int_{-\infty}^{\infty}\!\!\int_{-\infty}^{\infty}
  (a^{2}+b^{2})\,d\lvert\Delta(a+ib)f\rvert^{2}.
\end{align*}
In the same way,
\[
  \lvert R^{*}f\rvert^{2}
  =
  \int_{-\infty}^{\infty}\!\!\int_{-\infty}^{\infty}
  (a^{2}+b^{2})\,d\lvert\Delta(a+ib)f\rvert^{2}.
\]
This proves \(\lvert Rf\rvert=\lvert R^{*}f\rvert\).

We turn to the second assertion.  The linearity of \(R\), as of \(R^{*}\),
follows, for example, from \TheoremRef{15} (maximality gives
\(R=\overline R\) and \(R^{*}=\overline R^{*}\)).  We now prove that
\(R\) is closed; the proof for \(R^{*}\) is identical.  Suppose that
\(f_n\to f\), that every \(Rf_n\), and hence every \(R^{*}f_n\), is
defined, and that \(Rf_n\to f^{*}\).  The \(Rf_n\) satisfy the Cauchy
condition,
\[
  \lvert R(f_m-f_n)\rvert\longrightarrow0
  \qquad(m,n\to\infty),
\]
and therefore so do the \(R^{*}f_n\), because
\[
  \lvert R(f_m-f_n)\rvert
  =
  \lvert R^{*}(f_m-f_n)\rvert .
\]
Consequently there exists an \(f^{**}\) such that
\(R^{*}f_n\to f^{**}\).  It follows at once that \(f,f^{*},f^{**}\)
satisfy the hypotheses of \TheoremRef{15}.  Thus \(Rf\) is defined and equals
\(f^{*}\).  Hence \(R\) is closed, as asserted.
\end{proof}

% [Source p. 418 / PDF 49]
\section*{Appendix I}
\label{app:bounded-spectral-families}
\addcontentsline{toc}{section}{Appendix I}

For completeness, we give here the familiar Hilbertian characterization
of the resolution of the identity of a bounded Hermitian operator.  Let
\(R\) be hypermaximal and let \(E(\lambda)\) be a resolution of the
identity.

By \Eref, \S\,II, Theorem~\(12\gamma\), the boundedness of \(R\) is
characterized by
\[
  \lvert(Rf,f)\rvert\leq c\lvert f\rvert^{2}
  \qquad\text{(\(c\) fixed),}
\]
and
\[
  (Rf,f)
  =
  \int_{-\infty}^{\infty}\lambda\,d(E(\lambda)f,f)
  =
  \int_{-\infty}^{\infty}\lambda\,d\lvert E(\lambda)f\rvert^{2}.
\]
If
\[
  E(\lambda)=
  \begin{cases}
    1,&\lambda\geq c,\\
    0,&\lambda\leq -c,
  \end{cases}
\]
then the integral from \(-\infty\) to \(\infty\) may be replaced by the
integral from \(-c\) to \(c\).  The integrand is then at most \(c\) in
absolute value, while \(\lvert E(\lambda)f\rvert^{2}\) never decreases;
hence the integral is, in absolute value, at most
\[
  c\int_{-\infty}^{\infty}d\lvert E(\lambda)f\rvert^{2}
  =
  c\lvert f\rvert^{2}.
\]
Thus \(R\) is bounded.

Conversely, suppose that \(E(\lambda)\neq0\) for arbitrarily small
\(\lambda\), respectively \(E(\lambda)\neq1\) for arbitrarily large
\(\lambda\); in other words, that \(E(\lambda)\) is not constant for
arbitrarily large \(\lvert\lambda\rvert\).  (Since
\(E(\lambda)f\to0\), respectively \(f\), as
\(\lambda\to-\infty\), respectively \(+\infty\), no constant value other
than \(0\), respectively \(1\), can occur.)  Then, for every \(D\), there
are \(\lambda,\mu\) such that
\[
  E(\lambda)\neq E(\mu),
  \qquad
  \lambda<\mu<-D
  \quad\text{or}\quad
  D<\lambda<\mu.
\]
Thus \(E(\mu)-E(\lambda)\) is a nonzero projection operator.  Let
\(f\neq0\) be a point of its closed linear manifold, so that
\[
  (E(\mu)-E(\lambda))f=f.
\]
Then\footnote{\label{fn:70}Indeed,
\[
  \lvert E(\mu)f\rvert^{2}-\lvert E(\lambda)f\rvert^{2}
  =
  \lvert(E(\mu)-E(\lambda))f\rvert^{2}
  =
  \lvert f\rvert^{2},
\]
while \(\lvert E(\mu)f\rvert^{2}\leq\lvert f\rvert^{2}\) and
\(\lvert E(\lambda)f\rvert^{2}\geq0\) (cf.\ \Eref, \S\,III).  Therefore
\(\lvert E(\mu)f\rvert=\lvert f\rvert\) and
\(\lvert E(\lambda)f\rvert=0\).  A fortiori, for
\(\lambda'\geq\mu\), respectively \(\lambda'\leq\lambda\), we have
\(\lvert E(\lambda')f\rvert=\lvert f\rvert\), respectively \(0\);
that is, \(E(\lambda')f=f\), respectively \(0\) (cf.\ the passage cited),
as asserted.}
\[
  E(\lambda')f=
  \begin{cases}
    f,&\lambda'\geq\mu,\\
    0,&\lambda'\leq\lambda.
  \end{cases}
\]
Consequently,
\begin{align*}
  \lvert(Rf,f)\rvert
  &=
  \left\lvert
  \int_{-\infty}^{\infty}
  \lambda'\,d\lvert E(\lambda')f\rvert^{2}
  \right\rvert\\
  &=
  \left\lvert
  \int_{\lambda}^{\mu}
  \lambda'\,d\lvert E(\lambda')f\rvert^{2}
  \right\rvert\\
  &\geq
  D\int_{\lambda}^{\mu}d\lvert E(\lambda')f\rvert^{2}\\
  &=
  D\int_{-\infty}^{\infty}d\lvert E(\lambda')f\rvert^{2}
  =
  D\lvert f\rvert^{2}.
\end{align*}
For \(D>c\), this is greater than \(c\lvert f\rvert^{2}\); hence \(R\) is
unbounded.

In precisely the same way one can read off the boundedness of a maximal
normal conjugate pair \(R,R^{*}\) from its complex resolution of the
identity \(\Delta(z)\).  It is defined by
\[
  \lvert Rf\rvert\leq c\lvert f\rvert
  \qquad\text{(\(c\) fixed),}
\]
and it is convenient to use the following relation (\TheoremRef{21}):
\[
  \lvert Rf\rvert^{2}
  =
  \iint\lvert z\rvert^{2}\,d\lvert\Delta(z)f\rvert^{2}.
\]

% [Source p. 419 / PDF 50]
\section*{Appendix II}
\label{app:laurent-matrices}
\addcontentsline{toc}{section}{Appendix II}

\subsection*{1. Operators of a prescribed class}
\label{app:II-1}

We shall apply the results of
\hyperref[sec:iv-normal-operators]{\S\,IV}--\hyperref[sec:v-spectral-form]{V},
in particular to Laurent
matrices and their generalizations.

Let \(\mathcal P\) be a subset of \(\mathfrak B\) for which
\(\mathcal P'\) is Abelian.  If \(R,R^{*}\) is a conjugate pair, we say
that it is \emph{of class \(\mathcal P\)} when
\[
  (R)'\supset\mathcal P,
\]
that is, when \(R\) commutes with every \(A,A^{*}\) as \(A\) ranges over
all of \(\mathcal P\).

Then \((R)''\subset\mathcal P'\), so \((R)''\) is Abelian; that is,
\(R,R^{*}\) is normal.  By \TheoremRef{16} it therefore has a unique maximal
extension \(\overline R,\overline R^{*}\), which is likewise normal.  In
the proof of \TheoremRef{14} we saw that
\[
  (\overline R)'\supset(R)',
\]
and hence \((\overline R)'\supset\mathcal P\); thus \(\overline R\) is
also of class \(\mathcal P\).  It is therefore enough to consider maximal
normal pairs \(R,R^{*}\).  We now assume \(R,R^{*}\) to be such a pair,
write \(\Delta(z)\) for its complex resolution of the identity, and let
\(\mathcal D\) be the set of all \(\Delta(z)\).  Since
\((R)'=\mathcal D'\) by \TheoremRef{20}, the pair \(R,R^{*}\) is of class
\(\mathcal P\) if and only if
\[
  \mathcal D'\supset\mathcal P.
\]
Because \((\Delta(z))'\supset\mathcal D'\supset\mathcal P\), every
\(\Delta(z)\) then belongs to class \(\mathcal P\); and since
\(\mathcal D'\) is the intersection of all \((\Delta(z))'\), this
condition is also sufficient.  Moreover,
\[
  \mathcal D'\supset\mathcal P
  \ \Longrightarrow\
  \mathcal P'\supset\mathcal D''\supset\mathcal D,
\]
whereas
\[
  \mathcal P'\supset\mathcal D
  \ \Longrightarrow\
  \mathcal D'\supset\mathcal P''\supset\mathcal P.
\]
Consequently \(\mathcal D\subset\mathcal P'\) as well; in other words,
it is necessary and sufficient that every \(\Delta(z)\) belong to
\(\mathcal P'\).

We now construct a set \(\mathcal P\) of the kind just described.  Let
\(G\) be a countably infinite Abelian group (the method can also be
generalized to continuous groups \(G\), but we shall not discuss this
here), with elements \(\alpha,\beta,\ldots\).  Let
\(\varphi_{\alpha},\varphi_{\beta},\ldots\) be a complete normalized
orthogonal system in \(\mathfrak H\), indexed not by the integers
\(1,2,\ldots\), but by the elements \(\alpha,\beta,\ldots\) of \(G\).

It is immediate that the operator \(U_{\alpha}\) defined by
\[
  U_{\alpha}
  \left(\sum_{\beta\in G}x_{\beta}\varphi_{\beta}\right)
  =
  \sum_{\beta\in G}x_{\alpha\beta}\varphi_{\beta},
  \qquad
  \left(\sum_{\beta\in G}\lvert x_{\beta}\rvert^{2}<\infty\right),
\]
is unitary.  We have
\[
  U_{\alpha}U_{\beta}=U_{\alpha\beta},
  \qquad
  U_{\alpha}^{*}=U_{\alpha}^{-1}=U_{\alpha^{-1}}.
\]
Let \(\mathcal P\) be the set of all \(U_{\alpha}\).  When does an
\(A\in\mathfrak B\) belong to \(\mathcal P'\)?  It must commute with every
\(U_{\alpha}\in\mathcal P\) (the relations
\(U_{\alpha}^{*}=U_{\alpha^{-1}}\) add nothing new), so that
\[
  AU_{\alpha}=U_{\alpha}A.
\]
This means
\[
  (AU_{\alpha}f,g)=(U_{\alpha}Af,g)
\]
for all \(f,g\).  Since both sides are linear and continuous in \(f\) and
in \(g\), it is enough to require this on a complete normalized
orthogonal system.  Taking \(f=\varphi_{\beta}\) and
\(g=\varphi_{\gamma}\), we obtain
\[
  (AU_{\alpha}\varphi_{\beta},\varphi_{\gamma})
  =
  (A\varphi_{\alpha^{-1}\beta},\varphi_{\gamma}),
\]
and
\[
  (U_{\alpha}A\varphi_{\beta},\varphi_{\gamma})
  =
  (A\varphi_{\beta},U_{\alpha}^{*}\varphi_{\gamma})
  =
  (A\varphi_{\beta},U_{\alpha^{-1}}\varphi_{\gamma})
  =
  (A\varphi_{\beta},\varphi_{\alpha\gamma}).
\]
Thus
\[
  (A\varphi_{\alpha^{-1}\beta},\varphi_{\gamma})
  =
  (A\varphi_{\beta},\varphi_{\alpha\gamma});
\]
equivalently, \((A\varphi_{\beta},\varphi_{\gamma})\) depends only on
\(\beta^{-1}\gamma\).

% [Source p. 420 / PDF 51]
If \(A,B\in\mathcal P'\), write
\[
  (A\varphi_{\varepsilon},\varphi_{\eta})
  =a(\varepsilon^{-1}\eta),
  \qquad
  (B\varphi_{\varepsilon},\varphi_{\eta})
  =b(\varepsilon^{-1}\eta).
\]
Then
\begin{align*}
  (AB\varphi_{\beta},\varphi_{\gamma})
  &=
  (B\varphi_{\beta},A^{*}\varphi_{\gamma})\\
  &=
  \sum_{\delta\in G}
  (B\varphi_{\beta},\varphi_{\delta})
  (\varphi_{\delta},A^{*}\varphi_{\gamma})\\
  &=
  \sum_{\delta\in G}
  (B\varphi_{\beta},\varphi_{\delta})
  (A\varphi_{\delta},\varphi_{\gamma})\\
  &=
  \sum_{\delta\in G}
  a(\delta^{-1}\gamma)b(\beta^{-1}\delta)\\
  &=
  \sum_{\substack{\varepsilon,\eta\in G\\
                   \varepsilon\eta=\beta^{-1}\gamma}}
  a(\varepsilon)b(\eta).
\end{align*}
Exactly the same expression results for
\[
  (BA\varphi_{\beta},\varphi_{\gamma})
  =
  \sum_{\substack{\varepsilon,\eta\in G\\
                   \varepsilon\eta=\beta^{-1}\gamma}}
  a(\varepsilon)b(\eta).
\]
Thus \((ABf,g)=(BAf,g)\) at least on the complete normalized orthogonal
system \(\{\varphi_{\alpha}\}\), and therefore, since both sides are
linear and continuous in \(f\) and in \(g\), for all \(f,g\).  Hence
\(AB=BA\): the operators \(A\) and \(B\) commute.  We have thereby shown
that \(\mathcal P'\), which contains \(A^{*}\) whenever it contains
\(A\), is Abelian.

\subsection*{2. Generalized Laurent matrices}
\label{app:II-2}

We now apply to the set \(\mathcal P\) just constructed what was said at
the beginning of \hyperref[app:II-1]{\S\,1}; that is, we investigate the pairs \(R,R^{*}\) of
this class.

Let \(a_{\alpha}\) be a sequence of numbers, likewise indexed by the
elements \(\alpha\in G\), subject only to
\[
  \sum_{\alpha\in G}\lvert a_{\alpha}\rvert^{2}<\infty.
\]
We define two operators \(R,R^{*}\), on the vectors
\(\varphi_{\alpha}\) and only on these, by
\[
  R\varphi_{\alpha}
  =
  \sum_{\beta\in G}a_{\alpha^{-1}\beta}\varphi_{\beta},
  \qquad
  R^{*}\varphi_{\alpha}
  =
  \sum_{\beta\in G}\overline{a_{\alpha\beta^{-1}}}\varphi_{\beta}.
\]
Indeed,
\[
  \sum_{\beta\in G}\lvert a_{\alpha^{-1}\beta}\rvert^{2}
  =
  \sum_{\beta\in G}\lvert a_{\alpha\beta^{-1}}\rvert^{2}
  =
  \sum_{\beta\in G}\lvert a_{\beta}\rvert^{2}
  <\infty.
\]
It is immediate that \(R,R^{*}\) form a conjugate pair in the sense of
\hyperref[def:conjugate-pair]{Definition~5}, and also that \(R\) commutes
with every \(U_{\alpha}\), and
hence with every \(U_{\alpha}^{*}=U_{\alpha^{-1}}\).  Thus
\((R)'\supset\mathcal P\).  The pair \(R,R^{*}\) is therefore of class
\(\mathcal P\); it is normal, and we may form
\(\overline R,\overline R^{*}\).  This pair is maximal normal and also
belongs to class \(\mathcal P\).

Using \TheoremRef{15}, we now determine the domain and range of
\(\overline R\) and \(\overline R^{*}\).  According to that theorem,
\(\overline Rf\) and \(\overline R^{*}f\) are defined if and only if
there exist two vectors \(f^{*},f^{**}\) such that, whenever \(Rg\) and
\(R^{*}g\) are defined,
\[
  (f,Rg)=(f^{**},g),
  \qquad
  (f,R^{*}g)=(f^{*},g).
\]
Since \(g\) must be one of the \(\varphi_{\alpha}\), writing
\[
  f=\sum_{\beta\in G}x_{\beta}\varphi_{\beta},
  \qquad
  \sum_{\beta\in G}\lvert x_{\beta}\rvert^{2}<\infty,
\]
this says
\[
  \sum_{\beta\in G}
  x_{\beta}\overline{a_{\alpha^{-1}\beta}}
  =
  (f^{**},\varphi_{\alpha}),
  \qquad
  \sum_{\beta\in G}
  x_{\beta}a_{\alpha\beta^{-1}}
  =
  (f^{*},\varphi_{\alpha}).
\]
Put
\[
  y_{\alpha}
  =
  \sum_{\beta\in G}a_{\alpha\beta^{-1}}x_{\beta},
  \qquad
  z_{\alpha}
  =
  \sum_{\beta\in G}\overline{a_{\alpha^{-1}\beta}}x_{\beta}.
\]

% [Source p. 421 / PDF 52]
Then we must have
\[
  (f^{*},\varphi_{\alpha})=y_{\alpha},
  \qquad
  (f^{**},\varphi_{\alpha})=z_{\alpha}.
\]
By \Eref, \S\,I, Theorems~5 and~7, this is possible if and only if
\[
  \sum_{\alpha\in G}\lvert y_{\alpha}\rvert^{2}<\infty,
  \qquad
  \sum_{\alpha\in G}\lvert z_{\alpha}\rvert^{2}<\infty,
\]
and in that case
\[
  f^{*}=\sum_{\alpha\in G}y_{\alpha}\varphi_{\alpha},
  \qquad
  f^{**}=\sum_{\alpha\in G}z_{\alpha}\varphi_{\alpha}.
\]
Our result is therefore the following.  For
\[
  f=\sum_{\alpha\in G}x_{\alpha}\varphi_{\alpha},
  \qquad
  \sum_{\alpha\in G}\lvert x_{\alpha}\rvert^{2}<\infty,
\]
the vectors \(\overline Rf\) and \(\overline R^{*}f\) are defined if and
only if the two sums
\[
  \sum_{\alpha\in G}\lvert y_{\alpha}\rvert^{2},
  \qquad
  \sum_{\alpha\in G}\lvert z_{\alpha}\rvert^{2}
\]
are finite, where \(y_{\alpha}\) and \(z_{\alpha}\) are as defined above.
In that case
\[
  \overline Rf=\sum_{\alpha\in G}y_{\alpha}\varphi_{\alpha},
  \qquad
  \overline R^{*}f=\sum_{\alpha\in G}z_{\alpha}\varphi_{\alpha}.
\]

On the other hand, by \TheoremRef{17} and \TheoremRef{18} the pair
\(\overline R,\overline R^{*}\) has exactly one resolution of the
identity \(\Delta(z)\); as established in \hyperref[app:II-1]{\S\,1}, each \(\Delta(z)\)
belongs to \(\mathcal P'\).  In the present case this means, as we saw
there, that
\[
  (\Delta(z)\varphi_{\alpha},\varphi_{\beta})
\]
depends only on \(\alpha^{-1}\beta\).  This is the general element of the
matrix of the projection operator \(\Delta(z)\) in the complete
normalized orthogonal system
\(\varphi_{\alpha},\varphi_{\beta},\ldots\).

If, for example, we choose for \(G\) the infinite cyclic group (the
integers \(0,\pm1,\pm2,\ldots\), with addition as the composition law),
we obtain a general eigenvalue theory of Laurent matrices,\footnote{\label{fn:71}These
and similar classes of matrices were related by Toeplitz to questions in
the theory of functions.  Cf.\ also the work cited above, \NoteRef{24}.}
including non-real and unbounded ones.  Other choices of \(G\) lead to
analogous but more general classes of matrices, which are now equally
under our control.

\section*{Appendix III}
\label{app:commutativity-unbounded}
\addcontentsline{toc}{section}{Appendix III}

\subsection*{1. A matrix criterion for commutativity}
\label{app:III-1}

We shall analyze somewhat more closely the notion of commutativity for
unbounded operators.  We defined the normality of \(R\), that is, the
commutativity of \(R,R^{*}\), by requiring \((R)''\) to be Abelian.  Thus,
if we form the Hermitian operators
\[
  S_{1}=\frac{R+R^{*}}{2},
  \qquad
  S_{2}=\frac{R-R^{*}}{2i},
\]
their commutativity is expressed by the Abelian character of
\((S_{1},S_{2})''\).\footnote{\label{fn:72}In the proof of \TheoremRef{14} we saw that
\((R)'=(R,R^{*})'=(S_{1},S_{2})'\), and therefore
\((R)''=(S_{1},S_{2})''\).}
It is natural, however, to try to use for this purpose the usual
definition of commutativity of matrices.

Let, for example, \(R,S\) be two Hermitian operators and let
\(\varphi_{1},\varphi_{2},\ldots\) be a complete normalized orthogonal
system in which both possess matrices, say
\(\{a_{\mu\nu}\}\) and \(\{b_{\mu\nu}\}\).\footnote{\label{fn:73}Cf.\ \Eref, Appendix~III,
and also the author's paper ``On the Theory of Unbounded Matrices,''
forthcoming in the \emph{Journal f\"ur die reine und angewandte
Mathematik}.}\TranslatorNote{The cited paper subsequently appeared as
J.~von Neumann, ``Zur Theorie der unbeschr\"ankten Matrizen,''
\emph{Journal f\"ur die reine und angewandte Mathematik} \textbf{161}
(1929), 208--236,
\href{https://doi.org/10.1515/crll.1929.161.208}
{doi:10.1515/crll.1929.161.208}.}
Since all row and column sums of squares of absolute values
\[
  \sum_{\mu=1}^{\infty}\lvert a_{\mu\nu}\rvert^{2}
  =
  \sum_{\mu=1}^{\infty}\lvert a_{\nu\mu}\rvert^{2},
  \qquad
  \sum_{\mu=1}^{\infty}\lvert b_{\mu\nu}\rvert^{2}
  =
  \sum_{\mu=1}^{\infty}\lvert b_{\nu\mu}\rvert^{2}
\]
are finite, the following series converge absolutely as well:

% [Source p. 422 / PDF 53]
\[
  \sum_{\rho=1}^{\infty}a_{\mu\rho}b_{\rho\nu},
  \qquad
  \sum_{\rho=1}^{\infty}b_{\mu\rho}a_{\rho\nu},
\]
and one might try to define commutativity by
\[
  \sum_{\rho=1}^{\infty}a_{\mu\rho}b_{\rho\nu}
  =
  \sum_{\rho=1}^{\infty}b_{\mu\rho}a_{\rho\nu}
  \qquad(\mu,\nu=1,2,\ldots).
\]

Since
\[
  a_{\mu\nu}=(R\varphi_{\mu},\varphi_{\nu}),
  \qquad
  b_{\mu\nu}=(S\varphi_{\mu},\varphi_{\nu})
\]
(cf., for example, \Eref, Appendix~III), we have
\begin{align*}
  \sum_{\rho=1}^{\infty}a_{\mu\rho}b_{\rho\nu}
  &=
  \sum_{\rho=1}^{\infty}
  (R\varphi_{\mu},\varphi_{\rho})
  (S\varphi_{\rho},\varphi_{\nu})\\
  &=
  \sum_{\rho=1}^{\infty}
  (R\varphi_{\mu},\varphi_{\rho})
  (\varphi_{\rho},S\varphi_{\nu})\\
  &=
  (R\varphi_{\mu},S\varphi_{\nu}),
\end{align*}
and similarly
\begin{align*}
  \sum_{\rho=1}^{\infty}b_{\mu\rho}a_{\rho\nu}
  &=
  \sum_{\rho=1}^{\infty}
  (S\varphi_{\mu},\varphi_{\rho})
  (R\varphi_{\rho},\varphi_{\nu})\\
  &=
  \sum_{\rho=1}^{\infty}
  (S\varphi_{\mu},\varphi_{\rho})
  (\varphi_{\rho},R\varphi_{\nu})\\
  &=
  (S\varphi_{\mu},R\varphi_{\nu}).
\end{align*}
Thus the commutativity condition is simply
\[
  (R\varphi_{\mu},S\varphi_{\nu})
  =
  (S\varphi_{\mu},R\varphi_{\nu}).
\]

If all \(R(S\varphi_{\mu})\) and \(S(R\varphi_{\mu})\) are defined, this
may be written
\[
  (SR\varphi_{\mu},\varphi_{\nu})
  =
  (RS\varphi_{\mu},\varphi_{\nu}),
  \qquad
  ((RS-SR)\varphi_{\mu},\varphi_{\nu})=0.
\]
Since the \(\varphi_{\nu}\) form a complete normalized orthogonal
system, \((RS-SR)\varphi_{\mu}=0\).  If any closed linear operator \(T\)
is an extension of \(i(RS-SR)\), then \(Tf=0\) for every
\(f=\varphi_{\mu}\), and hence for every finite linear aggregate of these
vectors, that is, on an everywhere dense set.  Since \(T\) is closed, it
follows that \(Tf\) is defined for every \(f\), and equals \(0\); thus
\(T=0\).  On the other hand, \(i(RS-SR)\) is plainly Hermitian: it is
defined for every \(\varphi_{\mu}\), so its domain spans the closed linear
manifold \(\mathfrak H\).  We may therefore form
\[
  T=\widetilde{i(RS-SR)},
\]
and this operator must be \(0\).

In this case we may therefore, with some justification, speak of
commutativity of \(R\) and \(S\).  In particular, this applies to
\(R,S\in\mathfrak B\), which are bounded.  For then \(i(RS-SR)\) is
defined everywhere and is itself closed and linear, so
\[
  i(RS-SR)=0,\qquad RS=SR.
\]
But what of arbitrary \(R,S\), for which the vectors
\(R(S\varphi_{\mu})\) and \(S(R\varphi_{\mu})\) need not all be defined?
To what extent does
\[
  (R\varphi_{\mu},S\varphi_{\nu})
  =
  (S\varphi_{\mu},R\varphi_{\nu})
  \qquad(\mu,\nu=1,2,\ldots)
\]
express a reasonable commutativity of \(R,S\), when \(R,S\) have
matrices in the complete normalized orthogonal system
\(\varphi_{1},\varphi_{2},\ldots\) (cf.\ \NoteRef{73})?

We shall show by a counterexample that, in this generality, the displayed
relation can hardly have any significance.

First, however, let us mention the notion of normality (for conjugate
pairs \(R,R^{*}\)) to which this notion of commutativity (for Hermitian
operators) leads.  We must form
\[
  S_{1}=\frac{R+R^{*}}{2},
  \qquad
  S_{2}=\frac{R-R^{*}}{2i},
\]
and require
\[
  (S_{1}\varphi_{\mu},S_{2}\varphi_{\nu})
  =
  (S_{2}\varphi_{\mu},S_{1}\varphi_{\nu}),
\]

% [Source p. 423 / PDF 54]
or, equivalently,
\[
  (R\varphi_{\mu},R\varphi_{\nu})
  =
  (R^{*}\varphi_{\mu},R^{*}\varphi_{\nu}).
\]
This plainly means
\[
  \lvert Rf\rvert=\lvert R^{*}f\rvert
\]
for every finite linear aggregate of the \(\varphi_{\mu}\).  We know from
\TheoremRef{21} that this is necessary for normality.

\subsection*{2. A counterexample}
\label{app:III-2}

Let \(\varphi_{1},\varphi_{2},\ldots\) be a complete normalized
orthogonal system.  Define two matrices \(\{a_{\mu\nu}\}\) and
\(\{b_{\mu\nu}\}\) as follows.  Unless both \(\mu\) and \(\nu\) are equal
to either \(2n-1\) or \(2n\), for the same \(n=1,2,\ldots\), set
\[
  a_{\mu\nu}=b_{\mu\nu}=0.
\]
Within each such \(2\times2\) block, set
\[
  a_{2n-1,2n-1}
  =
  a_{2n,2n-1}
  =
  a_{2n-1,2n}
  =
  a_{2n,2n}
  =
  n,
\]
\[
  b_{2n-1,2n-1}=b_{2n,2n}=n,
  \qquad
  b_{2n-1,2n}=-b_{2n,2n-1}=in.
\]
Let the closed linear Hermitian operators associated with the
(plainly Hermitian-square-summable) matrices
\(\{a_{\mu\nu}\}\) and \(\{b_{\mu\nu}\}\), relative to
\(\{\varphi_{1},\varphi_{2},\ldots\}\), be \(R\) and \(S\), respectively
(cf.\ \NoteRef{73}).

The formulas\TranslatorNote{In the displayed definition of \(\chi_{2n}\),
the source prints \(\varphi_{2n}\) in the first term.  The correction to
\(\varphi_{2n-1}\) is forced both by orthonormality and by the immediately
following diagonalization.}
\[
  \psi_{2n-1}
  =
  \frac{1}{\sqrt2}\varphi_{2n-1}
  +
  \frac{1}{\sqrt2}\varphi_{2n},
  \qquad
  \psi_{2n}
  =
  \frac{1}{\sqrt2}\varphi_{2n-1}
  -
  \frac{1}{\sqrt2}\varphi_{2n},
\]
\[
  \chi_{2n-1}
  =
  \frac{1}{\sqrt2}\varphi_{2n-1}
  +
  \frac{i}{\sqrt2}\varphi_{2n},
  \qquad
  \chi_{2n}
  =
  \frac{1}{\sqrt2}\varphi_{2n-1}
  -
  \frac{i}{\sqrt2}\varphi_{2n}
\]
define two further complete normalized orthogonal systems
\(\psi_{1},\psi_{2},\ldots\) and
\(\chi_{1},\chi_{2},\ldots\), and
\[
  R\psi_{2n-1}=2n\,\psi_{2n-1},
  \qquad
  R\psi_{2n}=0,
\]
\[
  S\chi_{2n-1}=2n\,\chi_{2n-1},
  \qquad
  S\chi_{2n}=0.
\]
The operators \(R,S\) are closed and linear.  In the complete normalized
orthogonal systems above, \(R\) and \(S\), respectively, therefore have
diagonal matrices; hence both are hypermaximal.\footnote{\label{fn:74}More generally,
if
\[
  T\xi_n=\alpha_n\xi_n
  \qquad(n=1,2,\ldots),
\]
where \(T\) is a closed linear operator and
\(\xi_1,\xi_2,\ldots\) is a complete normalized orthogonal system, then
\(T\) is hypermaximal.  Indeed, the closed linear manifolds
\(\mathfrak E,\mathfrak F\) (cf.\ \Eref, Chapter~V) contain all
\[
  (T\pm i1)\xi_n=(\alpha_n\pm i)\xi_n,
\]
and hence all \(\xi_n\).  Both manifolds are therefore equal to
\(\mathfrak H\).}

In particular, for
\[
  f=\sum_{\nu=1}^{\infty}y_{\nu}\psi_{\nu},
  \qquad
  \sum_{\nu=1}^{\infty}\lvert y_{\nu}\rvert^{2}<\infty,
\]
the vector \(Rf\) is defined if and only if
\[
  \sum_{n=1}^{\infty}4n^{2}\lvert y_{2n-1}\rvert^{2}<\infty.
\]
If instead
\[
  f=\sum_{\nu=1}^{\infty}x_{\nu}\varphi_{\nu},
  \qquad
  \sum_{\nu=1}^{\infty}\lvert x_{\nu}\rvert^{2}<\infty,
\]
then
\[
  y_{2n-1}
  =
  \frac{1}{\sqrt2}x_{2n-1}
  +
  \frac{1}{\sqrt2}x_{2n},
\]
so the characteristic condition is
\[
  \sum_{n=1}^{\infty}
  n^{2}\lvert x_{2n-1}+x_{2n}\rvert^{2}<\infty.
\]
Similarly, for
\[
  f=\sum_{\nu=1}^{\infty}z_{\nu}\chi_{\nu},
  \qquad
  \sum_{\nu=1}^{\infty}\lvert z_{\nu}\rvert^{2}<\infty,
\]
the vector \(Sf\) is defined precisely when
\[
  \sum_{n=1}^{\infty}4n^{2}\lvert z_{2n-1}\rvert^{2}<\infty.
\]
For \(f=\sum_{\nu=1}^{\infty}x_{\nu}\varphi_{\nu}\), where
\[
  z_{2n-1}
  =
  \frac{1}{\sqrt2}x_{2n-1}
  -
  \frac{i}{\sqrt2}x_{2n},
\]
this is characterized by
\[
  \sum_{n=1}^{\infty}
  n^{2}\lvert x_{2n-1}-ix_{2n}\rvert^{2}<\infty.
\]
The finiteness of both sums is, as one readily sees, equivalent to the
finiteness of
\[
  \sum_{n=1}^{\infty}
  n^{2}\bigl(\lvert x_{2n-1}\rvert^{2}
             +\lvert x_{2n}\rvert^{2}\bigr).
\]

% [Source p. 424 / PDF 55]
If, in addition, \(R(Sf)\) and \(S(Rf)\) are to be defined, write
\[
  Rf=\sum_{\nu=1}^{\infty}u_{\nu}\varphi_{\nu},
  \qquad
  Sf=\sum_{\nu=1}^{\infty}v_{\nu}\varphi_{\nu},
\]
where
\[
  u_{2n-1}=n x_{2n-1}+n x_{2n},
  \qquad
  u_{2n}=n x_{2n-1}+n x_{2n},
\]
\[
  v_{2n-1}=n x_{2n-1}-in x_{2n},
  \qquad
  v_{2n}=in x_{2n-1}+n x_{2n}.
\]
We must then require the finiteness of
\[
  \sum_{n=1}^{\infty}
  n^{2}\lvert v_{2n-1}+v_{2n}\rvert^{2},
  \qquad
  \sum_{n=1}^{\infty}
  n^{2}\lvert u_{2n-1}-iu_{2n}\rvert^{2},
\]
that is, of
\[
  \sum_{n=1}^{\infty}
  n^{4}\lvert x_{2n-1}+x_{2n}\rvert^{2},
  \qquad
  \sum_{n=1}^{\infty}
  n^{4}\lvert x_{2n-1}-ix_{2n}\rvert^{2}.
\]
The finiteness of these two sums is readily seen to be equivalent to that
of
\[
  \sum_{n=1}^{\infty}
  n^{4}\bigl(\lvert x_{2n-1}\rvert^{2}
             +\lvert x_{2n}\rvert^{2}\bigr).
\]

For such \(f\), we now calculate
\[
  i(RS-SR)f
  =
  \sum_{\nu=1}^{\infty}w_{\nu}\varphi_{\nu}.
\]
The result is\TranslatorNote{After defining the coefficients \(w_\nu\),
the source prints \(z_{2n-1}\) and \(z_{2n}\) in these two formulas.  They
have been corrected to \(w_{2n-1}\) and \(w_{2n}\).}
\[
  w_{2n-1}=-2n^{2}x_{2n-1},
  \qquad
  w_{2n}=2n^{2}x_{2n}.
\]
Thus \(i(RS-SR)\) is a closed linear Hermitian operator and has a
diagonal matrix in the system
\(\varphi_{1},\varphi_{2},\ldots\); it is therefore hypermaximal
(cf.\ \NoteRef{74}).  Plainly it is not \(0\); indeed,
\[
  f\neq0\quad\Longrightarrow\quad i(RS-SR)f\neq0.
\]

There can consequently be no question of commutativity of \(R\) and
\(S\).  Nevertheless, we shall find a complete normalized orthogonal
system \(\omega_{1},\omega_{2},\ldots\) in which both \(R\) and \(S\)
have matrices and the commutativity relations of \hyperref[app:III-1]{\S\,1} are satisfied.

\subsection*{3. Construction of the common matrix system}
\label{app:III-3}

For \(R\) and \(S\) to possess matrices \(\{a_{\mu\nu}\}\) and
\(\{b_{\mu\nu}\}\), respectively, in a complete normalized orthogonal
system \(\omega_{1},\omega_{2},\ldots\), means the following
(cf.\ \NoteRef{73}).  Necessarily
\[
  a_{\mu\nu}=(R\omega_{\mu},\omega_{\nu}),
  \qquad
  b_{\mu\nu}=(S\omega_{\mu},\omega_{\nu}).
\]
The elementary operators of these matrices, \(R'\) and \(S'\), are the
restrictions of \(R\) and \(S\), respectively, to the set
\(\omega_{1},\omega_{2},\ldots\), and their associated closed linear
operators are \(\widetilde{R'}\) and \(\widetilde{S'}\).  One requires
\[
  \widetilde{R'}=R,
  \qquad
  \widetilde{S'}=S.
\]
That is, for every \(f\) for which \(Rf\), respectively \(Sf\), is
defined, there must be a sequence \(f_{1},f_{2},\ldots\) in the linear
manifold---not merely the closed linear manifold---spanned by the
\(\omega_{1},\omega_{2},\ldots\), such that
\[
  f_n\to f,\qquad Rf_n\to Rf,
\]
respectively
\[
  f_n\to f,\qquad Sf_n\to Sf.
\]
If both \(Rf\) and \(Sf\) are defined, the two approximating sequences
may be different.

Suppose that \(\omega_{1},\omega_{2},\ldots\) is obtained from an
everywhere dense sequence \(g_{1},g_{2},\ldots\) by E.~Schmidt's
orthogonalization procedure.\footnote{\label{fn:75}Cf.\ E.~Schmidt at the place cited
in \NoteRef{52}; see also \Eref, \S\,I, Theorem~8.}
The required condition certainly holds if the sequences
\(f_{1},f_{2},\ldots\) just mentioned, with
\[
  f_n\to f,\qquad Rf_n\to Rf,
\]
respectively
\[
  f_n\to f,\qquad Sf_n\to Sf,
\]
can be chosen as subsequences of \(g_{1},g_{2},\ldots\).  Moreover, if
\[
  (Rg_p,Sg_q)=(Sg_p,Rg_q)
\]
always holds, then it holds also for the
\(\omega_{1},\omega_{2},\ldots\), which are linear aggregates of the
\(g\)'s:
\[
  (R\omega_{\mu},S\omega_{\nu})
  =
  (S\omega_{\mu},R\omega_{\nu}).
\]
Thus \(R,S\) satisfy the commutativity relation of \hyperref[app:III-1]{\S\,1}.  Notice that we
have assumed \(Rg_p\) and \(Sg_p\) to be defined, but not
\(R(Sg_p)\) and \(S(Rg_p)\).

% [Source p. 425 / PDF 56]
Let \(\overline g_{1},\overline g_{2},\ldots\) be the set of all linear
aggregates of finitely many
\(\varphi_{1},\varphi_{2},\ldots\), with rational coefficients, written
as a sequence.  Since \(R,S\) have matrices in the system
\(\varphi_{1},\varphi_{2},\ldots\), it is immediate that the sequences
\(f_{1},f_{2},\ldots\) mentioned above can be selected from
\(\overline g_{1},\overline g_{2},\ldots\); the rationality condition on
the coefficients causes no difficulty.  We now choose
\(h_{1},h_{2},\ldots\) so that all \(Rh_n,Sh_n\) are defined and, with
strong convergence as always in \(\mathfrak H\),
\[
  h_n\longrightarrow0,\qquad
  Rh_{2n-1}\longrightarrow0,\qquad
  Sh_{2n}\longrightarrow0
  \qquad(n\to\infty).
\]
The sequences \(f_{1},f_{2},\ldots\) above can then always be selected
as subsequences of
\[
  \overline g_{1}+h_{1},\ 
  \overline g_{2}+h_{3},\ldots,
  \qquad
  \overline g_{1}+h_{2},\
  \overline g_{2}+h_{4},\ldots.
\]
For brevity, write
\[
  \overline g_{1},\overline g_{1},
  \overline g_{2},\overline g_{2},\ldots
\]
as
\(\widetilde g_{1},\widetilde g_{2},\ldots\), and choose
\[
  g_{1},g_{2},\ldots
  =
  \widetilde g_{1}+h_{1},
  \widetilde g_{2}+h_{2},\ldots .
\]
Everything will therefore be accomplished once
\[
  (Rg_p,Sg_q)=(Sg_p,Rg_q)
\]
has been enforced by a suitable choice of \(h_{1},h_{2},\ldots\).
There is still some freedom in the choice of the \(h\)'s, whereas the
\(\widetilde g_{1},\widetilde g_{2},\ldots\) have already been fixed.

We reindex the sequence
\(\varphi_{1},\varphi_{2},\ldots\) as follows.  Divide it into pairs
\(\varphi_{2n-1},\varphi_{2n}\), replace the single index
\(n=1,2,\ldots\) by the triple \(p,q,\mu\), where each index ranges over
\(1,2,\ldots\), and denote \(\varphi_{2n-1},\varphi_{2n}\) by
\[
  \varphi'_{pq\mu},\qquad \varphi''_{pq\mu},
\]
respectively.  Denote the integer \(n\) belonging to \(p,q,\mu\) by
\(n(pq\mu)\), and arrange the indexing so that
\[
  n(pq\mu)\geq\mu^{2}
\]
always.  Define
\[
  h_p
  =
  \sum_{q,\mu=1}^{\infty}
  \frac{1}{pq\,[n(pq\mu)]^{3/2}}\,\varphi'_{pq\mu}
  +
  \sum_{r=1}^{p-1}x_{rp}v_{rp},
\]
where\TranslatorNote{In the even-\(p\) branch the source prints
\(\varphi'_{r,q,\rho(rp)}\) in the first term.  The free index \(q\) has
been corrected to \(p\), as required by the definition of \(v_{rp}\).}
\[
  v_{rp}
  =
  \begin{cases}
    \varphi'_{r,p,\rho(rp)}-\varphi''_{r,p,\rho(rp)},
      &p\ \text{odd},\\
    \varphi'_{r,p,\rho(rp)}
      -i\varphi''_{r,p,\rho(rp)},
      &p\ \text{even}.
  \end{cases}
\]

The quantities \(x_{rp}\) and \(\rho(rp)\), for \(r<p\), will be chosen
later.

First consider
\[
  \xi_p
  =
  \sum_{q,\mu=1}^{\infty}
  \frac{1}{pq\,[n(pq\mu)]^{3/2}}\,\varphi'_{pq\mu}.
\]
These sums converge, since the sums of the squares of the absolute
values of their coefficients satisfy
\begin{align*}
  \sum_{q,\mu=1}^{\infty}
  \frac{1}{p^{2}q^{2}[n(pq\mu)]^{3}}
  &\leq
  \sum_{q,\mu=1}^{\infty}
  \frac{1}{p^{2}q^{2}\mu^{6}}\\
  &=
  \frac{1}{p^{2}}
  \left(\sum_{q=1}^{\infty}\frac1{q^{2}}\right)
  \left(\sum_{\mu=1}^{\infty}\frac1{\mu^{6}}\right).
\end{align*}
Thus \(\xi_p\to0\) as \(p\to\infty\).  Moreover,
\begin{align*}
  \sum_{q,\mu=1}^{\infty}
  \frac{[n(pq\mu)]^{2}}{p^{2}q^{2}[n(pq\mu)]^{3}}
  &=
  \sum_{q,\mu=1}^{\infty}
  \frac{1}{p^{2}q^{2}n(pq\mu)}\\
  &\leq
  \sum_{q,\mu=1}^{\infty}
  \frac{1}{p^{2}q^{2}\mu^{2}}\\
  &=
  \frac{1}{p^{2}}
  \left(\sum_{q=1}^{\infty}\frac1{q^{2}}\right)
  \left(\sum_{\mu=1}^{\infty}\frac1{\mu^{2}}\right)
\end{align*}
is finite, and even tends to \(0\) with \(p\to\infty\).  Hence
\(R\xi_p\) and \(S\xi_p\) are defined.

% [Source p. 426 / PDF 57]
The sum
\[
  \sum_{r=1}^{p-1}x_{rp}v_{rp}
\]
is finite, and \(R,S\) are therefore applicable to it.  This sum tends
to \(0\) as \(p\to\infty\) whenever
\[
  \sum_{r=1}^{p-1}\lvert x_{rp}\rvert^{2}\longrightarrow0,
\]
which certainly holds, for example, if
\[
  \lvert x_{rp}\rvert\leq\frac1{rp}.
\]
Furthermore, for odd \(p\) we have \(Rv_{rp}=0\), and for even \(p\) we
have \(Sv_{rp}=0\).  Thus applying \(R\), respectively \(S\), to
\(\sum_{r=1}^{p-1}x_{rp}v_{rp}\) gives \(0\).  In
summary:\TranslatorNote{At this point the source prints \(h_p\to\infty\);
the construction and the estimate immediately above require \(h_p\to0\).}
\[
  Rh_p\ \text{and}\ Sh_p\ \text{are always defined},\qquad
  h_p\to0,
\]
\[
  Rh_p\to0\quad\text{for odd \(p\)},\qquad
  Sh_p\to0\quad\text{for even \(p\)}.
\]

It remains to ensure, by a suitable choice of \(x_{rp}\) and
\(\rho(rp)\), that
\[
  (Rg_p,Sg_q)-(Sg_p,Rg_q)=0.
\]
It is plainly enough to consider \(p<q\), as we shall do, while retaining
\(\lvert x_{rp}\rvert\leq1/(rp)\).  Substitute
\[
  g_p
  =
  \widetilde g_p
  +
  \sum_{t,\mu=1}^{\infty}
  \frac{1}{pt\,[n(pt\mu)]^{3/2}}\varphi'_{pt\mu}
  +
  \sum_{r=1}^{p-1}x_{rp}v_{rp},
\]
\[
  g_q
  =
  \widetilde g_q
  +
  \sum_{\mu,\nu=1}^{\infty}
  \frac{1}{q\mu\,[n(q\mu\nu)]^{3/2}}\varphi'_{q\mu\nu}
  +
  \sum_{s=1}^{q-1}x_{sq}v_{sq}
\]
in the formula above and expand.  In view of the existing
orthogonalities, and since \(p<q\), the result is
\begin{align*}
  &(\{SR-RS\}\widetilde g_p,\widetilde g_q)
  +(\{SR-RS\}\widetilde g_p,\xi_q)
  +(\xi_p,\{RS-SR\}\widetilde g_q)\\
  &\quad+
  \sum_{s=1}^{q-1}\overline{x}_{sq}
  (\{SR-RS\}\widetilde g_p,v_{sq})\\
  &\quad+
  \sum_{r=1}^{p-1}x_{rp}
  (v_{rp},\{RS-SR\}\widetilde g_q)\\
  &\quad+
  \frac{\overline{x}_{pq}}
  {pq\,[n(p,q,\rho(pq))]^{3/2}}
  \bigl(\{SR-RS\}\varphi'_{p,q,\rho(pq)},v_{pq}\bigr)
  =0.
\end{align*}
For brevity put
\[
  i(RS-SR)=T.
\]
We also introduce the following notation.  The vector
\(\widetilde g_l\) is a linear aggregate of finitely many
\(\varphi_n\), and hence of finitely many
\(\varphi'_{pq\mu}\) and \(\varphi''_{pq\mu}\); let \(\mu(l)\) denote the
largest \(\mu\) occurring in it.  Require henceforth that
\[
  \rho(rt)>\max\bigl(\mu(1),\ldots,\mu(t-1)\bigr)
\]
always.  The fourth term in the preceding expression then vanishes
identically, leaving
\begin{align*}
  &(T\widetilde g_p,\widetilde g_q)
  +(T\widetilde g_p,\xi_q)
  +(\xi_p,T\widetilde g_q)
  +\sum_{r=1}^{p-1}x_{rp}(v_{rp},T\widetilde g_q)\\
  &\qquad+
  \frac{\overline{x}_{pq}}
  {pq\,[n(p,q,\rho(pq))]^{3/2}}
  (T\varphi'_{p,q,\rho(pq)},v_{pq})
  =0.
\end{align*}
Since
\[
  (T\varphi'_{p,q,\rho(pq)},v_{pq})
  =
  -2[n(p,q,\rho(pq))]^{2},
\]
we obtain\TranslatorNote{The source prints a minus sign and omits the
factor \(pq\) in this recursion.  Solving the preceding equation gives
the corrected formula displayed here; the ensuing estimate and choice
of \(\rho(pq)\) have been corrected consistently.}
\[
  \overline{x}_{pq}
  =
  \frac{pq}{2[n(p,q,\rho(pq))]^{1/2}}
  \left[
  \begin{aligned}
    &(T\widetilde g_p,\widetilde g_q)
    +(T\widetilde g_p,\xi_q)
    +(\xi_p,T\widetilde g_q)\\
    &\quad+
    \sum_{r=1}^{p-1}x_{rp}(v_{rp},T\widetilde g_q)
  \end{aligned}
  \right].
\]

% [Source p. 427 / PDF 58]
Here \(x_{pq}\) is expressed in terms of the \(x_{rp}\) and \(v_{rp}\),
and hence of the \(\rho(rp)\), for \(r=1,\ldots,p-1\), together with
\(\rho(pq)\).  The latter remains arbitrary.  If the \(x_{rs}\) and
\(\rho(rs)\) are ordered by increasing \(r+s\), the displayed equation
is therefore a recursion for the \(x_{rs}\), while the \(\rho(rs)\) may
be chosen freely.  It remains only to maintain
\[
  \lvert x_{pq}\rvert\leq\frac1{pq},
  \qquad
  \rho(pq)\geq
  \max\bigl(\mu(1),\ldots,\mu(q-1)\bigr).
\]
Choose the \(\rho(rs)\) in the same order.  If \(C\) denotes the absolute
value of the expression in brackets in the formula for
\(\overline{x}_{pq}\), then \(C\) depends only on the \(x_{rs}\) and
\(\rho(rs)\) for which \(r+s<p+q\), and not on \(\rho(pq)\).  Hence
\[
  \lvert x_{pq}\rvert
  =
  \frac{pq\,C}{2[n(p,q,\rho(pq))]^{1/2}}
  \leq
  \frac{pq\,C}{2\rho(pq)}.
\]
It is therefore enough to choose
\[
  \rho(pq)\geq\frac12 Cp^{2}q^{2},
  \qquad
  \rho(pq)\geq
  \max\bigl(\mu(1),\ldots,\mu(q-1)\bigr),
\]
and our objective is attained.

We have now constructed the sequences
\[
  h_{1},h_{2},\ldots,
  \qquad
  g_{1},g_{2},\ldots,
\]
and the system
\[
  \omega_{1},\omega_{2},\ldots
\]
which we sought.  The desired construction is therefore complete.

\begin{center}
  \emph{(Received February 20, 1929.)}
\end{center}

\end{document}
